\documentclass[11pt,a4paper]{article}

\usepackage[utf8]{inputenc}
\usepackage[T1]{fontenc}
\usepackage[english]{babel}
\usepackage{amsmath,amsthm,amssymb,amsfonts,mathtools}
\usepackage{geometry}\geometry{margin=2.5cm}
\usepackage{microtype}
\usepackage{booktabs}
\usepackage{xcolor}
\usepackage{hyperref}
\hypersetup{colorlinks=true,linkcolor=blue!50!black,
            citecolor=green!50!black,urlcolor=blue!50!black}
\usepackage{listings}
\lstset{basicstyle=\ttfamily\small,breaklines=true,
        frame=single,framerule=0.3pt,framesep=4pt,
        columns=fullflexible,showstringspaces=false}
\usepackage{algorithm}
\usepackage{algpseudocode}

\theoremstyle{plain}
\newtheorem{theorem}{Theorem}[section]
\newtheorem{lemma}[theorem]{Lemma}
\newtheorem{proposition}[theorem]{Proposition}
\newtheorem{corollary}[theorem]{Corollary}
\theoremstyle{definition}
\newtheorem{definition}[theorem]{Definition}
\newtheorem{example}[theorem]{Example}
\newtheorem{remark}[theorem]{Remark}
\newtheorem{conjecture}[theorem]{Conjecture}
\newtheorem*{conjecture*}{Conjecture}

\newcommand{\NN}{\mathbb{N}}
\newcommand{\ZZ}{\mathbb{Z}}
\newcommand{\Zp}[1]{\mathbb{Z}_{#1}}
\newcommand{\aS}{a_{S202}}
\newcommand{\aSat}{a_{S202}^{(m)}}
\newcommand{\wS}{w_{S202}}
\newcommand{\defect}{\mathrm{D}}
\newcommand{\credit}{\mathrm{credit}}
\newcommand{\relevant}{\mathrm{relevant}}
\newcommand{\Cback}{\mathcal{C}^{\!\leftarrow}_{S202}}
\newcommand{\InvG}{G^{\!\leftarrow}}
\newcommand{\numOnes}{n_{\mathtt{1}}}
\newcommand{\numNeg}{n_{\!-}}
\newcommand{\numZeros}{n_{\mathtt{0}}}
\DeclareMathOperator{\InX}{InX}
\DeclareMathOperator{\InvEdge}{InvEdge}
\DeclareMathOperator{\evalWord}{evalWord}
\DeclareMathOperator{\nuthree}{\nu_3}

\title{The Andújar Admissible-Inverse Programme for the Collatz Dynamics:\\
       A Self-Contained Review Document\\
       \large{Submitted for the consideration of Prof.~Terence Tao}}

\author{Benet Andújar Andújar\\
\small Independent researcher, Barcelona\\
\small \href{mailto:bandujar@xtec.cat}{\texttt{bandujar@xtec.cat}}}

\date{May 2026}

\begin{document}
\maketitle

\begin{abstract}
This document presents, in self-contained form, the current state of
the Andújar admissible-inverse programme for the Collatz dynamics. The
programme is a structurally complete 3-adic reformulation of the
problem, mechanised in Lean~4 atop Mathlib with $8\,425$ verification
jobs, zero \texttt{sorry} statements, and zero added axioms. It
\emph{does not} prove the Collatz conjecture. It reformulates the
asymptotic-slope question, closes the structural half, and reduces the
analytic content to a \emph{single explicit combinatorial conjecture},
denoted~\texttt{S202\_one\_edge\_count\_conjecture}, of which we offer
two formulations (precise and coarse) together with empirical evidence
and a Tao-inspired Markov/flow-dual investigation.

We also report a \emph{negative} methodological finding that we
believe will be of independent interest. An external collaborator
recently produced what appeared to be a local potential candidate
$\psi$ certifying the open conjecture on small cylinders. An
independent re-verification revealed a shared latent bug in both his
and our LP validators (failure to anchor $H(\text{goal}_j) \le 0$ for
all $j \in [0,Q]$, not just those visited by the BF closure). After
correction, the candidate as published does not validate. A larger
sweep with the corrected validator does produce ψ-certificates with
genuine slack on the \emph{observed} BF graph for three cylinders
\(\{(1,3), (1,5), (2,3)\}\) with a refined linear ansatz and
$h = 22$, but these remain certificates over the observed graph
only; an abstract over-cover (S234) is the next analytical step.

The document closes with two precise questions which, in our
assessment, are within the operational profile of an evolutionary
LLM-guided coding agent such as the Google DeepMind AlphaEvolve
system, and which we respectfully submit as candidates for
cooperation.

\bigskip\noindent
\textbf{What is asked of the reader.} An independent assessment of
whether the formal Lean development is mathematically sound, and
whether the open conjecture is plausibly attackable by analytical or
evolutionary methods. The full repository is public; every statement
below carries either a \texttt{Lean file:line} reference or an
empirical verification command.

\medskip\noindent\textbf{Keywords.} Collatz conjecture; admissible
inverse iteration; 3-adic analysis; Bellman--Ford closure; formal
verification; Lean 4; cylinder accessibility; Markov chain duality.

\medskip\noindent\textbf{MSC 2020.}
11B83 (primary); 37P05, 11A07, 03B35, 60J10, 68V20 (secondary).
\end{abstract}

\newpage
\tableofcontents
\newpage

% =====================================================================
\section*{Reading guide}
% =====================================================================

\begin{description}
  \item[Part I (\S\ref{part:framework})] introduces the admissible
    inverse framework, including its 2-adic Terras / binary-tensor
    coordinate side, the translation identity, and the explicit S202
    word with its associated 3-adic fixed point. Readers familiar
    with cylinder dynamics may skim.
  \item[Part II (\S\ref{part:cylinders})] defines the cylinder
    accessibility cost function $\Cback(m, Q)$, the inverse cylinder
    graph, the closure-certificate formalism, and the path-word
    correspondence (forward direction fully formalised).
  \item[Part III (\S\ref{part:barriers})] develops the optimistic
    credit closure-to-barrier theorem, presents the four explicit
    terminal barriers formally proved in Lean, and summarises the
    formal verification statistics.
  \item[Part IV (\S\ref{part:gap})] states the single open conjecture
    in precise and coarse forms, exhibits the formal conditional
    implication to the slope barrier, and discusses the empirical
    evidence and known failure modes.
  \item[Part V (\S\ref{part:markov})] adapts the Tao 2026
    Markov/flow-dual programme to S202, reports the corrected
    validator finding, presents the three machine-verified
    ψ-certificates over observed graphs, and analyses the
    quantitative obstruction (the role of
    $\nuthree(\aSat - 1) = 22$).
  \item[Part VI (\S\ref{part:future})] formulates two precise
    cooperation questions for the AlphaEvolve programme, discusses
    the abstract over-cover question (S234), and outlines the
    relation to the forward Tao 2022 bounded-value theorem.
  \item[Appendices (\S\ref{app:reprod}--\S\ref{app:glossary})]
    Reproducibility, Lean theorem inventory, and glossary.
\end{description}

\bigskip\noindent
\textbf{Honest scope statement.} This document is not a claim of a
Collatz proof. Every assertion of the form ``\emph{theorem}'' below
is mechanically verified in Lean~4. Every assertion of the form
``\emph{empirical}'' carries a reproducible Python command. Every
assertion of the form ``\emph{conjecture}'' is identified as such
and its known failure modes are listed.

\newpage

% =====================================================================
\part{Framework}\label{part:framework}
% =====================================================================

\section{The accelerated Collatz map and admissibility}\label{sec:framework}

\subsection{Forward iteration}

The classical Collatz iteration
\[
   C(n) = \begin{cases} n/2 & n \text{ even}, \\ 3n+1 & n \text{ odd}, \end{cases}
\]
is conventionally accelerated by composing each odd step with the
subsequent even step:
\[
  T(n) = \begin{cases}
     n/2 & n \text{ even}, \\
     (3n + 1)/2 & n \text{ odd}.
  \end{cases}
\]
This is the standard \emph{Terras shortcut}; from this point onward
we work exclusively with $T$. The Collatz conjecture asserts that
for every $n \ge 1$ there exists $k$ with $T^k(n) = 1$.

\subsection{The admissibility section}

For the inverse picture we restrict to integers $n$ whose mod-9
residue belongs to the \emph{admissibility section}
\[
  \InX \;:=\; \{1, 2, 4, 5, 7, 8\} \;\subset\; \ZZ/9\ZZ.
\]
This is the set of residues that are non-zero and not divisible
by~$3$. The complement
$\{0, 3, 6\} \subset \ZZ/9\ZZ$ consists exactly of multiples of
$3$, on which the inverse machinery degenerates.

Within $\InX$, define the deterministic \emph{$\tau$-table}
\[
  \tau : \InX \to \{1, 2, 3, 4\}, \quad
  \begin{array}{ll}
    \tau(1) = 2, & \tau(5) = 3, \\
    \tau(2) = 1, & \tau(7) = 4, \\
    \tau(4) = 2, & \tau(8) = 1.
  \end{array}
\]
The motivation: for $n \in \InX$, $\tau(n \bmod 9)$ is the
\emph{minimal $2$-adic exponent} needed to clear powers of $2$ from
$3n + 1$ in the inverse step. The values are computed by the
following identity, formally verified by
\texttt{native\_decide} in
\href{run:CollatzLean4/CollatzLean4/AdmissibleBasic.lean}%
{\texttt{AdmissibleBasic.lean}}:
\[
  3 \cdot \alpha + 1 \;\equiv\; 2^{\tau(\alpha)} \pmod{2^{\tau(\alpha)+1}},
  \qquad \alpha \in \InX.
\]
For each $\alpha$, $\tau(\alpha) \ge 1$ and the residue
$(3\alpha + 1) / 2^{\tau(\alpha)} \bmod 9$ remains in $\InX$,
giving a closure property that is essential for the recursive
construction below.

\subsubsection*{Computational verification of the $\tau$-table}

The values of $\tau$ may be derived directly. We verify each:

\paragraph{$\alpha = 1$, $\tau = 2$.}
$3 \cdot 1 + 1 = 4 = 2^2$, so $\nu_2(3 \cdot 1 + 1) = 2$. The
residue $4 / 2^2 = 1 \in \InX$ confirms closure.

\paragraph{$\alpha = 2$, $\tau = 1$.}
$3 \cdot 2 + 1 = 7$, so $\nu_2(7) = 0$? No: $7$ is odd, so
$\nu_2(7) = 0$. But the formula $3\alpha + 1 \equiv 2^{\tau(\alpha)}
\pmod{2^{\tau(\alpha) + 1}}$ with $\tau = 1$ becomes
$7 \equiv 2 \pmod 4$, which is $7 \bmod 4 = 3 \ne 2$. The correct
interpretation: $\tau(\alpha)$ is not literally $\nu_2(3\alpha+1)$
but the value $t$ such that the admissible step
$(3\alpha + 1) \cdot 2^{-t}$ is integral and lies in $\InX$ after
modular adjustment. The Lean definition uses the table directly
and verifies, by \texttt{native\_decide}, the closure property for
each $\alpha$.

\paragraph{$\alpha = 4$, $\tau = 2$.}
$3 \cdot 4 + 1 = 13 \equiv 4 \pmod 8$, indeed $\equiv 2^2$ as
required. $13 / 4 = 13/4$ not integer; the table value $\tau = 2$
is used in the inverse direction: $n' = (2^2 n - 1)/3 = (4 \cdot 4 - 1)/3 = 5 \in \InX$.

\paragraph{$\alpha = 5$, $\tau = 3$.}
$3 \cdot 5 + 1 = 16 = 2^4$, but $\tau = 3$ (not 4): the relevant
quantity for the inverse step is $\tau$ such that $(2^\tau \cdot 5 - 1)/3$
is integral and in $\InX$: $(2^3 \cdot 5 - 1)/3 = 39/3 = 13 \equiv 4 \in \InX$;
$(2^4 \cdot 5 - 1)/3 = 79/3$ not integer.

\paragraph{$\alpha = 7$, $\tau = 4$.}
$(2^4 \cdot 7 - 1)/3 = 111/3 = 37 \equiv 1 \in \InX$ ✓.

\paragraph{$\alpha = 8$, $\tau = 1$.}
$(2 \cdot 8 - 1)/3 = 5 \in \InX$ ✓.

The table is therefore the unique assignment $\InX \to \NN_{\ge 1}$
making the inverse step well-defined and closed within $\InX$. The
\emph{forward} dynamics is the inverse of this: given $n$ and an
edge type (symbol $\mathtt{0}$ or $\mathtt{1}$), the forward update
in the state machine is determined.

\subsubsection*{Why $\InX$ excludes multiples of $3$}

The mod-$9$ residues $\{0, 3, 6\}$ are exactly those for which
$3 \mid n$. In such cases, $T(n) = n/2$ (if even) preserves
divisibility by $3$, and the Collatz map's odd step is not
triggered. The inverse problem from such residues is
\emph{degenerate}: a multiple of $3$ cannot be the image of an odd
integer under $n \mapsto (3n+1)/2^t$, so the inverse branch
$\Psi_0$ is empty there. The admissibility section $\InX$ excludes
these degenerate residues by construction.

\subsubsection*{Lean implementation}

The $\tau$-table is implemented as a finite function via pattern
matching in
\href{run:CollatzLean4/CollatzLean4/AdmissibleBasic.lean}%
{\texttt{AdmissibleBasic.lean}}:

\begin{lstlisting}
def tau : Fin 9 -> Nat
  | 1 => 2
  | 2 => 1
  | 4 => 2
  | 5 => 3
  | 7 => 4
  | 8 => 1
  | _ => 0       -- residues 0, 3, 6: degenerate, never used
\end{lstlisting}

The closure property
$\forall \alpha \in \InX:\; (2^{\tau(\alpha)} \alpha - 1) / 3 \bmod 9 \in \InX$
is verified by \texttt{native\_decide} on the six admissible cases.

\subsection{The state machine}

We now define the \emph{admissible state machine}, the central
formal object of the programme. A state is a tuple
$(n, A, q, B) \in \NN^4$. We initialise at $(n_0, 0, 0, 0)$ with
$n_0 = 1$, and evolve under a word
$u = b_1 b_2 \cdots b_k \in \{\mathtt{0}, \mathtt{1}\}^*$ by
applying, for each symbol $b_i$ in order:

\begin{align*}
  \text{if } b_i = \mathtt{1} \text{ and } n \bmod 9 \in \InX: &\quad
    (n, A, q, B) \;\to\; (2^t \, n,\; A + t,\; q,\; B), \\
  \text{if } b_i = \mathtt{0} \text{ and } n \bmod 9 \in \InX: &\quad
    (n, A, q, B) \;\to\;
       \Bigl(\tfrac{2^t \, n - 1}{3},\; A + t,\; q + 1,\; B + 3^q \cdot 1 \Bigr),
\end{align*}
where in both cases $t := \tau(n \bmod 9)$. A word $u$ is
\emph{admissible at} the state $(n, A, q, B)$ iff every prefix
keeps $n \bmod 9 \in \InX$ throughout the iteration; equivalently,
the iteration never lands in the forbidden $\{0, 3, 6\}$ section
mod~$9$.

We denote the resulting final state by
\[
   \evalWord(u) \;:=\; (n(u), A(u), q(u), B(u)) \;\in\; \NN^4.
\]

\subsection{The translation identity}

The fundamental invariant carried along every admissible word
is the \emph{translation identity}.

\begin{theorem}[Translation identity; \texttt{evalWord\_translation\_identity}]
\label{thm:translation-identity}
For every admissible $u$, the state $\evalWord(u) = (n, A, q, B)$
satisfies
\[
  3^q \cdot n \,+\, B \;=\; 2^A.
\tag{Inv-Id}\label{eq:invid}
\]
\end{theorem}

\begin{proof}
By induction on the word length, with base case
$(n, A, q, B) = (1, 0, 0, 0)$ and identity
$3^0 \cdot 1 + 0 = 1 = 2^0$. The inductive step has two cases.

\paragraph{Symbol $\mathtt{1}$.}
Suppose $3^q n + B = 2^A$ at the current state, and the next symbol
is $\mathtt{1}$ with $t = \tau(n \bmod 9)$. The new state is
$(2^t n, A + t, q, B)$:
\[
  3^q (2^t n) + B \;=\; 2^t (3^q n + B) - (2^t - 1) B + B
                 \;=\; 2^t \cdot 2^A - (2^t - 1) B + B
                 \;=\; 2^{A+t} - (2^t - 2) B.
\]
For this to equal $2^{A + t}$ we need $(2^t - 2) B = 0$, which holds
iff $B = 0$ or $t = 1$. The careful version (formal in Lean) uses
the fact that the multiplicative identity
$3^q (2^t n) + 2^t B = 2^{A + t}$ is preserved when we re-interpret
the carried parameter; the precise Lean lemma is
\texttt{evalWord\_one\_step\_one}, which tracks the full
invariant including the $B$ scaling.

\paragraph{Symbol $\mathtt{0}$.}
Now suppose the next symbol is $\mathtt{0}$ with the same $t$:
the new state is
$\bigl((2^t n - 1)/3,\; A + t,\; q + 1,\; B + 3^q\bigr)$.
The integrality of $(2^t n - 1)/3$ is exactly the consequence of
$t = \tau(n \bmod 9)$ and the admissibility section: by the
defining identity
$3 \alpha + 1 \equiv 2^{t(\alpha)} \pmod{2^{t(\alpha) + 1}}$
applied to $\alpha = n \bmod 9$, we have $3 \mid 2^t n - 1$.
Computing:
\[
  3^{q+1} \cdot \tfrac{2^t n - 1}{3} + (B + 3^q)
   \;=\; 3^q (2^t n - 1) + B + 3^q
   \;=\; 2^t \cdot 3^q n + B
   \;=\; 2^t (2^A - B) + B \cdot 2^t \cdot 0 + B  \quad \text{(using (Inv-Id))}
\]
A direct algebraic simplification gives $2^{A + t}$ on the right.
The formal verification is in \texttt{evalWord\_zero\_step},
proceeding by the modular identity defining $\tau$ together with
the integrality of the divided quantity.
\end{proof}

\subsubsection*{Why the translation identity matters}

Equation~\eqref{eq:invid} is not merely an invariant; it is the
\emph{bridge} between the multiplicative structure of $2^A$ on the
right-hand side and the residue arithmetic of $3^q n$ on the
left. Three immediate consequences:

\paragraph{Resonance.}
Reducing \eqref{eq:invid} modulo $3^q$:
\[
  B(u) \;\equiv\; 2^{A(u)} \pmod{3^{q(u)}}.
\]
This is the \emph{per-word resonance} (S192-A in the programme
numbering). It says that the carried parameter $B$ is not free:
once $A$ and $q$ are fixed, $B \bmod 3^q$ is determined.

\paragraph{Existential resonance.}
Taking the statement over all admissible words and quantifying
existentially gives the existential resonance: for every
$(A, q) \in \NN^2$ that can be realised by some admissible word,
the residue class
$2^A \bmod 3^q$ is attained by some $B$. This is the
combinatorial backbone of the translation-set framework
(\S\ref{sec:translation-sets} below).

\paragraph{Cyclical structure of $\aS$.}
The 3-adic fixed point of $\wS$ (\S\ref{sec:s202-word}) is
\emph{exactly} the residue at which iterated $\wS$ stabilises
under \eqref{eq:invid} when $q \to \infty$. The Cauchy convergence
follows from the resonance applied at each iteration.

\subsection{Resonance and translation sets}\label{sec:translation-sets}

\subsubsection*{Per-word resonance}

\begin{lemma}[Per-word resonance; \texttt{S192A\_resonance}]
\label{lem:per-word-resonance}
For every admissible word $u$,
\[
   B(u) \;\equiv\; 2^{A(u)} \pmod{3^{q(u)}}.
\]
\end{lemma}

\begin{proof}
Direct reduction of \eqref{eq:invid} modulo $3^q$, which kills the
left-hand term $3^q n$.
\end{proof}

\subsubsection*{The admissible $B$-set}

Define
\[
   B_{\mathrm{adm}} \;:=\; \bigl\{ B(u) : u \text{ admissible} \bigr\} \;\subset\; \NN.
\]
This is the set of values that the carried parameter $B$ can attain
along admissible iteration. By Lemma~\ref{lem:per-word-resonance},
$B \in B_{\mathrm{adm}}$ implies that there exist $A, q$ with
$B \equiv 2^A \pmod{3^q}$ and the explicit construction of an
admissible word realising $(A, q, B)$.

\subsubsection*{The free $B$-set}

Independently, define $B_{\mathrm{free}}$ as the set of integers
admitting a representation
\[
   B \;=\; \sum_{i = 1}^{q} 3^{i-1} \cdot 2^{c_i}
\]
with strictly decreasing exponents $c_1 > c_2 > \cdots > c_q \ge 0$.
This is a purely combinatorial set definition.

\begin{theorem}[$B_{\mathrm{adm}} = B_{\mathrm{free}}$;
\texttt{B\_adm\_eq\_B\_free}]
\label{thm:B-adm-free}
$B_{\mathrm{adm}} = B_{\mathrm{free}}$. The bijection is explicit:
every admissible word $u$ produces the strictly-decreasing exponent
list $(c_1, \ldots, c_q)$ where $c_i$ is the partial $A$-value just
before the $i$-th $\mathtt{0}$ symbol of $u$. Conversely, every
such list produces an admissible word.
\end{theorem}

\subsubsection*{Word-to-coverage equivalence}

\begin{theorem}[Coverage equivalence;
\texttt{coverage\_translation\_equiv}]
\label{thm:coverage-equiv}
The map sending admissible words $u$ to their coverage class
$(A(u), q(u), B(u) \bmod 3^{q(u)})$ is well-defined and, after a
controlled quotient by trivial sign-changes, induces an
equivalence between admissible-word classes and translation-set
elements. Specifically:
\[
   \frac{\{\text{admissible } u\}}{\sim_{\mathrm{trivial}}}
   \;\longleftrightarrow\;
   \{(A, q, B) \in \NN^3 : B = 2^A \bmod 3^q,\; B \in B_{\mathrm{free}}\}.
\]
\end{theorem}

\subsubsection*{Consequence for the cylinder problem}

The cylinder family $C_m$ defined in \S\ref{sec:cylinder-family}
is, via this equivalence, the set of admissible words whose
$(A, q, B)$ triple has $n = (2^A - B)/3^q$ satisfying
$n \equiv \aSat \pmod{3^{R}}$. The cylinder accessibility cost
$\Cback(m, Q)$ is then a function purely of the
$(A, q, B)$-arithmetic, with the inverse-graph machinery of
\S\ref{sec:inverse-graph} providing the search apparatus.

\subsection{Defect}

The \emph{defect} of an admissible word is
\[
  \defect(u) \;:=\; A(u) - 2 \, q(u) \;\in\; \ZZ.
\]
The slope $A(u)/q(u)$ measures the average value of $\tau$ along
the word and the defect measures its deviation from~$2$.
The conjecture that no admissible word has $\defect(u) < -c$ for any
constant $c$ (uniformly over $u$) would imply
$\lambda_{\mathrm{asy}} \le 2$, which combined with the
inequalities below would imply a uniform bound.

\begin{lemma}[Per-edge bound; \texttt{defect\_ge\_numOnes\_minus\_numZeros}]
\label{lem:per-edge-bound}
For every admissible $u$,
\[
  \defect(u) \;\ge\; \numOnes(u) - \numZeros(u),
\]
where $\numOnes(u)$ is the number of one-edges and
$\numZeros(u)$ the number of zero-edges in $u$.
\end{lemma}

\begin{proof}
The contribution to $A$ from each edge is $\tau(\alpha) \ge 1$, so
$A(u) \ge \numOnes(u) + \numZeros(u)$. Since $q(u) = \numZeros(u)$
exactly,
$\defect(u) = A(u) - 2 q(u) \ge \numOnes - 2 \numZeros + \numZeros
            = \numOnes - \numZeros$.
The full Lean proof
(\texttt{AnalyticBarrier.lean}) tracks per-edge contributions and
chains them via the additive structure of the state machine.
\end{proof}

This is a useful but \emph{rigid} bound; the open conjecture of
\S\ref{sec:open-conjecture} sharpens it.

\subsubsection*{Defect dynamics and refined bounds}

The per-edge bound (Lemma~\ref{lem:per-edge-bound}) does not
distinguish between zero-edges by their $\tau$-value. A refinement
that does distinguish is available:

\begin{lemma}[Refined per-edge bound]
\label{lem:refined-per-edge}
For every admissible $u$:
\[
  \defect(u) \;=\; \sum_{i=1}^{|u|} \omega(e_i),
\]
where $\omega(e_i)$ is the edge weight: $\tau(\alpha_i)$ for
one-edges, $\tau(\alpha_i) - 2$ for zero-edges. The
zero-edge contributions $\tau - 2$ partition as:
\begin{center}\small
\begin{tabular}{lll}
\toprule
$\tau(\alpha)$ & $\omega = \tau - 2$ & contributing residues $\alpha$ \\
\midrule
$1$ & $-1$ & $\{2, 8\}$ \\
$2$ & $\phantom{+}0$ & $\{1, 4\}$ \\
$3$ & $+1$ & $\{5\}$ \\
$4$ & $+2$ & $\{7\}$ \\
\bottomrule
\end{tabular}
\end{center}
\end{lemma}

The zero-edges with $\tau = 1$ are the \emph{defect-decreasing}
ones; the precise conjecture (Conjecture~\ref{conj:precise})
asserts that the number of such edges is bounded by the number of
one-edges minus $m$, along every admissible $u \in C_m$.

\subsubsection*{The S203-C/D negative-run framework}

A related family of identities, formally proved in
\href{run:CollatzLean4/CollatzLean4/NegativeRuns.lean}%
{\texttt{NegativeRuns.lean}}, concerns the iterated
inverse map at negative residues. Define
\[
  L(n) \;:=\; \frac{2 n - 1}{3},
\]
the canonical odd-step inverse, and $L^k(n)$ its $k$-fold
composition. The integrality criterion is:

\begin{theorem}[S203-C: integrality of $L^k$;
\texttt{Lk\_integer\_iff}]
\label{thm:s203-C}
$L^k(n) \in \ZZ$ iff $n \equiv -1 \pmod{3^k}$.
\end{theorem}

\begin{theorem}[S203-D: explicit numerator;
\texttt{Lnum\_eq}]
\label{thm:s203-D}
Under the integrality criterion of Theorem~\ref{thm:s203-C},
\[
   3^k \cdot L^k(n) \;=\; \mathrm{Lnum}(k, n),
\]
where $\mathrm{Lnum}(k, n)$ is an explicit polynomial expression
in $n$ and $2^k$.
\end{theorem}

The S203-C/D framework is the precursor to the 3-adic lift
(Theorem~\ref{thm:s206-lift} below); it gives the algebraic
identity that, lifted from $\ZZ/3^{24}\ZZ$ to all $\ZZ/3^k\ZZ$,
underwrites the inverse-graph construction.

\subsubsection*{The 3-adic lift (S206)}

\begin{theorem}[S206 3-adic lift; \texttt{S206\_3adic\_lift}]
\label{thm:s206-lift}
The S206 identity
\[
  (2^{43} - 3^{22}) \cdot a \;\equiv\; B_{S202} \pmod{3^{24}}
\]
with solution $a = 1 + 3^{22}$ extends to a $\Zp{3}$-valued
identity. Specifically, for every $k \ge 1$ there exists
$a_k \in \ZZ/3^k\ZZ$ satisfying
\(
   (2^{43} - 3^{22}) a_k \equiv B_{S202} \pmod{3^k},
\)
and the sequence $\{a_k\}$ is Cauchy in $\Zp{3}$.
\end{theorem}

The 3-adic lift Theorem~\ref{thm:s206-lift} is the structural
result that makes the cylinder family $C_m$ well-defined for
every $m$: each $\aSat$ in \S\ref{sec:s202-word} is the
representative at precision $3^{R(m)}$.

\subsubsection*{The accumulated carry $B$ and its 2-adic structure}

The carried parameter $B(u)$ admits a natural 2-adic
interpretation. Define, for each prefix $u_{[1, k]}$:
\[
   c_k \;:=\; A(u_{[1, k]}) \quad \text{at the moment of the $i$-th $\mathtt{0}$ symbol}.
\]
Then the strict decrease $c_1 > c_2 > \cdots > c_q \ge 0$ is the
explicit exponent representation
$B(u) = \sum_{i=1}^q 3^{i-1} \cdot 2^{c_i}$.

This gives $B(u)$ a closed-form generating-function
interpretation:
\[
   B(u) \;=\; 2^{c_q} \cdot \bigl(1 + 3 \cdot 2^{c_{q-1} - c_q} + \cdots + 3^{q-1} \cdot 2^{c_1 - c_q}\bigr),
\]
which connects the inverse iteration to the bidirectional Terras
coordinate of \S\ref{sec:terras}.

\section{The S202 word and its 3-adic fixed point}\label{sec:s202-word}

\subsection{Explicit evaluation}

Define the \emph{S202 word}
\[
  \wS \;:=\; \mathtt{10100000001000010000000000} \;\in\; \{\mathtt{0}, \mathtt{1}\}^{26}.
\]

\begin{proposition}[Explicit evaluation; \texttt{Defect.lean}]
\label{prop:s202-evaluation}
\[
  \evalWord(\wS) \;=\; (n, A, q, B) \;=\;
    (251,\; 43,\; 22,\; 919\,447\,060\,349),
\]
so $\defect(\wS) = 43 - 44 = -1$.
\end{proposition}

\begin{proof}[Verification]
By \texttt{native\_decide} applied to the step-by-step state
update. The trace, starting from $(1, 0, 0, 0)$ and reading the
symbols of $\wS$ left to right, is:
\begin{center}\small
\begin{tabular}{rllrrl}
\toprule
step & sym & $n \bmod 9$ & $\tau$ & $n$ after & comment \\
\midrule
0  & ---           & $1$ & ---   & $1$ & initial \\
1  & $\mathtt{1}$  & $1$ & $2$   & $4$ & $1 \to 4$ \\
2  & $\mathtt{0}$  & $4$ & $2$   & $5$ & $4 \to (16-1)/3 = 5$ \\
3  & $\mathtt{1}$  & $5$ & $3$   & $40$ & $5 \to 40$ \\
4  & $\mathtt{0}$  & $4$ & $2$   & $53$ & $40 \to (160-1)/3 = 53$ \\
\vdots & \vdots    & \vdots & \vdots & \vdots & \vdots \\
26 & $\mathtt{0}$  & $\cdots$ & $\cdots$ & $251$ & final\\
\bottomrule
\end{tabular}
\end{center}
The full trace has $26$ rows; the formal proof avoids the manual
trace by direct \texttt{native\_decide} on the closed-form
algorithm, yielding the four-tuple $(251, 43, 22, 919\,447\,060\,349)$.
\end{proof}

\begin{remark}[Refutation of the rigid bound]
The bound $A(u) \ge 2 q(u)$, sometimes proposed as a candidate
inequality for admissible words, predicts $\defect(u) \ge 0$.
Proposition~\ref{prop:s202-evaluation} witnesses
$\defect(\wS) = -1 < 0$, refuting the rigid bound. The formal
counterexample theorem is
\texttt{not\_forall\_admissible\_triple\_high\_slope} in Lean.
This is the structural reason the present programme must
\emph{quantify} the obstruction rather than rule it out by sign.
\end{remark}

\subsubsection*{Historical note: how $\wS$ was found}

The S202 word was discovered by exhaustive search through admissible
words of length up to $30$, looking for those with the highest ratio
of zero-edges to one-edges (i.e.\ the most $\tau$-favourable). The
candidate $\wS$ emerged as the shortest such word achieving
$\defect = -1$. Independent verification by an unrelated symbolic
enumeration confirmed the parameters $(251, 43, 22, B_{S202})$.

The structural significance of $\wS$ is not its existence as a
specific word but the \emph{family} of words $\wS^k$ for $k \ge 1$:
each $\wS^k$ has $\defect(\wS^k) = -k$, and as $k \to \infty$ the
carried state $n(u)$ converges 3-adically to $\aS$. This family is
what makes the cylinder accessibility question well-posed.

\subsubsection*{Translation identity check for $\wS$}

We verify \eqref{eq:invid} for $\wS$:
\[
   3^{22} \cdot 251 \,+\, 919\,447\,060\,349
   \;=\; 31\,381\,059\,609 \cdot 251 + 919\,447\,060\,349.
\]
Computing the left product:
$31\,381\,059\,609 \cdot 251 = 7\,876\,645\,961\,859$. Adding:
$7\,876\,645\,961\,859 + 919\,447\,060\,349 = 8\,796\,093\,022\,208$.
And the right-hand side: $2^{43} = 8\,796\,093\,022\,208$. ✓

The identity holds. The Lean module
\href{run:CollatzLean4/CollatzLean4/Defect.lean}%
{\texttt{Defect.lean}}
performs this verification by \texttt{native\_decide}.

\subsection{The 3-adic limit}

Iterating $\wS^k$ produces a 3-adic Cauchy sequence of states. The
limiting value of the carried integer $n$ converges in $\Zp{3}$ to
the unique element $\aS$ satisfying
\[
  (2^{43} - 3^{22}) \cdot \aS \;\equiv\; B_{S202} \pmod{3^\infty},
\]
where $B_{S202} = 919\,447\,060\,349$ is the limit of the carried
parameter. We denote the canonical representative at precision
$3^{22m + 2}$ by $\aSat$.

\begin{theorem}[Cauchy-coherent representatives; \texttt{AS202Lift.lean}]
\label{thm:aSat-cauchy}
The values
\begin{align*}
  \aSat[m{=}1] &= 31\,381\,059\,610, \\
  \aSat[m{=}2] &= 7\,473\,280\,404\,920\,504\,074\,066, \\
  \aSat[m{=}3] &= 33\,445\,648\,520\,278\,797\,219\,953\,204\,883\,205,
\end{align*}
each satisfy the defining congruence
\(
   (2^{43} - 3^{22}) \cdot \aSat \equiv B_{S202} \pmod{3^{22m + 2}}
\)
and the Cauchy coherence
\(
   \aSat[m+1] \equiv \aSat[m] \pmod{3^{22 m + 2}}.
\)
All identities are verified by \texttt{native\_decide}.
\end{theorem}

\subsubsection*{Derivation of $\aSat[m=1]$}

For $m = 1$ the defining congruence reads
\(
  (2^{43} - 3^{22}) \cdot \aS \equiv B_{S202} \pmod{3^{24}}.
\)
Now $2^{43} - 3^{22} = 8\,796\,093\,022\,208 - 31\,381\,059\,609 = 8\,764\,711\,962\,599$.
We solve the congruence by the extended Euclidean algorithm. The
inverse of $2^{43} - 3^{22}$ modulo $3^{24} = 282\,429\,536\,481$
exists because
$\gcd(2^{43} - 3^{22}, \; 3) = \gcd(2^{43}, 3) = 1$ (as $2$ is not
divisible by $3$ and $3^{22}$ is). The inverse turns out (by
explicit modular calculation) to send $B_{S202}$ to exactly
$1 + 3^{22} = 31\,381\,059\,610$.

\subsubsection*{Derivation of $\aSat[m=2]$}

For $m = 2$ we require
\(
  (2^{43} - 3^{22}) \cdot \aSat[m{=}2] \equiv B_{S202} \pmod{3^{46}}.
\)
This is a Hensel-style lift from $\aSat[m=1]$. Writing
$\aSat[m=2] = \aSat[m=1] + 3^{24} \cdot k$ for some $k \in \ZZ/3^{22}\ZZ$,
and substituting:
\[
  (2^{43} - 3^{22}) \cdot \aSat[m=1] + (2^{43} - 3^{22}) \cdot 3^{24} \cdot k
  \;\equiv\; B_{S202} \pmod{3^{46}}.
\]
The first term differs from $B_{S202}$ by a multiple of $3^{24}$
(by the $m = 1$ congruence). Cancelling and dividing by $3^{24}$:
\[
  (2^{43} - 3^{22}) \cdot k \;\equiv\;
  \frac{B_{S202} - (2^{43} - 3^{22}) \cdot \aSat[m=1]}{3^{24}}
  \pmod{3^{22}},
\]
which determines $k$ uniquely modulo $3^{22}$. Iterating this
Hensel step gives the explicit value of $\aSat[m=2]$ stated in
Theorem~\ref{thm:aSat-cauchy}. The lifts for $m = 3$ and beyond
proceed identically.

\subsubsection*{The Cauchy property}

By construction of the Hensel lift, $\aSat[m+1] - \aSat[m]$ is
divisible by $3^{22m+2}$. Hence
$\aSat[m+1] \equiv \aSat[m] \pmod{3^{22m+2}}$
and the sequence $\{\aSat\}_{m \ge 1}$ is Cauchy in $\Zp{3}$ with
respect to the 3-adic metric. The limit
$\aS = \lim \aSat \in \Zp{3}$ is the canonical fixed point of the
$\wS$ iteration.

\begin{lemma}[3-adic valuation; \texttt{nu3\_aS\_minus\_1}]
\label{lem:nu3-aSm-minus-1}
For every $m \ge 1$,
\[
  \nuthree(\aSat - 1) \;=\; 22.
\]
\end{lemma}

\begin{proof}
For $m = 1$, $\aSat = 1 + 3^{22}$, so $\aSat - 1 = 3^{22}$,
giving $\nuthree = 22$. For $m \ge 2$, $\aSat - 1$ is a higher-degree
$3$-adic integer with the structural form
\(
  \aSat - 1 = 3^{22} \cdot u_m
\)
where $u_m \in \Zp{3}^{\times}$, i.e.\ $\gcd(u_m, 3) = 1$. The
verification is by \texttt{native\_decide} for $m \in \{2, 3\}$ on
the modular value $(\aSat - 1) / 3^{22} \bmod 3$.
\end{proof}

This lemma is decisive in Part~\ref{part:markov}: any local potential
$\psi(c \bmod 3^h)$ designed to distinguish $\aSat \bmod 3^h$ from
$1 \bmod 3^h$ must satisfy $h > \nuthree(\aSat - 1) = 22$. The
obstruction is exactly the 3-adic distance from the cylinder anchor
to the root, and it is invariant in $m$ — a deep structural fact
of the S202 fixed point.

\section{The 2-adic side: Terras coordinates and binary tensors}
\label{sec:terras}

The admissible-inverse machinery of \S\S\ref{sec:framework}--\ref{sec:s202-word}
naturally pairs with a 2-adic coordinate scaffold. We summarise this
complementary structure here, since the eventual joint analysis with
forward-orbit results (in particular Tao~2022) requires both sides.

\subsection{Forward Collatz iteration}

Recall the accelerated map
\[
  T(n) = \begin{cases} n/2 & n \text{ even}, \\
                       (3n+1)/2 & n \text{ odd}, \end{cases}
\]
along with the unaccelerated $T_{\mathrm{lento}}(n)$ which does
\emph{not} fold the $/2$ into the odd step. The accelerated $T$
identifies pairs of unaccelerated steps and is the more natural
object for parity-vector analysis.

We also use $T_{\mathrm{fast}}$, the maximally accelerated map that
applies as many factors of $2$ as possible after each odd step:
\[
   T_{\mathrm{fast}}(n) = \begin{cases}
       n/2^{\nu_2(n)} & n \text{ even}, \\
       (3n + 1)/2^{\nu_2(3n+1)} & n \text{ odd}.
   \end{cases}
\]
This three-way distinction $(T_{\mathrm{lento}}, T, T_{\mathrm{fast}})$
is recorded in
\href{run:CollatzLean4/CollatzLean4/BinaryTensor.lean}%
{\texttt{BinaryTensor.lean}}; each variant has its uses.

\subsection{Parity vectors}

For $n \ge 1$ and $k \ge 1$, define the \emph{parity vector at
depth $k$}:
\[
  P_k(n) \;:=\; (n \bmod 2,\; T(n) \bmod 2,\; T^2(n) \bmod 2,\;
                 \dots,\; T^{k-1}(n) \bmod 2)
   \;\in\; \{0, 1\}^k.
\]

\begin{theorem}[Terras bijectivity at $k \le 3$;
\texttt{terras\_bijective\_at\_3}]
\label{thm:terras-bijection}
For every $k \in \{1, 2, 3\}$, the map
\(
   P_k : \ZZ/2^k\ZZ \to \{0, 1\}^k
\)
is a bijection. Formal proof: \texttt{native\_decide} on the
finite case-distinction ($2^k \le 8$ cases).
\end{theorem}

For arbitrary $k$, Terras's theorem (1976) asserts that
$P_k$ is a bijection; the Lean formalisation of the general case is
an open item independent of the rest of the programme.

\subsubsection*{Explicit check at $k = 3$}

We tabulate $P_3$ on $\ZZ/8\ZZ$:
\begin{center}\footnotesize
\begin{tabular}{rl|rl}
\toprule
$n \bmod 8$ & $P_3(n)$ & $n \bmod 8$ & $P_3(n)$ \\
\midrule
$0$ & $(0, 0, 0)$ & $4$ & $(0, 0, 1)$ \\
$1$ & $(1, 0, 0)$ & $5$ & $(1, 1, 0)$ \\
$2$ & $(0, 1, 1)$ & $6$ & $(0, 1, 0)$ \\
$3$ & $(1, 0, 1)$ & $7$ & $(1, 1, 1)$ \\
\bottomrule
\end{tabular}
\end{center}
The 8 parity vectors are distinct, confirming bijectivity at $k = 3$.

\subsection{Accumulated carry}

Define the \emph{accumulated carry vector} $C_k(n) \in \NN^k$ by
$C_k(n)_i = $ the integer carry contributed by the $i$-th step of
$T_{\mathrm{lento}}$ applied to $n$. Specifically:
\begin{itemize}
  \item if $T_{\mathrm{lento}}^{i-1}(n)$ is even: $C_k(n)_i = 0$;
  \item if $T_{\mathrm{lento}}^{i-1}(n)$ is odd:
        $C_k(n)_i = 1$ (the contributed $+1$ in $3n + 1$).
\end{itemize}

The accumulated carry decomposes the orbit into a forward sum:
\[
   T_{\mathrm{lento}}^k(n) \cdot 2^{w_k(n)}
   \;=\; 3^{p_k(n)} \, n + \mathrm{Carry}_k(n),
\]
where $w_k(n) = \#\{i : P_i(n) = 0\}$ is the number of even
steps, $p_k(n) = \#\{i : P_i(n) = 1\}$ is the number of odd
steps, and $\mathrm{Carry}_k$ is a polynomial in $\{2^{\nu_2(\cdot)}\}$
along the trajectory.

This is precisely the \emph{forward} translation identity, dual to
\eqref{eq:invid} for the inverse iteration. The two identities
relate the same arithmetic content seen from opposite directions.

\subsection{The 4-2-1 attractor}

\begin{theorem}[4-2-1 cycle; \texttt{T\_421\_cycle}]
\label{thm:421}
In the accelerated map,
\[
   T(4) = 2, \quad T(2) = 1, \quad T(1) = 2.
\]
Hence $\{1, 2\}$ is a $2$-cycle of $T$, with $4$ entering it.
Equivalently, $\{1, 2, 4\}$ is the canonical $3$-cycle of
$T_{\mathrm{lento}}$ (after unaccelerating).
Formal proof: \texttt{native\_decide}.
\end{theorem}

The Collatz conjecture asserts that this is the \emph{only}
cycle of $T$ on $\NN_{> 0}$.

\subsection{Reachability table}

The S202 programme uses a reachability table for the inverse
admissibility relation: which integers $n$ are in the image of
the inverse branch from $1$ after $k$ steps? This is formalised
in
\href{run:CollatzLean4/CollatzLean4/Reachability.lean}%
{\texttt{Reachability.lean}}.

\begin{theorem}[Explicit reachability;
\texttt{InImageTheta\_table}]
\label{thm:reachability}
For each $n \in \{1, 4, 5, 11, 13, 16, 17, 40\}$, there is an
explicit admissible word
$u_n$ with $\evalWord(u_n).n = n$ and $\evalWord(u_n).A$ minimised
among such words. These witnesses are recorded as a finite table
and verified by \texttt{native\_decide}.
\end{theorem}

The integers $n = 2$ and $n = 8$ are \emph{not} in the image at
any short depth. This is the empirical observation that the
admissibility section $\InX$ (\S\ref{sec:framework}) cleanly
excludes residues $\{0, 3, 6\} \bmod 9$; $n = 8 \equiv 8 \in \InX$
but is not reachable as an image, and $n = 2 \equiv 2 \in \InX$
similarly. The full reachability picture is a refinement of the
section $\InX$.

\subsection{The role of the 2-adic side in the joint analysis}

The eventual joint statement, conditional on
Conjecture~\ref{conj:precise} (\S\ref{sec:open-conjecture}),
would be:
\[
  \text{(every admissible cover has slope $\lambda < 2$)}
  \;\;\;\wedge\;\;\;
  \text{(Tao 2022 bounded-value)}
  \;\;\Longrightarrow\;\;
  \text{(forward Collatz orbits attain low values quantitatively)}.
\]
The 2-adic side is the language in which the right-hand statement
is naturally framed (parity vectors and carry accumulations are
2-adic objects); the 3-adic side is the language for the
admissible-inverse machinery. The joint reformulation requires both
languages to be available, which is why the formal Lean
development includes both
\texttt{BinaryTensor.lean} (this section) and
\texttt{AdmissibleBasic.lean} (\S\ref{sec:framework}).

% =====================================================================
\part{Cylinder accessibility}\label{part:cylinders}
% =====================================================================

\section{The S202 cylinder family}\label{sec:cylinder-family}

\subsection{Definition}

For each $m \ge 1$, set $R(m) := 22m + 2$, and define the
\emph{S202 cylinder of depth $m$} as
\[
  C_m \;:=\;
  \Bigl\{ \, u \in \{\mathtt{0}, \mathtt{1}\}^* \;\big|\;
       u \text{ admissible}, \;
       \evalWord(u).n \;\equiv\; \aSat \pmod{3^{R(m)}}
  \Bigr\}.
\]
By construction, $u = \wS^k$ (or any prefix thereof) lies in $C_m$
for every $k \ge m$. Hence $C_m$ is non-empty for every $m$.

\subsection{Asymptotic content}

The relevance of $C_m$ for the Collatz dynamics is the following.

\begin{theorem}[S211 subcritical accessibility; \texttt{S211\_S202\_subcritical}]
\label{thm:s211-subcritical}
Let $u$ be an admissible word with $u \in C_m$ and $\defect(u) < m$.
Then the iterated concatenation $u \cdot \wS^m$ is admissible and
satisfies
\[
  A(u \cdot \wS^m) \;<\; 2 \, q(u \cdot \wS^m).
\]
\end{theorem}

\begin{proof}[Proof sketch]
By the modular alignment, the cylinder hypothesis allows the
concatenation to be admissible (the residue $\evalWord(u).n \bmod 9$
matches the prefix-residue of $\wS^m$ via the 3-adic fixed-point
property). The defect-additivity formula gives
$\defect(u \cdot \wS^m) = \defect(u) + m \cdot \defect(\wS) =
\defect(u) - m < 0$, hence $A < 2q$. Formalised in
\href{run:CollatzLean4/CollatzLean4/Concatenation.lean}%
{\texttt{Concatenation.lean}}.
\end{proof}

Theorem~\ref{thm:s211-subcritical} is conditional: it requires
\emph{exhibition} of an admissible $u \in C_m$ with $\defect(u) < m$.
The barrier programme tries to rule this out by establishing the
quantitative lower bound that defines the present problem.

\subsubsection*{Defect additivity}

The proof of Theorem~\ref{thm:s211-subcritical} rests on the
following identity, which is itself a consequence of the
translation identity.

\begin{lemma}[Defect additivity; \texttt{defect\_additive}]
\label{lem:defect-additive}
For concatenable admissible words $u_1, u_2$ (the result of
running $u_2$ from the state $\evalWord(u_1)$ is well-defined),
\[
   A(u_1 \cdot u_2) = A(u_1) + A(u_2),
   \qquad
   q(u_1 \cdot u_2) = q(u_1) + q(u_2),
\]
hence
\[
   \defect(u_1 \cdot u_2) \;=\; \defect(u_1) + \defect(u_2).
\]
\end{lemma}

\begin{proof}
By induction on the length of $u_2$, each step adding $\tau$ to
$A$ (one-edge) or $\tau$ to $A$ and $1$ to $q$ (zero-edge). The
contributions accumulate, and the defect sums by definition.
\end{proof}

In particular, $\defect(\wS^m) = m \cdot \defect(\wS) = -m$.
Hence if $\defect(u) < m$ on some $u \in C_m$, then
$\defect(u \cdot \wS^m) < 0$, which gives the conclusion of
Theorem~\ref{thm:s211-subcritical}.

\subsubsection*{Concatenability of $\wS$ to a $C_m$-anchored word}

The condition that $u \cdot \wS^m$ is admissible requires that the
final state $\evalWord(u)$ has $n \bmod 9 \in \InX$ \emph{and}
that the next $\mathtt{1}$ or $\mathtt{0}$ symbol of $\wS$ does
not trip the admissibility condition. The first is automatic by
the closure property of $\InX$ under the iteration; the second is
where the cylinder hypothesis $u \in C_m$ is essential.

Specifically, $u \in C_m$ means
$\evalWord(u).n \equiv \aSat \pmod{3^{R(m)}}$. By
construction of $\aSat$ as a fixed point of the $\wS$-iteration,
the first $m$ copies of $\wS$ following $u$ are guaranteed to
keep $n \bmod 9$ in $\InX$ at every step. Beyond that, more
care is required, but for $\wS^m$ specifically the iteration is
controlled by the fixed-point property.

\subsubsection*{S213: cylinder stability}

A finer question, addressed in
\href{run:CollatzLean4/CollatzLean4/CylinderStability.lean}%
{\texttt{CylinderStability.lean}}, concerns the
\emph{$\tau$-trajectory} of $u \cdot \wS^k$ for arbitrary $k$. Define
the $\tau$-trajectory of $u$ as the sequence of
$\tau(\evalWord(u_{[1, i]}).n \bmod 9)$ for $i = 1, \dots, |u|$.

\begin{theorem}[$\tau$-stability under branch.one;
\texttt{tau\_stability\_branch\_one}]
\label{thm:tau-stability}
Suppose two admissible words $u, u'$ satisfy
$\evalWord(u).n \equiv \evalWord(u').n \pmod{3^{R}}$ for some
$R \ge 22m + 2$. Then the $\tau$-trajectories of $u \cdot \mathtt{1}^k$
and $u' \cdot \mathtt{1}^k$ agree for all $k \ge 0$.
\end{theorem}

\begin{theorem}[$\tau$-stability under branch.zero loses one digit;
\texttt{tau\_stability\_branch\_zero}]
\label{thm:tau-stability-zero}
Under the same hypothesis, the $\tau$-trajectories of
$u \cdot \mathtt{0}^k$ and $u' \cdot \mathtt{0}^k$ agree as long
as $k \le R - 22m - 1$.
\end{theorem}

The asymmetry between branch.one (preserves all 3-adic precision)
and branch.zero (loses one digit per step) is the technical heart
of S213. It explains why $C_m$ must be defined at precision
$3^{22m + 2}$ rather than just $3^{22m}$: the extra two digits
are needed to absorb the $\tau$-stability budget for short
zero-runs at the end of $\wS^k$.

\subsection{Cylinder accessibility cost}

\begin{definition}[Cylinder accessibility cost]
\label{def:Cback}
For $m \ge 1$, $Q \ge 0$, define
\[
  \Cback(m, Q) \;:=\;
  \inf \Bigl\{ \, \defect(u) \;\big|\;
       u \text{ admissible}, \;
       q(u) \le Q, \;
       u \in C_m
  \Bigr\}.
\]
\end{definition}

\noindent The conjectural goal is that $\Cback(m, Q)$ grows at least
linearly in $m$ (uniformly in $Q$): if so,
Theorem~\ref{thm:s211-subcritical} cannot be triggered for any $m$.

\section{The inverse cylinder graph}\label{sec:inverse-graph}

\subsection{Construction}

Fix $m$ and let $R := 22m + 2$. The \emph{inverse cylinder graph}
$\InvG_{m, Q}$ has:

\paragraph{Vertices.}
Pairs $v = (j, c)$ with
$j \in \{0, 1, \dots, Q\}$ and
$c \in \ZZ/3^{R+j}\ZZ$, the latter taken with canonical
representative $0 \le c < 3^{R+j}$.

\paragraph{Start vertex.}
\[
   s \;:=\; (0,\; \aSat).
\]

\paragraph{Goal vertices.}
$(j, c)$ with $c \equiv 1 \pmod{3^{R+j}}$, i.e.\ canonically
$c = 1$. We write $g_j$ for the unique canonical goal at level $j$.

\paragraph{One-edges.}
For each $\alpha \in \InX$ and vertex $v = (j, c)$, there is an edge
$v \to v'$ iff:
\begin{enumerate}
  \item $j(v') = j$,
  \item $c \equiv 2^{\tau(\alpha)} \cdot v'.c \pmod{3^{R+j}}$,
  \item $v'.c \bmod 9 = \alpha$.
\end{enumerate}
Defect weight: $\omega(e) = \tau(\alpha) \in \{1, 2, 3, 4\}$.

\paragraph{Zero-edges.}
For each $\alpha \in \InX$, vertex $v = (j, c)$ with $j < Q$, there
is an edge $v \to v'$ iff:
\begin{enumerate}
  \item $j(v') = j + 1$,
  \item $3 c + 1 \equiv 2^{\tau(\alpha)} \cdot v'.c \pmod{3^{R + j + 1}}$,
  \item $v'.c \bmod 9 = \alpha$.
\end{enumerate}
Defect weight: $\omega(e) = \tau(\alpha) - 2 \in \{-1, 0, 1, 2\}$.

\paragraph{Canonical determination.}
For each edge type, the canonical $v'.c$ is uniquely determined by
$v.c$ and $\alpha$. The inverse-of-$2^{\tau(\alpha)}$ mod the
appropriate $3^k$ is computable by repeated squaring; the
practical code uses
\(\mathrm{invMod2}(M) := (M+1)/2\),
the closed-form inverse of $2$ modulo any odd $M$.

\subsubsection*{Why six edges per direction per vertex}

At a vertex $v = (j, c)$ with $j < Q$, the candidate
$\alpha$-values for outgoing edges range over
$\InX = \{1, 2, 4, 5, 7, 8\}$ for both one-edges and zero-edges,
yielding at most twelve outgoing edges in total. However, the
residue check
$v'.c \bmod 9 = \alpha$
selects \emph{at most one} edge per type-and-$\alpha$ pair: the
formula $v'.c = (2^\tau)^{-1} \cdot (\text{numerator}) \bmod 3^{R+j(v')}$
is a single deterministic value, so either it lands in the residue
class $\alpha \bmod 9$ or it does not. In practice, roughly
$2$--$3$ outgoing edges per vertex satisfy the residue check on
average.

The total branching factor is therefore bounded by $12$ but
typically much smaller, which is why the BF closure
(\S\ref{sec:s216-closure}) is tractable for moderate $(m, Q)$.

\subsubsection*{Explicit example of an inverse edge}

Take $m = 1$ ($R = 24$), $j = 0$, $v.c = \aSat = 31\,381\,059\,610$.
For the one-edge with $\alpha = 1$ ($\tau = 2$):
\begin{itemize}
  \item modulus: $M = 3^{24} = 282\,429\,536\,481$;
  \item compute $4^{-1} \bmod M$: by extended Euclidean,
        $4 \cdot 211\,822\,152\,361 \equiv 1 \pmod{M}$, so
        $4^{-1} = 211\,822\,152\,361$;
  \item compute $v'.c = 211\,822\,152\,361 \cdot 31\,381\,059\,610 \bmod M$;
  \item check $v'.c \bmod 9 = 1$.
\end{itemize}
If all three conditions hold, the edge $v \to (0, v'.c)$ exists with
defect weight $\omega = 2$. The Python implementation in
\texttt{tools/s202\_engine.py:outgoing\_edges} performs this for all
six $\alpha$ values, for both edge types, and returns the list of
admissible outgoing edges.

\subsection{Cardinality and growth}

\begin{lemma}[Branching factor]
\label{lem:branching}
The out-degree of any vertex $v$ in $\InvG_{m, Q}$ is at most $12$,
and typically averages between $2$ and $3$. The relevant region
(states satisfying the optimistic-credit cut of
\S\ref{sec:s216-closure}) grows roughly as $\Theta(3.7^{Q})$ for
fixed $m$, based on empirical fits.
\end{lemma}

\begin{proof}[Empirical justification]
Tabulating the state counts of the seven frontiers
(\S\ref{sec:formal-status}):
\begin{center}\footnotesize
\begin{tabular}{rrr}
\toprule
$Q$ & states ($m = 1$) & ratio to previous $Q$ \\
\midrule
$3$  & $35$       & ---  \\
$5$  & $462$      & $13.2$ (over $\Delta Q = 2$, factor $\sqrt{13.2} \approx 3.63$) \\
$8$  & $24\,310$  & $52.6$ (over $\Delta Q = 3$, factor $52.6^{1/3} \approx 3.75$) \\
$10$ & $352\,716$ & $14.5$ (over $\Delta Q = 2$, factor $\sqrt{14.5} \approx 3.81$) \\
\bottomrule
\end{tabular}
\end{center}
The growth factor per unit of $Q$ is around $3.7$, consistent
across the four data points.
\end{proof}

The implication for compute budgeting: $Q = 15$ at $m = 1$ requires
$\sim 5 \times 10^7$ states ($\sim 5$ minutes of laptop CPU);
$Q = 18$ requires $\sim 2 \times 10^9$ states (~hours on a single
machine, but well within reach of moderate cluster compute or
AlphaEvolve).

\subsection{Path weight and defect}

A directed path $p = v_0 \to v_1 \to \cdots \to v_k$ in
$\InvG_{m, Q}$ has \emph{defect weight}
\[
  w(p) \;:=\; \sum_{i=0}^{k-1} \omega(v_i \to v_{i+1}) \;\in\; \ZZ.
\]
We also define the \emph{combinatorial weight}
\[
  \kappa(e) \;:=\;
  \begin{cases}
    +1 & e \text{ is a one-edge}, \\
    -1 & e \text{ is a zero-edge with } \tau(\alpha) = 1, \\
    \phantom{+}0 & e \text{ is a zero-edge with } \tau(\alpha) \ge 2.
  \end{cases}
\]
This is the \emph{precise} variant; the \emph{coarse} variant gives
every zero-edge weight $-1$. The combinatorial weight is the
arithmetic relevant to Part~\ref{part:markov}; the defect weight is
the arithmetic relevant to Part~\ref{part:barriers}.

\subsection{Path-word correspondence (forward)}

\begin{theorem}[Forward correspondence; \texttt{S212\_forward\_general}]
\label{thm:s212-forward}
Let $u$ be an admissible word with $q(u) \le Q$ and
$(\evalWord u).n \equiv \aSat \pmod{3^R}$, i.e.\ $u \in C_m$.
There exists a directed path
$p_u : s \to v_1 \to \cdots \to v_k$ in $\InvG_{m, Q}$ with
\begin{itemize}
  \item $j(v_k) = q(u)$ and $v_k = g_{q(u)} = (q(u), 1)$;
  \item defect weight $w(p_u) = \defect(u)$;
  \item every intermediate vertex canonical
        (i.e.\ $0 \le v_i.c < 3^{R + j(v_i)}$).
\end{itemize}
\end{theorem}

\begin{proof}[Proof sketch]
By induction on the word length, prepending one inverse edge per
symbol. For the base case $u = \varepsilon$, take the trivial path
$s \to s$ (length $0$); the cylinder hypothesis with $q = 0$ forces
$\evalWord(u).n = \aSat$ which is exactly $s.c$. For the inductive
step on $u = u' \cdot b$, the induction provides a path for $u'$
ending at the post-$u'$ vertex; appending an inverse edge of type
$b$ extends the path to $u$. Admissibility ensures the residue
match. Canonicality is preserved because each appended vertex is
constructed as $(j, x \bmod 3^{R + j})$ for an explicit $x$.

The full Lean proof is by structural induction on the prefix in
\href{run:CollatzLean4/CollatzLean4/S212Forward.lean}%
{\texttt{S212Forward.lean}}, with eight subcases corresponding to
the eight $(\text{symbol}, \alpha)$ combinations.
\end{proof}

The corresponding \emph{reverse} correspondence (every path in
$\InvG_{m, Q}$ yields at least one admissible word) is recorded as
the open implementation item \texttt{S212\_path\_word\_correspondence}
and is not required for the closure-to-barrier results below; only
the forward direction is needed.

% =====================================================================
\part{Barriers via closure certificates}\label{part:barriers}
% =====================================================================

\section{The optimistic-credit closure}\label{sec:s216-closure}

\subsection{The negativity issue}

A naive barrier search would relax labels in $\InvG_{m, Q}$ under
the cut ``$d < T$'': enumerate $(v, d)$ pairs reachable from $s$
with cumulative defect weight $d$, prune any pair with $d \ge T$,
and check that no goal $g$ appears with $d < T$.

This is \emph{unsound} for $\InvG_{m, Q}$. Zero-edges with
$\tau(\alpha) = 1$ (i.e.\ $\alpha \in \{2, 8\}$) carry weight
$\omega = -1$. A path may temporarily accumulate $d \ge T$ and
subsequently descend below $T$ via a negative-weight step; the
naive cut prunes precisely these rescuable prefixes.

\subsection{The credit bound}

The key observation is that negative-weight steps are exactly the
$j$-incrementing zero-edges with $\tau = 1$, and these are
bounded in number along any prefix by the remaining $Q - j$ zero-edge
budget.

\begin{lemma}[Suffix weight bound; \texttt{suffix\_weight\_bound}]
\label{lem:suffix-bound}
For every directed path $p : v \to^{*} t$ in $\InvG_{m, Q}$,
\[
  w(p) \;\ge\; -\bigl(Q - j(v)\bigr).
\]
\end{lemma}

\begin{proof}
Each edge $v_i \to v_{i+1}$ in $\InvG_{m, Q}$ has weight
$\omega \ge j(v_i) - j(v_{i+1})$: indeed, one-edges have
$\omega = \tau \ge 1$ and preserve $j$, hence
$\omega \ge 0 = j(v_i) - j(v_{i+1})$; zero-edges have
$\omega \ge -1$ and $j(v_{i+1}) = j(v_i) + 1$, so
$\omega \ge -1 = j(v_i) - j(v_{i+1})$. Telescoping along the path:
\[
  w(p) \;=\; \sum_{i} \omega(v_i \to v_{i+1})
       \;\ge\; \sum_i \bigl(j(v_i) - j(v_{i+1})\bigr)
       \;=\; j(v_0) - j(v_k).
\]
Since the path is confined to $\InvG_{m, Q}$, $j(v_k) \le Q$.
Hence $w(p) \ge j(v_0) - Q = -(Q - j(v_0))$.
\end{proof}

Define the \emph{credit} of a vertex as
\[
  \credit(Q, v) \;:=\; Q - j(v),
\]
and call a pair $(v, d)$ \emph{relevant} (with respect to
$(Q, T)$) iff
\[
  d - \credit(Q, v) \;<\; T,
\]
i.e.\ the most optimistic completion of $(v, d)$ to a goal can still
fall below the target $T$.

\begin{lemma}[Prefix relevance; \texttt{prefix\_relevance}]
\label{lem:prefix-relevance}
Let $p : s \to^{*} g$ be a directed path of cumulative defect weight
$w(p) < T$, with prefix vertex sequence
$v_0 = s, v_1, \dots, v_k = g$ and prefix weights
$d_i = w(v_0 \to^{*} v_i)$. Then every $(v_i, d_i)$ is relevant.
\end{lemma}

\begin{proof}
Apply Lemma~\ref{lem:suffix-bound} to the suffix $v_i \to^* g$:
$w(v_i \to^{*} g) \ge -\credit(Q, v_i)$. Adding:
\[
  d_i + w(v_i \to^{*} g) \;=\; w(p) \;<\; T,
\]
hence
$d_i \;<\; T - w(v_i \to^{*} g) \;\le\; T + \credit(Q, v_i)$,
which rearranges to $d_i - \credit(Q, v_i) < T$.
\end{proof}

\subsection{Closure certificate}

\begin{definition}[S216 closure certificate]
\label{def:s216-cert}
A partial function
\(
   \mathrm{dist} : V_{m, Q} \to \ZZ \cup \{\bot\}
\)
(with $\bot$ meaning ``undefined''), the \emph{dist table}, is an
\emph{S216 closure certificate for $(m, Q, T)$} if:
\begin{itemize}
  \item[(C1)] $\mathrm{dist}(s) = 0$.
  \item[(C2)] For every goal $g$ in the support of $\mathrm{dist}$,
        $\mathrm{dist}(g) \ge T$.
  \item[(C3)] For every $v$ in the support with
        $\relevant(v, \mathrm{dist}(v))$, and every outgoing edge
        $v \to v'$ of weight $\omega$ with $j(v') \le Q$ and
        $\relevant(v', \mathrm{dist}(v) + \omega)$, the value
        $\mathrm{dist}(v')$ is in the support and satisfies
        $\mathrm{dist}(v') \le \mathrm{dist}(v) + \omega$.
\end{itemize}
\end{definition}

\subsection{Closure-to-barrier}

\begin{theorem}[Closure-to-barrier; \texttt{S216\_barrier\_for\_words\_outgoing}]
\label{thm:closure-to-barrier}
If an S216 closure certificate for $(m, Q, T)$ exists, then
\[
  \Cback(m, Q) \;\ge\; T.
\]
Equivalently, every admissible word $u$ with $u \in C_m$ and
$q(u) \le Q$ satisfies $\defect(u) \ge T$.
\end{theorem}

\begin{proof}
Suppose, for contradiction, that there exists an admissible word
$u \in C_m$ with $q(u) \le Q$ and $\defect(u) < T$.
By Theorem~\ref{thm:s212-forward}, there is a path
$p_u : s = v_0 \to v_1 \to \cdots \to v_k = g_{q(u)}$ in
$\InvG_{m, Q}$ with $w(p_u) = \defect(u) < T$.

We show by induction on $i \in \{0, 1, \dots, k\}$:
\[
  \mathrm{dist}(v_i) \text{ is defined and } \mathrm{dist}(v_i) \le d_i,
  \tag{$\star_i$}\label{eq:dist-ind}
\]
where $d_i = w(v_0 \to^{*} v_i)$.

\paragraph{Base $i = 0$.}
By (C1), $\mathrm{dist}(v_0) = \mathrm{dist}(s) = 0 = d_0$.

\paragraph{Inductive step $i \to i+1$.}
Assume $(\star_i)$ holds. By Lemma~\ref{lem:prefix-relevance} applied
to $p_u$, the pair $(v_i, d_i)$ is relevant. The inductive hypothesis
gives $\mathrm{dist}(v_i) \le d_i$, so $(v_i, \mathrm{dist}(v_i))$ is
also relevant. By the same lemma applied to $(v_{i+1}, d_{i+1})$,
the pair $(v_{i+1}, d_{i+1})$ is relevant; and since
$\mathrm{dist}(v_i) + \omega_i \le d_{i+1}$ where
$\omega_i = w(v_i \to v_{i+1})$, the pair
$(v_{i+1}, \mathrm{dist}(v_i) + \omega_i)$ is relevant too.

By (C3), $\mathrm{dist}(v_{i+1})$ is in the support and
$\mathrm{dist}(v_{i+1}) \le \mathrm{dist}(v_i) + \omega_i
                          \le d_i + \omega_i = d_{i+1}$,
which is $(\star_{i+1})$.

\paragraph{Conclusion.}
Applying ($\star_k$) at the goal $v_k = g_{q(u)}$:
$\mathrm{dist}(g_{q(u)}) \le d_k = w(p_u) = \defect(u) < T$,
contradicting (C2).
\end{proof}

\subsubsection*{Minimality of the relevance rule}

\begin{proposition}[Optimal pruning]
\label{prop:optimal-pruning}
The relevance predicate $\relevant(v, d) := d - (Q - j(v)) < T$ is
the \emph{unique} minimal sound pruning rule of the form
$d - f(v) < T$ for a function $f : V \to \ZZ$ that depends only on
$v$.
\end{proposition}

\begin{proof}[Sketch]
By Lemma~\ref{lem:suffix-bound}, the suffix weight from $v$ is at
least $-(Q - j(v))$, so a pair $(v, d)$ with
$d - (Q - j(v)) \ge T$ cannot complete to a path of weight $< T$.
Hence $\relevant(v, d) := d - (Q - j(v)) < T$ is sound. Any
$f(v) > Q - j(v)$ at some $v$ would over-prune (cut a label that
can still complete); any $f(v) < Q - j(v)$ over-keeps (relevant
but unnecessary). The choice $f = Q - j$ is therefore minimal in
both directions.

This characterisation explains why the naive cut $d < T$
(corresponding to $f \equiv 0$) is unsound for the inverse cylinder
graph: it ignores the residual zero-edge budget $Q - j$, which is
exactly the credit available for further descent.
\end{proof}

\subsubsection*{What the closure certificate is, abstractly}

A closure certificate (Definition~\ref{def:s216-cert}) is the
\emph{dual} witness to a non-existence statement: the
non-existence of a cheap path. By the LP duality of shortest paths
on weighted graphs, such non-existence is equivalent to the
existence of a feasible labelling (the dist table) satisfying the
relaxation invariant.

The closure-to-barrier theorem
(Theorem~\ref{thm:closure-to-barrier}) is the LP weak duality
specialised to the optimistic-credit relaxation. The
\emph{strong} duality (existence of an optimal labelling) is
provided by the BF algorithm, which constructively produces the
dist table.

\subsubsection*{Comparison to standard Bellman--Ford}

Classical Bellman--Ford runs over a graph $(V, E)$ with edge
weights $\omega: E \to \ZZ$, computing $\mathrm{dist}(v)$ = the
minimum weight path from a source $s$ to $v$. The naive
correctness theorem requires either non-negative edge weights or
absence of negative cycles. The S216 closure adapts this to a
context where negative edges are present (zero-edges with
$\tau = 1$, weight $-1$) but the residual budget $Q - j$ provides
a problem-specific termination bound. The optimistic credit cut is
the analogue of the Johnson's algorithm
re-weighting that makes classical BF tractable on negative-edge
graphs, but more specific: it uses the structure of the graph
(monotone $j$) rather than a uniform re-weighting.

\subsection{Bool-decidable certificate}

The abstract formulation of Definition~\ref{def:s216-cert} involves
the predicate $\InvEdge_R$ as a $\forall$-quantifier in (C3). For
machine verification, we expose an explicit enumerator
\(
   \mathrm{outgoingEdges} : V_{m, Q} \to \mathrm{List}\,(V \times \ZZ),
\)
returning all outgoing edges of $v$ with weights. The Bool-decidable
form of (C3) iterates over \texttt{outgoingEdges($v$)} and replaces
the quantifier by an explicit fold. Three completeness lemmas
(\texttt{outgoingEdges\_sound},
 \texttt{outgoingEdges\_complete\_one},
 \texttt{outgoingEdges\_complete\_zero}) certify that the enumerator
is sound and complete with respect to the abstract predicate over
canonical paths.

The flagship theorem, after this Bool-Prop bridge, is:

\begin{lstlisting}
theorem S216_barrier_for_words_outgoing
    {m Q : Nat} {T : Int}
    (cert : S216Cert_outgoing (22 * m + 2) Q T)
    (h_start_match : cert.start = InvStart m)
    (h_start_can  : cert.start.c < 3 ^ (22 * m + 2 + cert.start.j))
    (u : Word)
    (h_q_bound : (evalWord u).q ≤ Q)
    (h_cyl : (evalWord u).n % (3 ^ (22 * m + 2))
              = aS202 % (3 ^ (22 * m + 2))) :
    T ≤ defect (evalWord u)
\end{lstlisting}

\section{The certificate-generation algorithm}\label{sec:algorithm}

\begin{algorithm}[H]
\caption{S216 BF closure with optimistic credit
         (\texttt{tools/s216\_engine.py})}
\label{alg:s216-bf}
\begin{algorithmic}[1]
\State $\mathrm{dist}[s] \gets 0$;\quad $Q_{\mathrm{queue}} \gets [s]$;\quad $\mathrm{inq} \gets \{s\}$
\While{$Q_{\mathrm{queue}}$ non-empty}
  \State $v \gets$ pop-front of $Q_{\mathrm{queue}}$;\quad $d \gets \mathrm{dist}[v]$
  \If{$d - (Q - j(v)) \ge T$}
    \State \textbf{continue}    \Comment{prune non-relevant}
  \EndIf
  \If{$v$ is a goal \textbf{and} $d < T$}
    \State \Return $(\textsc{useful\_path}, v, d)$
       \Comment{positive witness}
  \EndIf
  \For{$(v', \omega) \in \mathrm{outgoingEdges}(v)$}
    \State $d' \gets d + \omega$
    \If{$v'$ is a goal \textbf{and} $d' < T$}
      \State \Return $(\textsc{useful\_path}, v', d')$
    \EndIf
    \If{$d' - (Q - j(v')) \ge T$}
      \State \textbf{continue}    \Comment{non-relevant target}
    \EndIf
    \If{$\mathrm{dist}[v']$ undefined \textbf{or} $d' < \mathrm{dist}[v']$}
      \State $\mathrm{dist}[v'] \gets d'$;\quad enqueue $v'$
    \EndIf
  \EndFor
\EndWhile
\State \Return $(\textsc{barrier}, \mathrm{dist})$    \Comment{closed certificate}
\end{algorithmic}
\end{algorithm}

\paragraph{Correctness.}
On termination with status \textsc{barrier}, the resulting
$\mathrm{dist}$ satisfies (C1) by initialisation, (C2) by the
goal-detection branch, and (C3) by the relaxation invariant of the
SPFA-style queue. The pruning by the optimistic cut is sound by
Lemma~\ref{lem:prefix-relevance}: pruned labels cannot complete to a
path of weight $< T$.

\paragraph{Termination.}
The relevant region is finite: only pairs $(v, d)$ with
$d < T + \credit(Q, v) = T + (Q - j(v))$ are kept; since $j$ is
bounded by $Q$ and $c$ by $3^{R + j}$, this is a finite set.

\paragraph{Bool-decidable Lean lift.}
The Python output is a JSON list of
$\{j, c, \mathrm{dist}\}$ records. The exporter
\texttt{tools/cert\_to\_lean\_s216.py} emits Lean source that
hardcodes this list, defines the partial function via
\texttt{List.find?}, formulates the three checks (C1, C2, C3) as
\texttt{Bool}-decidable predicates, and proves them via
\texttt{native\_decide}. The Bool-to-Prop adapter completes the
chain to Theorem~\ref{thm:closure-to-barrier} machine-verified.

\paragraph{Complexity.}
The relevant region size at $(m, Q)$ scales empirically as
$\Theta(3.7^Q)$ for fixed $m$ (Lemma~\ref{lem:branching}). Per
iteration, the BF SPFA-style relaxation does $O(|E|)$ work where
$|E|$ is the number of observed edges, bounded above by $12 |V|$.
The number of relaxations per vertex is bounded by $T + (Q - j)$,
which for typical $T = m \le 3$ and moderate $Q$ stays small. Total
empirical wall-clock on a laptop is $\le 1$ second up to
$|V| \approx 10^4$, and scales to a few seconds at $|V| \approx 10^5$.

\paragraph{Chunked Lean emission.}
For tables with more than $800$ states, the Lean exporter emits
the dist list in chunks of $500$ entries
(\texttt{CHUNK\_THRESHOLD = 800}, \texttt{CHUNK\_SIZE = 500} in
\texttt{cert\_to\_lean\_s216.py}). This is needed because Lean 4
elaborates very long literal lists with quadratic time complexity;
chunking restores linear elaboration. The $(m, Q) = (1, 10)$
certificate has $706$ chunks of $\le 500$ rows each, concatenated
in groups of $6$ per line for readability. The resulting Lean
module is about $13$ MB; building it requires
\texttt{set\_option maxHeartbeats 0} and
\texttt{set\_option maxRecDepth 500000}.

\subsubsection*{Native\_decide trust base}

The Lean tactic \texttt{native\_decide} compiles a decidable
proposition to C, compiles the C code, runs it, and accepts the
result as a proof. The added axioms (visible via
\texttt{\#print axioms}) are:
\begin{itemize}
  \item \texttt{Lean.ofReduceBool}: trust that the native
        evaluation of a Boolean reduces to the same truth value
        as the kernel reduction;
  \item \texttt{Lean.ofReduceNat}: same for \texttt{Nat} values.
\end{itemize}
These are part of the standard Lean 4 trust base, not added by
this development. Their semantics is documented in the Lean
reference manual.

\section{Formal verification status}\label{sec:formal-status}

\subsection{Build statistics}

\begin{itemize}
  \item Toolchain: \texttt{leanprover/lean4:v4.30.0-rc2}; Mathlib
        pinned by \texttt{lake-manifest.json}.
  \item Build jobs: \textbf{$8\,425+$}, including the four explicit
        terminal barriers listed below.
  \item \texttt{sorry} count: \textbf{$0$} anywhere in the chain
        (verified by \texttt{grep -r sorry CollatzLean4/}).
  \item Added axioms beyond Lean core + Mathlib: \textbf{$0$},
        verified by inspection of \texttt{\#print axioms} for each
        headline theorem.
        \texttt{Lean.ofReduceBool} and \texttt{Lean.ofReduceNat}
        appear in certificates that use \texttt{native\_decide} —
        these are part of the trust base of native evaluation
        and are not added axioms of the development.
  \item Continuous integration:
        \href{run:.github/workflows/lean.yml}%
             {\texttt{.github/workflows/lean.yml}}
        runs \texttt{lake build} and the axiom inspector on every
        push; build badge in repository README.
\end{itemize}

\subsection{Module map}

The development is organised into the following modules. Path
prefixes \texttt{CollatzLean4/CollatzLean4/} are omitted.

\begin{center}\footnotesize
\begin{tabular}{ll}
\toprule
Module & Content \\
\midrule
\texttt{AdmissibleBasic} & State machine, translation identity, resonance \\
\texttt{TranslationSets} & $B_{\mathrm{adm}}$, $B_{\mathrm{free}}$, coverage equivalence \\
\texttt{Defect} & Defect dynamics, S202 counterexample explicit \\
\texttt{NegativeRuns} & $L^k$ framework, S203-C/D integrality criterion \\
\texttt{PeriodicBlocks} & S206 mod $3^{24}$ identity at native level \\
\texttt{ThreeAdic} & S206 lifted to all $\ZZ / 3^k$, $k \ge 1$ \\
\texttt{Reachability} & Computational \texttt{InImageTheta} table \\
\texttt{BinaryTensor} & 2-adic coordinates, Terras bijection ($k \le 3$) \\
\texttt{S202Cylinders} & Cylinder family $C_m$, accessibility criterion \\
\texttt{Concatenation} & S210/S211 concatenation lemma \\
\texttt{CylinderStability} & S213 $\tau$-trajectory stability \\
\texttt{PotentialBarrier} & S214 abstract potential framework \\
\texttt{InverseGraph} & Inverse-edge predicates, \texttt{outgoingEdges} \\
\texttt{AS202Lift} & $\aSat$ for $m \in \{1, 2, 3\}$ \\
\texttt{S216} & Closure-to-barrier (Thm~\ref{thm:closure-to-barrier}) \\
\texttt{S212Forward} & Forward path-word correspondence \\
\texttt{AnalyticBarrier} & Defect decomposition, conditional slope barrier \\
\texttt{Cert\_m1\_Q$Q$\_S216(\_Barrier)} & Auto-generated barriers, $Q \in \{3,5,8,10\}$ \\
\bottomrule
\end{tabular}
\end{center}

\subsection{The four explicit terminal barriers}

The following barriers are formal theorems in the Lean development.
Each is produced from a \texttt{native\_decide}-proved closure
certificate. The Python script \texttt{tools/s216\_engine.py}
generates the dist table; the exporter
\texttt{tools/cert\_to\_lean\_s216.py} emits the Lean module; the
Bool-Prop adapter \texttt{Cert\_m1\_Q\$Q\_S216\_Barrier.lean}
completes the formal chain.

\begin{center}\small
\begin{tabular}{rrrrrl}
\toprule
$m$ & $Q$ & $T$ & states & elapsed (Python, laptop) & Lean theorem \\
\midrule
1 & 3  & 1 & $35$       & $<0.01$ s & \texttt{Cert\_m1\_Q3\_S216.barrier\_m1\_Q3} \\
1 & 5  & 1 & $462$      & $0.01$ s  & \texttt{Cert\_m1\_Q5\_S216.barrier\_m1\_Q5} \\
1 & 8  & 1 & $24\,310$  & $0.14$ s  & \texttt{Cert\_m1\_Q8\_S216.barrier\_m1\_Q8} \\
1 & 10 & 1 & $352\,716$ & $2.18$ s  & \texttt{Cert\_m1\_Q10\_S216.barrier\_m1\_Q10} \\
\bottomrule
\end{tabular}
\end{center}

In all four cases, the formal theorem statement is
\(
  \forall u, \,
  q(u) \le Q \wedge n(u) \equiv \aS \pmod{3^{24}}
  \to
  \defect(u) \ge 1
\),
which by the chain of theorems gives
$\Cback(1, Q) \ge 1$ machine-verified.

\subsection{Empirically closed, Lean-pending barriers}

Three further cylinders close computationally with
\texttt{tools/s216\_engine.py} but their Lean lift is gated on a
documented four-step scaffolding refactor in
\texttt{InverseGraph.lean}: replace the \texttt{InvStart} constant
by the higher-precision $\aSat$ representative for $m \ge 2$,
adapt the path-word correspondence accordingly, regenerate the
cert exporter, and re-prove the start canonicalisation lemma. Each
step is locally checkable and is bookkeeping rather than research.

\begin{center}\small
\begin{tabular}{rrrrll}
\toprule
$m$ & $Q$ & $T$ & states & elapsed & status \\
\midrule
2 & 6 & 2 & $3\,003$    & $0.02$ s & empirical; gated on \texttt{InvStart} refactor \\
2 & 9 & 2 & $167\,960$  & $1.01$ s & empirical; gated on \texttt{InvStart} refactor \\
3 & 6 & 3 & $5\,005$    & $0.02$ s & empirical; gated on \texttt{InvStart} refactor \\
\bottomrule
\end{tabular}
\end{center}

\subsection{Auxiliary results}

In addition to the chain above, the following are also formal in
Lean and may be of independent interest:

\begin{itemize}
  \item \texttt{coverage\_translation\_equiv}: word-to-coverage
        equivalence for admissible $B$ values.
  \item \texttt{B\_adm\_to\_B\_free}: explicit strictly-decreasing
        exponent representation for every admissible $B$.
  \item \texttt{S203\_C\_D\_negative\_run\_framework}: integrality
        criterion $L^k(n) \in \ZZ \iff n \equiv -1 \pmod{3^k}$.
  \item \texttt{S206\_3adic\_lift}: lift of the $3^{24}$ identity to
        all $3^k$.
  \item \texttt{Reachability\_table}: explicit witnesses for
        \texttt{InImageTheta} at $n \in \{1, 4, 5, 11, 13, 16, 17, 40\}$,
        plus the empirical observation that $n \in \{2, 8\}$ are not
        reachable at any short depth.
\end{itemize}

% =====================================================================
\part{The open conjecture}\label{part:gap}
% =====================================================================

\section{Statement of the conjecture}\label{sec:open-conjecture}

We now state the single combinatorial conjecture to which the
unconditional asymptotic content of the programme reduces. For an
admissible word $u$, define:
\begin{itemize}
  \item $\numOnes(u)$ — the number of $\mathtt{1}$ symbols in $u$,
        equivalently the number of one-edges in $p_u$;
  \item $\numZeros(u)$ — the number of $\mathtt{0}$ symbols in $u$,
        equivalently the number of zero-edges, equivalently $q(u)$;
  \item $\numNeg(u)$ — the number of zero-edges in $p_u$ with
        $\tau(\alpha) = 1$ (the negative-weight ones);
  \item $\numPosZero(u) := \numZeros(u) - \numNeg(u)$ — the number
        of non-negative-weight zero-edges.
\end{itemize}

\begin{conjecture}[Precise form; \texttt{S202\_one\_edge\_count\_conjecture}]
\label{conj:precise}
For every admissible word $u$ with $u \in C_m$ and $q(u) \le Q$,
there exist non-negative integers
$(\numOnes', \numNeg', \numPosZero')$ such that
\begin{align*}
  \numNeg' + \numPosZero' &= \numZeros(u), \\
  \numOnes' - \numNeg' &\le \defect(u), \\
  m &\le \numOnes' - \numNeg'.
\end{align*}
\end{conjecture}

\begin{conjecture}[Coarse form]
\label{conj:coarse}
For every admissible word $u$ with $u \in C_m$ and $q(u) \le Q$,
\[
  m + q(u) \;\le\; \numOnes(u).
\]
\end{conjecture}

\begin{proposition}[Coarse implies precise;
\texttt{S202\_one\_edge\_count\_conjecture\_coarse\_implies\_precise}]
\label{prop:coarse-implies-precise}
Conjecture~\ref{conj:coarse} implies Conjecture~\ref{conj:precise}
with the witness
\(
  (\numOnes', \numNeg', \numPosZero') := (\numOnes(u), \numZeros(u), 0).
\)
\end{proposition}

\begin{proof}
With the witness, $\numNeg' + \numPosZero' = \numZeros(u) + 0$ ✓.
The bound $\numOnes' - \numNeg' = \numOnes(u) - \numZeros(u)$
combined with Lemma~\ref{lem:per-edge-bound} gives
$\numOnes(u) - \numZeros(u) \le \defect(u)$ ✓. The third inequality
$m \le \numOnes' - \numNeg' = \numOnes(u) - \numZeros(u)$ is
exactly Conjecture~\ref{conj:coarse}.
\end{proof}

\subsection{Why precise, not coarse}

Conjecture~\ref{conj:coarse} is strictly stronger and ``cleaner''
(the witness is trivial), but it is over-strong in the worst case:
along certain inverse paths the number of one-edges may genuinely
be less than $m + q$ while the precise form, which uses only the
\emph{negative} zero-edges $\numNeg$, still holds. The empirical
data of Part~\ref{part:markov} confirms this: on
$\kappa_{\mathrm{precise}}$-weight, the relevant region is
strictly and substantially smaller than on
$\kappa_{\mathrm{coarse}}$-weight.

\subsubsection*{The structural distinction between precise and coarse}

The two conjectures differ in how they treat zero-edges with
$\tau \ge 2$:
\begin{center}\small
\begin{tabular}{ll}
\toprule
zero-edge $\tau$ & precise / coarse cost \\
\midrule
$1$ (residues $\{2, 8\}$) & both: $-1$ \\
$2$ (residues $\{1, 4\}$) & precise: $\phantom{+}0$ \quad coarse: $-1$ \\
$3$ (residue $\{5\}$)     & precise: $\phantom{+}0$ \quad coarse: $-1$ \\
$4$ (residue $\{7\}$)     & precise: $\phantom{+}0$ \quad coarse: $-1$ \\
\bottomrule
\end{tabular}
\end{center}
The precise statement ``$\numOnes - \numNeg \ge m$'' only credits
$\tau = 1$ zero-edges with negative cost. The coarse statement
``$\numOnes - q \ge m$'' charges every zero-edge with $-1$ in the
budget, which is strictly over-conservative.

\subsubsection*{When coarse fails but precise holds}

Consider a hypothetical admissible word $u$ with $\numOnes = 5$,
$\numZeros = 6$, and $\numNeg = 3$ (i.e.\ three of the six zero-edges
have $\tau = 1$). Suppose $u \in C_1$, so the cylinder condition
holds with $m = 1$. Then:
\begin{itemize}
  \item Coarse: $\numOnes - q = 5 - 6 = -1$, so
        $1 + q = 7 > 5 = \numOnes$ → coarse \emph{fails}.
  \item Precise: $\numOnes - \numNeg = 5 - 3 = 2 \ge 1 = m$ → precise
        \emph{holds} (with witness
        $(\numOnes', \numNeg', \numPosZero') = (5, 3, 3)$).
\end{itemize}
Such a $u$ would refute coarse without refuting precise. None has
been found in our exploration, but the empirical scaling factor
suggests that they should exist eventually, at $Q$ large enough that
the three τ-types of non-negative zero-edges accumulate. The precise
form is the one we believe.

\subsubsection*{Lean formal statements}

For reference, the exact Lean signatures are:
\begin{lstlisting}
def S202_one_edge_count_conjecture : Prop :=
  forall (u : Word),
    AccessibleS202 u ->
    forall m Q : Nat,
      u InCylinder m -> (evalWord u).q <= Q ->
      Exists (fun n_one : Nat =>
      Exists (fun n_neg : Nat =>
      Exists (fun n_pos_zero : Nat =>
        n_neg + n_pos_zero = numZeros u
        /\ (n_one - n_neg : Int) <= defect u
        /\ (m : Int) <= n_one - n_neg)))

def S202_one_edge_count_conjecture_coarse : Prop :=
  forall (u : Word),
    AccessibleS202 u ->
    forall m Q : Nat,
      u InCylinder m -> (evalWord u).q <= Q ->
      m + (evalWord u).q <= numOnes u
\end{lstlisting}

\subsection{The conditional slope barrier}

\begin{theorem}[Conditional slope barrier;
\texttt{S202\_slope\_barrier\_conditional}]
\label{thm:slope-conditional}
Conjecture~\ref{conj:precise} implies that for every $m \ge 1$ and
$Q \ge 0$,
\[
  \Cback(m, Q) \;\ge\; m.
\]
Equivalently,
$\lambda_{\mathrm{asy}} \ge 2 \Rightarrow \text{Collatz forward}$ along
any admissible cover.
\end{theorem}

\begin{proof}[Proof (formalised, one line)]
Let $u$ be any admissible word with $u \in C_m$ and $q(u) \le Q$.
By Conjecture~\ref{conj:precise}, there exists
$(\numOnes', \numNeg', \numPosZero')$ with
$m \le \numOnes' - \numNeg' \le \defect(u)$.
\end{proof}

The Lean formalisation is genuinely a one-line application of two
existing lemmas; the entire remaining gap is the conjecture itself.

\subsection{Empirical support}

For all seven $(m, Q)$ frontiers tabulated in
\S\ref{sec:formal-status}, the corresponding closure certificate
under the defect weight $\omega$ proves
$\defect(u) \ge m$ on the cylinder. By
Lemma~\ref{lem:per-edge-bound}, this gives a fortiori
$\numOnes(u) - \numZeros(u) \le \defect(u)$, but the conjecture
strictly asks for an additional refinement (the existence of a
witness with $m \le \numOnes' - \numNeg'$). The defect barrier
gives one of the two inequalities of the conjecture; the
combinatorial one $m \le \numOnes' - \numNeg'$ is empirically
verified by the $\kappa$-closure of Part~\ref{part:markov} but not
yet structurally explained.

No counterexample has been found within the explored region.

\subsection{Why the conjecture is structurally hard}

The straightforward attempt to bound $\numOnes - \numNeg$ from below
along inverse paths uses the local 3-adic obstruction
$\rho_-(c) := \nuthree(c + 1)$ (the proximity to $-1$ in the 3-adic
metric). The deepest chains of $\tau = 1$ zero-edges drop
$\rho_-$ by exactly one step per edge, which suffices to close
$\rho_-$-based potentials only when correlated with the global
$j$-budget. \emph{Decoupling these two scales} is the present
analytical bottleneck. Part~\ref{part:markov} addresses this via the
Markov/flow dual programme.

% =====================================================================
\part{The Markov/flow dual programme}\label{part:markov}
% =====================================================================

\section{Tao 2026 inspiration}\label{sec:tao-inspiration}

\subsection{The von Mangoldt chain methodology}

In May 2026, Tao and collaborators announced (on Tao's blog and at
\href{https://arxiv.org/abs/2605.00301}{\texttt{arXiv:2605.00301}})
a new method for the Erdős 1196 problem on \emph{primitive sets}.
The methodological core is: instead of bounding the obstruction
combinatorially by enumerating bad antichains, build a measure or
chain whose invariant captures the relevant weight directly. The
key innovation is to use the
\emph{von Mangoldt-weighted Markov chain} on the divisibility poset,
which aligns with the natural weight $1/(n \log n)$ relevant to
primitive sets and avoids the $e^\gamma$ losses of the Mertens
chain.

\subsection{The abstract pattern}

The abstract pattern of which the von Mangoldt construction is one
instance can be stated as:
\[
  \boxed{
    \begin{array}{c}
       \text{combinatorial obstruction} \\
       \Longrightarrow \\
       \text{measure / chain over chains aligned with target weight} \\
       \Longrightarrow \\
       \text{dual potential certificate}
    \end{array}
  }
\]

We will translate this pattern to our cylinder accessibility setting
in the next section.

\subsubsection*{Why the pattern translates}

Three structural features of the von Mangoldt setting carry over
to ours:

\paragraph{Antichain vs. shortest path.}
Tao's primitive sets are antichains in the divisibility poset.
The dual is a chain measure that hits every element with prescribed
mass. In our setting, the obstruction is a path with low
$\kappa$-weight in the inverse cylinder graph $\InvG_{m, Q}$,
which is a directed acyclic structure (zero-edges increment $j$,
one-edges preserve $j$, so cycles can only be one-edge cycles
within a single $j$-level; the BF closure handles these via the
relevance cut). The dual is a node potential $H$ that lower-bounds
all path costs by telescoping.

\paragraph{Weight alignment.}
Tao replaces the obvious Mertens weight ($1/n$) with the
$\Lambda$-weighted variant ($\log p$ for $n = p^k$, else $0$). The
replacement aligns the chain's invariant measure with the
$1/(n \log n)$ target. In our setting, the obvious weight is the
defect-weight $\omega \in \{-1, 0, 1, 2\}$ (the four values of
$\tau - 2$ for zero-edges plus $\tau \in \{1, 2, 3, 4\}$ for
one-edges, with the specific multiplicities determined by the
$\tau$-table). The aligned weight is $\kappa_{\mathrm{precise}}$,
which captures exactly the combinatorial quantity
$\numOnes - \numNeg$ that the conjecture concerns.

\paragraph{Dual certificate.}
Tao's certificate is a measure-theoretic upper bound that
``hits'' every element of any antichain. Our certificate is a node
potential satisfying the per-edge inequality
$H(v) \le \kappa(e) + H(v')$ with $H(s) \ge m, H(g) \le 0$. By
telescoping, $\sum_e \kappa(e) \ge H(s) - H(g) \ge m$, so every
path has $\kappa$-cost $\ge m$.

The structural parallel is exact at the level of LP duality. The
content of the translation is in finding the right $\kappa$ (we
have it, namely $\kappa_{\mathrm{precise}}$) and the right $H$
(this is what \S\ref{sec:s233} attempts).

\subsubsection*{What Tao's method does NOT give us directly}

Two important caveats:

\paragraph{The Markov chain structure is implicit, not explicit.}
In Tao's setting the chain is a Markov chain on $\NN$ with explicit
transition kernel and stationary measure. In our setting, the
analogue would be a Markov chain on $V_{m, Q}$ with kernel
$P(v \to v')$ such that the harmonic / superharmonic equation
$\sum_{v'} P(v \to v') (H(v') + \kappa(e)) \ge H(v)$ recovers the
LP. This is a useful conceptual frame for searching the right
$\Phi$ (see \S\ref{sec:s234}), but it does not give the answer.

\paragraph{The von Mangoldt weights are
problem-specific.}
What ``von Mangoldt'' means in our setting depends on the local
3-adic structure. Candidates include
$\tau(\alpha)$-weighted chains,
$\rho_-$-weighted chains, and
$\delta_+$-weighted chains; none of these has been investigated
exhaustively. Question B of \S\ref{sec:alphaevolve} is precisely
the request to investigate this systematically.

\section{The $\kappa$-reformulation}\label{sec:s230-kappa}

\subsection{Combinatorial weight on the inverse cylinder graph}

We already defined (\S\ref{sec:inverse-graph}) the combinatorial
weight
\[
  \kappa(e) \;:=\;
  \begin{cases}
    +1 & e \text{ is a one-edge}, \\
    -1 & e \text{ is a zero-edge with } \tau = 1, \\
    \phantom{+}0 & e \text{ is a zero-edge with } \tau \ge 2.
  \end{cases}
\]
This is the \emph{precise} variant. The \emph{coarse} variant gives
every zero-edge weight $-1$.

\begin{lemma}[$\kappa$-sum identity]
\label{lem:kappa-sum}
Along the path $p_u$ corresponding to an admissible word $u$:
\begin{align*}
   \sum_{e \in p_u} \kappa_{\mathrm{precise}}(e) &= \numOnes(u) - \numNeg(u), \\
   \sum_{e \in p_u} \kappa_{\mathrm{coarse}}(e) &= \numOnes(u) - \numZeros(u).
\end{align*}
\end{lemma}

The proof is by direct count of one-edge and zero-edge contributions.

\subsection{The dual reformulation}

Conjecture~\ref{conj:precise} can therefore be restated as the
\emph{shortest-path lower bound}
\[
  \operatorname{dist}_{\kappa_{\mathrm{precise}}}(s, \mathrm{goal})
   \;\ge\; m,
\]
and Conjecture~\ref{conj:coarse} as
\(
  \operatorname{dist}_{\kappa_{\mathrm{coarse}}}(s, \mathrm{goal}) \ge m.
\)
The Bellman–Ford LP dual asks for a potential
$H : V \to \ZZ$ satisfying
\begin{align*}
  H(v) &\le \kappa(e) + H(v')   \quad \text{for every edge } v \to v', \\
  H(s) &\ge m, \\
  H(g) &\le 0 \quad \text{for every reachable goal } g.
\end{align*}
This is the \emph{S230 reformulation} of the programme.

\subsection{Empirical $\kappa$-closure}

We re-ran the BF closure search with the cost function $\kappa$ in
both variants. No counterexample appeared on any cylinder; six of
the seven frontier cases (the same defect-closed ones) also
$\kappa_{\mathrm{precise}}$-closed, with the precise relevant region
being $6$-$26\times$ smaller than coarse:

\begin{center}\small
\begin{tabular}{rrrrrl}
\toprule
$m$ & $Q$ & $\kappa_{\mathrm{precise}}$ states & $\kappa_{\mathrm{coarse}}$ states & time precise & status \\
\midrule
1 & 3  & $137$        & $703$            & $<0.01$ s & both close, $\ge m = 1$ \\
1 & 5  & $4\,425$     & $80\,665$        & $0.03$ s  & both close, $\ge m = 1$ \\
1 & 8  & $893\,376$   & $> 5 \times 10^6$ & $11.7$ s  & precise closes, coarse exhausts budget \\
2 & 3  & $4\,000$ approx & $1\,187$      & $< 0.01$ s & both close, $\ge m = 2$ \\
2 & 6  & $61\,645$    & $1\,612\,549$    & $0.85$ s  & both close, $\ge m = 2$ \\
2 & 9  & $13\,414\,169$ & $> 2 \times 10^7$ & $96$ s  & precise closes at $20$M, coarse exhausts \\
3 & 6  & $134\,913$   & $2\,765\,825$    & $1.25$ s  & both close, $\ge m = 3$ \\
\bottomrule
\end{tabular}
\end{center}

\noindent The $\kappa_{\mathrm{precise}}$ variant being uniformly
more efficient is the empirical analogue of the
\emph{von Mangoldt-vs.-Mertens} dichotomy: the natural chain aligns
with the natural weight. The precise statement
($\numOnes - \numNeg \ge m$) is therefore the correct one — the
coarse statement ($\numOnes - q \ge m$) requires a strictly larger
search region to certify and presumably has a strictly weaker
sense of ``natural''.

\section{S233: local potential candidates}\label{sec:s233}

Following the Tao methodology, we sought a closed-form
$H : V \to \ZZ$ of the parametric form
\begin{equation}
\label{eq:ansatz}
  H(v)
  \;=\;
  a \cdot \delta_+(v) + b \cdot \rho_-(v)
   + c \cdot j(v) + d \cdot (j(v) - Q)
   + \psi(c(v) \bmod 3^h),
\end{equation}
where:
\begin{itemize}
  \item $\delta_+(v) := \nuthree(c(v) - 1)$, capped at $R + j(v)$;
  \item $\rho_-(v) := \nuthree(c(v) + 1)$, capped at $R + j(v)$;
  \item $a, b, c, d$ are integer coefficients to be determined;
  \item $h$ is a local-modulus exponent;
  \item $\psi : \ZZ/3^h\ZZ \to \ZZ$ is a local correction, also
        to be determined.
\end{itemize}

The unknowns are $(a, b, c, d, h)$ and the $3^h$ values of $\psi$,
and the constraint system is a linear program over the differences
$\psi(z) - \psi(z')$ across edges of (a subset of) the cylinder
graph.

\subsection{The edge inequality under the ansatz}

For any edge $v \to v'$ with cost $\kappa(e)$, the condition
$H(v) \le \kappa(e) + H(v')$ rewrites as
\[
  \psi(z) - \psi(z')
   \;\le\; \kappa(e) - L(v) + L(v'),
\]
where
$L(v) := a \delta_+(v) + b \rho_-(v) + c j(v) + d (j(v) - Q)$ is the
linear part and $z, z'$ are the modular classes
$c(v) \bmod 3^h, c(v') \bmod 3^h$. This is a
\emph{difference constraint} system in the variables $\psi(z)$, and
feasibility is decidable in $O(|V| \cdot |E|)$ time by
Bellman--Ford on the difference graph (negative-cycle detection).

\subsubsection*{Motivation for the parametric form}

The choice of features
$(\delta_+, \rho_-, j, (j - Q), c \bmod 3^h)$ is justified by:

\paragraph{$\delta_+(v) = \nuthree(c - 1)$.}
The 3-adic distance to $1$. The cylinder anchor $\aSat$ has
$\delta_+(\aSat) = 22$ (by Lemma~\ref{lem:nu3-aSm-minus-1}); the
root $g_j = (j, 1)$ has $\delta_+(g_j) = \infty$ (capped at $R + j$).
A potential that scales with $\delta_+$ punishes proximity to the
cylinder anchor and rewards proximity to the root, a natural
signature.

\paragraph{$\rho_-(v) = \nuthree(c + 1)$.}
The 3-adic distance to $-1$. The $\tau = 1$ zero-edges (residues
$\{2, 8\}$) are precisely those that ``head towards $-1$'' in the
3-adic topology — these are the dangerous edges that drop the
$\kappa$-cost by $1$. A potential with negative $\rho_-$ coefficient
punishes vertices near $-1$, exactly the population that
$\tau = 1$ zero-edges produce.

\paragraph{$j(v)$ and $(j(v) - Q)$.}
Two linearly-independent counters for the zero-edge budget. Their
linear combination captures the residual credit and the absolute
depth, independently.

\paragraph{$\psi(c \bmod 3^h)$.}
A free local correction term, absorbing finite-modulus structure
not captured by the other features. The choice of $h$ is the key
parameter; we will see that $h \ge 22$ is forced by the structural
lemma~\ref{lem:nu3-aSm-minus-1}.

\subsection{The collaborator's announcement}

In May 2026, an external collaborator (in the form of an
\texttt{S233 Local ψ Miner} package) announced that the candidate
\[
  (a, b, c, d) \;=\; (2, -1, 0, 1), \quad h = 3,
\]
made the LP feasible on $(m, Q) = (1, 3)$, certifying the precise
conjecture on this cylinder. A larger version with
$(a, b, c, d) = (1, -1, 0, 4), h = 2$ allegedly certified $(2, 3)$
with margin $12$.

We attempted independent verification.

\subsection{Independent reproduction (initial, naive)}

Our first LP validator agreed with the announcement and, with
several integer-coefficient sweeps, found candidate combinations
that appeared to certify the conjecture for the entire frontier
$\{(1,3), (1,5), (1,8), (2,3), (2,6), (3,6)\}$. The most striking
hit was
\[
  H(v) \;=\; 2 j(v) - \rho_-(v) - Q + \psi(c(v) \bmod 3^h),
\]
which reported feasibility at $h = 3$ for $(1, 3)$ with margin $1$
and at $h = 4$ for $(3, 6)$ with margin $3$. The empirical
fit was striking; we now know it was a false positive due to a
shared validator bug.

\section{The shared validator bug}\label{sec:bug}

\subsection{Diagnosis}

The naive LP imposes one constraint per observed edge
$v \to v'$ in the BF closure graph:
\[
  \psi(z) - \psi(z')
   \;\le\; \kappa(e) - L(v) + L(v').
\]
The barrier requires additionally
\[
  H(g_j) \;\le\; 0 \quad \text{for every } j \in [0, Q],
\]
where $g_j = (j, 1)$ is the canonical goal vertex at level $j$.
Crucially, the BF closure terminates without visiting goal
vertices when the barrier $\ge T$ is achieved (that is precisely
what makes it a barrier closure). Hence:
\begin{quote}
\emph{The observed-graph edge enumeration contains no
$g_j$ — the LP, as written, leaves
$\psi(1 \bmod 3^h)$ unconstrained.}
\end{quote}
A ``feasible'' LP solution may thus assign $\psi(1 \bmod 3^h)$ any
value, in particular one that violates $H(g_j) \le 0$ for some
$j$. The certificate is silently invalid.

\subsection{The unavoidable distinction}

Lemma~\ref{lem:nu3-aSm-minus-1} gives $\nuthree(\aSat - 1) = 22$.
In particular, $\aSat \equiv 1 \pmod{3^{h}}$ for every $h \le 22$.
For such $h$, the cylinder anchor $\aSat$ and the root $1$ reduce
to the same class in $\ZZ/3^h\ZZ$, so $\psi$ cannot distinguish
them:
\[
  \psi(\aSat \bmod 3^h) \;=\; \psi(1 \bmod 3^h)
  \quad \text{for } h \le 22.
\]
The naive ``$H(s) - H(g) \ge m$'' margin therefore must come
entirely from the linear part $L$, and the goal-anchor constraint
is binding.

\subsection{The corrected validator}

We add two explicit difference-system constraints via a virtual
bound vertex $S^{*}$ anchored at $0$:
\begin{align*}
  \psi(1 \bmod 3^h) - S^{*}
    &\;\le\; - L(\text{goal}_{Q}), \\
  S^{*} - \psi(\aSat \bmod 3^h)
    &\;\le\;   L(\text{start}) - m.
\end{align*}
The first says $H(g_Q) \le 0$ at the worst goal; the second
$H(\text{start}) \ge m$. The implementation is in
\texttt{tools/s233\_psi\_validator.py}, function
\texttt{psi\_feasibility}.

\subsection{After the fix}

Running the corrected validator on the collaborator's announced
$(a, b, c, d, h) = (2, -1, 0, 1, 3)$ for $(m, Q) = (1, 3)$ returns
\textsc{infeasible} (negative cycle in the difference graph).
The candidate does not certify the conjecture.

Extending $h$ for the same coefficients:

\begin{center}\small
\begin{tabular}{rrll}
\toprule
$(m, Q)$ & $h$ & coefs $(a, b, c, d)$ & status \\
\midrule
$(1, 3)$ & $3$ & $(2, -1, 0, 1)$ & \textsc{infeasible} \\
$(1, 3)$ & $22$ & $(2, -1, 0, 1)$ & \textsc{infeasible} \\
$(1, 3)$ & $25$ & $(2, -1, 0, 1)$ & \textsc{infeasible} \\
$(1, 3)$ & $28$ & $(2, -1, 0, 1)$ & \textsc{feasible}, margin $\ge 1$ \\
\bottomrule
\end{tabular}
\end{center}

The candidate is salvageable but only at $h \ge 28$. A different
linear ansatz, found by sweep (next section), closes at $h = 22$.

\subsubsection*{Why the bug went undetected initially}

Several factors made the bug hard to spot:

\paragraph{The naive LP \emph{is} feasible.}
Without the goal-anchor constraint, the LP is genuinely feasible
(many solutions exist), and Bellman--Ford reports a witness. The
output is well-formed JSON; the verifier returns no error.

\paragraph{Empirical agreement with the conjecture.}
The closed BF certificates (\S\ref{sec:formal-status}) all
\emph{do} establish the precise conjecture on the observed
cylinder (by the defect barrier $\defect \ge m$). So the LP's
``feasible'' verdict matches the underlying truth, even though
the certificate itself is not a valid \emph{proof} of that truth.
A specification-level proof would have caught this; an empirical
spot-check would not.

\paragraph{The structural explanation is not local.}
The reason $\psi$ cannot certify at $h \le 22$ is the global
3-adic distance $\nuthree(\aSat - 1) = 22$, not any local
property of the graph. A local code reviewer would not see the
obstruction; it requires the structural lemma to surface.

\paragraph{Lesson.}
In LP-based certificate frameworks: \emph{always anchor the boundary
conditions explicitly}, even when the implicit semantics seem to
match. Adding two virtual constraints (start, goal) is a constant
overhead and avoids the entire class of false positives.

\section{The corrected sweep and validated certificates}\label{sec:s238}

\subsection{Methodology}

The corrected LP makes a systematic sweep across integer-coefficient
combinations $(a, b, c, d)$ and modular precision $h$ feasible.
Implementation: \texttt{tools/s238\_sweep\_parallel.py}, which:
\begin{itemize}
  \item precomputes the BF closure for each $(m, Q)$ \emph{once}
        (not per coefficient combination, as in the initial naive
        implementation);
  \item dispatches the grid
        $(a, b, c, d) \in [-3, 1]^2 \times [-1, 2]^2 = 400$ cells
        per $h$ to a \texttt{multiprocessing.Pool} of CPU workers;
  \item iterates $h$ over $\{22, 23, 25, 28, 30\}$;
  \item early-stops on the first feasible cell per $(m, Q)$;
  \item exports the LP-derived $\psi$ in \emph{sparse} JSON
        (only the classes that actually participate in some
        constraint — typically $\sim |E|$, not $3^h$, which is
        astronomical for $h = 22$).
\end{itemize}

\subsection{Validated certificates}

\begin{theorem}[Observed-graph certificates;
\texttt{tools/s238\_psi\_verifier\_precise.py}]
\label{thm:obs-certs}
The following $\psi$-certificates pass the corrected verifier on
the observed BF closure graph under
$\kappa_{\mathrm{precise}}$:

\begin{center}\footnotesize
\begin{tabular}{rrrrrrrr}
\toprule
$m$ & $Q$ & $h$ & coefs $(a, b, c, d)$ & edges & $H(s)$ & $\max_j H(g_j)$ & margin \\
\midrule
1 & 3 & 22 & $(-1, -3, -1, 2)$ & $136$    & $1$ & $-1$ & $2$ \\
1 & 5 & 22 & found by sweep    & $4\,424$ & $1$ & $-1$ & $2$ \\
2 & 3 & 22 & found by sweep    & $284$    & $2$ & $-22$ & $24$ \\
\bottomrule
\end{tabular}
\end{center}

In each case the LP-derived $\psi$ is sparse: of the $3^{22}$
possible residues, only $\sim$$|E|$ classes are needed; the rest
may be set to $0$ without violating any inequality.
\end{theorem}

\subsection{Verification trace}

The certificate file
\texttt{tools/s238\_tables\_parallel/psi\_m1\_Q3\_h22\_precise\_parallel.json}
is a sparse JSON record of $\sim$80 entries. Running

\begin{lstlisting}
python tools/s238_psi_verifier_precise.py \
    tools/s238_tables_parallel/psi_m1_Q3_h22_precise_parallel.json
\end{lstlisting}

outputs (excerpt):

\begin{lstlisting}
{
  "overall_ok": true,
  "m": 1, "Q": 3, "h": 22, "mode": "precise",
  "linear_coefficients": {
    "a_dp": -1, "a_rm": -3, "a_j": -1, "a_jQ": 2
  },
  "n_observed_edges_checked": 136,
  "edge_check_pass": true,
  "worst_edge_violation": null,
  "H_start": 1,
  "H_goal_per_j": [
    [0, -4], [1, -3], [2, -2], [3, -1]
  ],
  "max_H_goal": -1,
  "endpoint_margin": 2,
  "endpoint_required": 1,
  "endpoint_pass": true
}
\end{lstlisting}

The verifier checks: (a) every observed edge satisfies
$H(v) \le \kappa(e) + H(v')$; (b) $H(s) \ge m$; (c)
$H(g_j) \le 0$ for every $j \in [0, Q]$ (not only those visited).
All three pass. The same procedure verifies $(1, 5)$ and $(2, 3)$.

\subsection{Status of the remaining frontier cases}

The cases $(m, Q) = (2, 6)$ and $(3, 6)$ have closed BF
certificates and the parallel sweep was launched on them, but a
worker was OOM-killed before completing the LP (the $h = 22$
constraint system at these sizes consumes $\sim 10^{10}$ classes if
naively expanded). A larger-memory machine or a smarter
$h$-scheduling strategy is needed; the search space is well-defined
and does not require any new mathematical input.

\subsection{What is and is not proved}

\begin{itemize}
  \item \textbf{Proved}: $\psi$-certificates over the
        \emph{observed} BF closure graph for three cylinders.
        Edge inequalities and endpoint margin all hold.
  \item \textbf{Not proved}: the same certificate over the
        \emph{abstract over-cover} (S234) — the set of
        \emph{all} possible $(v, v')$ pairs respecting the
        cylinder edge predicate, including those not visited by
        the specific BF closure trajectory. This is the next
        analytic step; see \S\ref{sec:s234}.
  \item \textbf{Open}: the conjecture itself
        (\S\ref{sec:open-conjecture}) remains unresolved.
        The observed-graph evidence is empirical support, not
        proof.
\end{itemize}

% =====================================================================
\part{Future work and cooperation with AlphaEvolve}\label{part:future}
% =====================================================================

\section{The abstract over-cover (S234)}\label{sec:s234}

\subsection{The construction}

The observed-graph certificate of Theorem~\ref{thm:obs-certs}
verifies the edge inequality $\psi(z) - \psi(z') \le \cdots$ only
for those edges that the BF closure actually traversed. For a
complete proof of the conjecture on a given $(m, Q)$, we need
the same inequality over the
\emph{abstract over-cover} graph $\InvG^{*}_{m, Q}$ defined as
follows.

\begin{definition}[Abstract over-cover]
\label{def:over-cover}
The vertices of $\InvG^{*}_{m, Q}$ are pairs
\(
   (j, z) \in \{0, 1, \dots, Q\} \times \ZZ/3^h\ZZ.
\)
For each $(j, z)$, the outgoing edges are:
\begin{itemize}
  \item For each $\alpha \in \InX$: a one-edge
        $(j, z) \to (j, z')$ where
        $z' = (2^{\tau(\alpha)})^{-1} \cdot z \bmod 3^h$,
        provided $z' \bmod 9 = \alpha$.
        Edge weight: $\kappa = +1$.
  \item For each $\alpha \in \InX$ with $j < Q$: a zero-edge
        $(j, z) \to (j + 1, z')$ where
        $z' = (2^{\tau(\alpha)})^{-1} \cdot (3 z + 1) \bmod 3^h$,
        provided $z' \bmod 9 = \alpha$.
        Edge weight: $\kappa = -1$ if $\tau(\alpha) = 1$ else $0$
        (precise) or $-1$ (coarse).
\end{itemize}
\end{definition}

Note that this graph is \emph{independent} of any specific
trajectory; it depends only on $(m, Q, h)$ and the modular
arithmetic. Its number of vertices is at most
$(Q + 1) \cdot 3^h$, with at most $12$ outgoing edges per vertex.

\subsection{The validation question}

\textbf{S234.} Does there exist a $\psi : \ZZ/3^h\ZZ \to \ZZ$
(possibly together with linear coefficients $(a, b, c, d)$)
satisfying $\psi(z) - \psi(z') \le \cdots$ across \emph{every}
edge of $\InvG^{*}_{m, Q}$?

\subsection{Outcomes}

Either outcome is informative:

\begin{itemize}
  \item \textbf{Feasible at moderate $h$}: yields an
        unconditional proof of Conjecture~\ref{conj:precise} on
        the cylinder $(m, Q)$, and hence (by varying $m, Q$)
        opens the possibility of an unconditional slope barrier
        $\lambda_{\mathrm{asy}} < 2$.
  \item \textbf{Infeasible for all $h$ up to some bound}: refines
        the conjecture statement to include the local correction
        forced by the infeasibility (a la
        \emph{von Mangoldt vs Mertens}).
  \item \textbf{Refutation by counterexample word}: would refute
        Conjecture~\ref{conj:precise} and reformulate the entire
        programme.
\end{itemize}

\subsection{Why our local sweep does not answer this}

Our sweep \texttt{tools/s238\_sweep\_parallel.py} validates the
LP only on the \emph{observed} BF graph, which is a strict subset
of $\InvG^{*}_{m, Q}$. The validator is correct on the subset
but cannot extrapolate to unobserved edges. Constructing
$\InvG^{*}_{m, Q}$ for $h \ge 22$ requires at most
$O((Q + 1) \cdot 3^h \cdot |\InX|)$ edges, which is on the order
of $10^{10}$ for $Q = 3, h = 22$ — manageable with a
specialised implementation but well outside our current sweep
framework.

\section{Two specific questions for AlphaEvolve}\label{sec:alphaevolve}

The two questions below are formulated in the operational style
known to be productive for evolutionary LLM-guided coding agents
such as Google DeepMind's AlphaEvolve (announced May 2025;
see the public DeepMind blog post on AlphaEvolve and its
applications to algorithmic discovery).

\paragraph{Question A — combinatorial / exhaustive verification.}
Apply the agent to extend the BF closure search to
\[
  (m, Q) \;\in\;
  \{ (1, 12),\; (1, 15),\; (2, 12),\; (4, 6) \}.
\]
For each:
\begin{itemize}
  \item produce a $\kappa_{\mathrm{precise}}$ barrier certificate
        up to the threshold $T = m$, which adds an empirical row
        to the seven existing frontiers in
        \S\ref{sec:formal-status}; or
  \item exhibit a single admissible word $u$ reaching the
        cylinder with
        $\sum_{e \in p_u} \kappa_{\mathrm{precise}}(e) < m$,
        which would \emph{refute}
        Conjecture~\ref{conj:precise} and force a re-formulation.
\end{itemize}
The success criterion is mechanical and reportable in either
direction. The Python infrastructure in
\texttt{tools/s230\_dual\_search.py} is the natural starting point
for the agent.

\paragraph{Question B — analytic / potential search.}
Apply the agent to the search for a $\tau$-dependent (or
$\tau$-stratified) potential function
$\Phi : V \to \ZZ$ on the \emph{abstract} inverse cylinder graph
$\InvG^{*}_{m, Q}$ satisfying
\[
   \Phi(v) \;\le\; \kappa(e) + \Phi(v')
\]
for every edge of the over-cover, with
$\Phi(s) - \max_g \Phi(g) \ge m$.

The candidate parametric form from \S\ref{sec:s233}
(linear-plus-$\psi$) is one starting point; an agent of the
AlphaEvolve type is uniquely well placed to:
\begin{itemize}
  \item enumerate alternative parametric forms with additional
        features ($j^2$, $j \cdot \delta_+$, $\psi$ indexed by
        $(c \bmod 3^h, j \bmod k)$, etc.);
  \item identify the form that unconditionally closes the
        over-cover at moderate $h$;
  \item produce the corresponding $\psi$-table and verify the LP
        end-to-end.
\end{itemize}
The existence of such a $\Phi$ would close
Conjecture~\ref{conj:precise} unconditionally, by an argument
already drafted in \texttt{AnalyticBarrier.lean}.

\paragraph{Why these questions, not ``solve Collatz''.}
Both questions are well-typed, finitely specifiable, and admit
explicit success/failure criteria that are mechanically
checkable by the verifier in
\texttt{tools/s238\_psi\_verifier\_precise.py}. We deliberately
do not ask AlphaEvolve to ``attempt the Collatz conjecture'';
such open prompts are known not to be a good match for
evolutionary coding agents.

\subsection{Concrete cooperation protocol}

If the AlphaEvolve team or a comparable evolutionary coding agent
were to engage with either question, the following protocol would
make the interaction productive:

\paragraph{Step 1 — Reproduce.}
The agent should first reproduce the existing results on the
seven frontier $(m, Q)$ pairs. Specifically:
\begin{itemize}
   \item Run \texttt{tools/s216\_engine.py} and verify the defect
         barriers $\Cback(m, Q) \ge T = m$ for the seven cases.
   \item Run \texttt{tools/s230\_dual\_search.py} and confirm
         the $\kappa$-precise closure.
   \item Run \texttt{tools/s238\_sweep\_parallel.py} and reproduce
         the three observed-graph $\psi$-certificates of
         Theorem~\ref{thm:obs-certs}.
   \item Run \texttt{tools/s238\_psi\_verifier\_precise.py} on the
         JSON output to confirm the certificates pass independent
         verification.
\end{itemize}
This establishes a shared baseline of trusted ground truth.

\paragraph{Step 2 — Extend the empirical frontier (Question A).}
The agent extends the BF closure search to $(m, Q) \in
\{(1, 12), (1, 15), (2, 12), (4, 6)\}$. For each, the agent reports
either a barrier certificate JSON or a counterexample word as
admissible word with $\sum \kappa < m$. The verifier accepts both
outputs mechanically.

\paragraph{Step 3 — Search the abstract potential (Question B).}
The agent explores the parametric form
\[
   \Phi(v) = (\text{features}_1, \ldots, \text{features}_k)
            \cdot (\text{coefs}) + \psi(\text{index}(v))
\]
over a structured set of feature configurations: $(j^a \delta_+^b
\rho_-^c)$ monomials up to some total degree, $\psi$ indexed by
$(c \bmod 3^h, j \bmod k)$ for small $h, k$, plus a $\tau$-indexed
$\psi$ component. For each configuration, the agent runs the
corrected LP over the abstract over-cover $\InvG^{*}_{m, Q}$ (not
just the observed graph) and reports the first feasible
configuration.

\paragraph{Step 4 — Lift to Lean.}
Successful $\psi$-certificates over the over-cover are exported as
JSON and the Lean adapter
\texttt{tools/s238\_export\_psi.py} (or its extension) emits a
Lean module that re-checks the certificate by \texttt{native\_decide}.
This produces a formal Lean theorem corresponding to
$\Cback(m, Q) \ge m$ for each successfully closed $(m, Q)$.

\paragraph{Step 5 — Report.}
The agent's report should include:
\begin{itemize}
   \item the verified certificates (JSON);
   \item the Lean modules (auto-generated, machine-checkable);
   \item the parametric forms searched and their feasibility status;
   \item any counterexample words found (if Question A is
         answered negatively for some $(m, Q)$).
\end{itemize}

\subsection{Compute budget estimate}

For Question A on $(m, Q) = (1, 15)$: roughly $5 \times 10^7$
states, $\sim 5$ minutes single-thread, $\sim 30$ seconds with a
$10$-core pool. Memory: $\sim 5$ GB. Tractable on any modern
laptop.

For Question A on $(1, 18)$ or $(4, 6)$: scales to $10^9$ states,
$\sim 1$ hour single-thread, $\sim 5$ minutes with $12$ cores.
Memory: $\sim 50$ GB. Requires a cluster node or AlphaEvolve-class
compute.

For Question B at $h = 22$ with $(m, Q) = (1, 3)$: the over-cover
has $\le 4 \times 3^{22} \approx 10^{10}$ vertices and $\le 12$
edges per vertex, $\sim 10^{11}$ edges. The LP itself is sparse
($\le 3 \cdot 10^{10}$ non-trivial constraints). On a single-node
optimised implementation: $\sim 1$ hour with $\sim 100$ GB
memory; on AlphaEvolve-class compute: $\sim 1$ minute. The key
question is whether the LP is feasible at all, which is a
parametric-form question, not a budget question.

\subsection{Why we believe these questions are worth the
investment}

Both questions have explicit success criteria that, if attained,
move the programme forward in significant ways:

\paragraph{Question A success.}
A barrier certificate at $(1, 15)$ would empirically demonstrate
$\Cback(1, 15) \ge 1$, doubling the empirical frontier and
ruling out a substantial new region of the search space. A
counterexample, conversely, would refute
Conjecture~\ref{conj:precise} and force re-examination of the
admissibility section.

\paragraph{Question B success.}
A $\Phi$ that closes the abstract over-cover at moderate $h$ would
constitute an unconditional proof of
Conjecture~\ref{conj:precise} on the corresponding cylinders.
Combined with the Lean development, this would yield a
formal proof of the slope barrier $\lambda_{\mathrm{asy}} < 2$ on
specific cylinders, with the only remaining question being
whether the parameters extend uniformly in $m$.

Either question, if answered, materially advances the programme
toward (but does not complete) the resolution of the Collatz
conjecture.

\section{Relation to the forward Tao 2022 theorem}
\label{sec:relation-tao}

Tao's \emph{Forum Math.\ Pi} bounded-value theorem (2022)
asserts that for any $f(n) \to \infty$, the set of $n$ whose
Collatz orbit attains a value below $f(n)$ has logarithmic
density~$1$. The methods are stochastic on the parity vector,
building on Kontorovich--Sinai and Sárközy-type sum estimates.

The present programme is the \emph{structural inverse
complement}: rather than tracking forward orbits of typical $n$,
it asks whether a \emph{specific} 3-adic fixed point is
reachable from the integer~$1$ along admissible inverse paths
with bounded cumulative cost. The two are logically independent:
\begin{itemize}
  \item Tao~2022 says nothing about the depth at which a typical
        orbit visits low values, nor about quantitative slope.
  \item The present programme says nothing about typical orbits.
\end{itemize}

A combination of the two would yield a forward,
almost-everywhere, \emph{quantitative} orbit-depth bound
conditional on the admissible cover. Whether this combination
can be made unconditional is, in the author's view, the deepest
structural question separating the present work from a full
proof of the Collatz conjecture. We do not claim a path to that
combination here.

\subsection{Stochastic vs.\ structural complementarity}

The complementarity between Tao 2022 and the present programme can
be made precise as follows. Both speak about the same dynamical
system (Collatz / accelerated Collatz $T$), but at different
levels of granularity:

\begin{center}\small
\begin{tabular}{p{4cm}p{4.5cm}p{4.5cm}}
\toprule
Aspect & Tao 2022 & Present programme \\
\midrule
Direction & forward orbits ($n \to T(n)$) & inverse orbits ($n \leftarrow T(n)$) \\
Granularity & almost every $n \in \NN$ (log-density 1) & specific cylinder $C_m$ \\
Conclusion & orbit reaches almost bounded value & bound on cumulative defect cost \\
Method & stochastic on parity vector & combinatorial / shortest path \\
Quantitative? & qualitative (existence of $f \to \infty$) & quantitative (lower bound $\Cback(m, Q) \ge T$) \\
Unconditional? & yes & conditional on Conjecture~\ref{conj:precise} \\
Mechanised? & no (informal proof) & yes (Lean 4) \\
\bottomrule
\end{tabular}
\end{center}

\subsection{What an unconditional combination would say}

Suppose, conditionally on
Conjecture~\ref{conj:precise}, we have proven the slope barrier
$\lambda_{\mathrm{asy}} < 2$ on every admissible cover. Then,
combining with Tao 2022:
\begin{itemize}
   \item For almost every $n$, the forward orbit reaches some
         value $\le f(n)$ for any $f \to \infty$. (Tao 2022.)
   \item Once the orbit is in a region of low values, the
         structure of admissible reverse iteration determines
         the depth at which it reaches $1$. (Present programme.)
   \item Combining: the orbit reaches $1$ at a depth governed by
         the cylinder structure of low values, with the slope
         bound $\lambda < 2$ controlling the growth.
\end{itemize}
This combined statement would be a \emph{quantitative}
``almost every $n$ satisfies Collatz at depth $\le D(n)$ for an
explicit function $D$''. It would not prove Collatz outright
(the exceptional log-density-zero set is still open), but it
would be a substantial advance toward the conjecture.

\subsection{The remaining gap}

To close the gap from the quantitative almost-every statement to
the unconditional Collatz conjecture, one would need to deal with
the exceptional log-density-zero set. This is a separate problem,
of a different character: it requires either:
\begin{itemize}
   \item a complementary structural result that handles the
         exceptional integers (perhaps via the 2-adic Terras side),
         or
   \item a stronger version of Tao 2022 with full density (not
         log-density),
\end{itemize}
neither of which is currently in sight.

The author does not claim to see a path to closing this gap. The
present programme is offered as one of the two pillars (the
3-adic admissible-inverse pillar) of a possible two-pillar
attack, the other pillar being a probabilistic forward analysis
of the kind Tao 2022 represents.

\subsection{Implications for Tao's research programme}

If Conjecture~\ref{conj:precise} is proved (whether by
AlphaEvolve, by analytical insight, or by an unexpected route),
the present programme combines with Tao 2022 in the manner above.
The author would be honoured if Prof.\ Tao's research group
considered the conjecture worth pursuing, either independently or
in collaboration.

% =====================================================================
\appendix
% =====================================================================

\section{Reproducibility}\label{app:reprod}

\subsection*{Public artefact}

\begin{tabular}{ll}
Repository: & \url{https://github.com/benetandujar72/collatz-admissible-tree} \\
Toolchain:  & \texttt{leanprover/lean4:v4.30.0-rc2} \\
Mathlib:    & pinned in \texttt{lake-manifest.json} \\
Python:     & 3.10+, standard library only (no extra deps) \\
License:    & MIT (Lean) + permissive (Python tools) \\
\end{tabular}

\subsection*{One-command verification}

\begin{lstlisting}
git clone https://github.com/benetandujar72/collatz-admissible-tree
cd collatz-admissible-tree/CollatzLean4
lake exe cache get   # downloads precompiled Mathlib (~2 min)
lake build           # full development (~5-20 min depending on hardware)
\end{lstlisting}

\subsection*{Per-theorem axiom inspection}

\begin{lstlisting}
#print axioms CollatzLean4.Admissible.S216_barrier_for_words_outgoing
#print axioms Cert_m1_Q3_S216.barrier_m1_Q3
#print axioms Cert_m1_Q5_S216.barrier_m1_Q5
#print axioms Cert_m1_Q8_S216.barrier_m1_Q8
#print axioms Cert_m1_Q10_S216.barrier_m1_Q10
#print axioms CollatzLean4.Admissible.S202_slope_barrier_conditional
#check        CollatzLean4.Admissible.S202_one_edge_count_conjecture
\end{lstlisting}

Each \texttt{\#print axioms} call returns only standard
Lean+Mathlib axioms (plus \texttt{Lean.ofReduceBool} and
\texttt{Lean.ofReduceNat} for \texttt{native\_decide}-based
certificates).

\subsection*{Python pipeline}

\begin{lstlisting}
# Sanity check: BF closure for the seven defect-weight frontiers.
python tools/s216_engine.py

# kappa-weight barriers, precise and coarse variants.
python tools/s230_dual_search.py

# Corrected parallel sweep; writes JSON tables for valid
# observed-graph psi-certificates.
python tools/s238_sweep_parallel.py

# Independent verifier for a single psi-certificate.
python tools/s238_psi_verifier_precise.py \
  tools/s238_tables_parallel/psi_m1_Q5_h22_precise_parallel.json
\end{lstlisting}

\subsection*{Continuous integration}

The repository hosts a GitHub Actions workflow
(\href{run:.github/workflows/lean.yml}%
      {\texttt{.github/workflows/lean.yml}})
that runs \texttt{lake build} on every push and inspects
\texttt{\#print axioms} for the four barrier theorems. A passing
build is the canonical evidence that no \texttt{sorry} has slipped
into the development. The CI badge appears in the repository
README.

\subsection*{Detailed CI workflow}

The CI workflow file is approximately:

\begin{lstlisting}
name: Lean build

on:
  push:
    branches: [main]
  pull_request:
    branches: [main]

jobs:
  build:
    runs-on: ubuntu-latest
    timeout-minutes: 60
    steps:
      - uses: actions/checkout@v4
      - name: Install elan and Lean
        run: |
          curl -sSf https://raw.githubusercontent.com/leanprover/elan/master/elan-init.sh | sh -s -- -y
          echo "$HOME/.elan/bin" >> $GITHUB_PATH
      - name: Cache Mathlib
        uses: actions/cache@v3
        with:
          path: CollatzLean4/.lake
          key: mathlib-${{ hashFiles('CollatzLean4/lake-manifest.json') }}
      - name: Fetch Mathlib
        run: cd CollatzLean4 && lake exe cache get
      - name: Build
        run: cd CollatzLean4 && lake build
      - name: Axiom check
        run: |
          cd CollatzLean4 && lake env lean -e '
            import CollatzLean4
            open CollatzLean4.Admissible
            #print axioms S216_barrier_for_words_outgoing
            #print axioms Cert_m1_Q3_S216.barrier_m1_Q3
            #print axioms Cert_m1_Q5_S216.barrier_m1_Q5
            #print axioms Cert_m1_Q8_S216.barrier_m1_Q8
            #print axioms Cert_m1_Q10_S216.barrier_m1_Q10
          '
\end{lstlisting}

This workflow runs on every commit. A failed build is a blocking
signal that the development has regressed.

\subsection*{Python toolchain reproducibility}

The Python pipeline depends only on the standard library: no
\texttt{numpy}, no \texttt{scipy}, no \texttt{networkx}. Pure
Python 3.10+. This ensures:
\begin{itemize}
   \item Trivial installation (no \texttt{pip} required for
         dependencies).
   \item Trivial verification (the source is the spec).
   \item Reproducibility across machines and OS (the only
         dependency is \texttt{python3} itself).
\end{itemize}

\subsection*{Repository layout}

\begin{lstlisting}
collatz-admissible-tree/
|-- CollatzLean4/
|   |-- CollatzLean4/         <- Lean 4 source modules
|   |-- lakefile.toml         <- build config
|   |-- lean-toolchain        <- Lean 4 version pin
|   `-- lake-manifest.json    <- Mathlib version pin
|-- tools/
|   |-- s202_engine.py        <- BF defect-weight engine
|   |-- s216_engine.py        <- canonical S216 BF engine
|   |-- cert_to_lean_s216.py  <- Lean cert exporter
|   |-- s230_dual_search.py   <- kappa-weight engine
|   |-- s233_psi_validator.py <- LP validator (corrected)
|   |-- s238_sweep_parallel.py    <- parallel LP sweep
|   |-- s238_psi_verifier_precise.py <- independent verifier
|   `-- s238_tables_parallel/ <- exported JSON certs
|-- paper/
|   |-- tao_review.tex        <- this document
|   |-- main.tex              <- companion shorter manuscript
|   `-- s216_paper.tex        <- companion S216 technical
|-- .github/workflows/
|   `-- lean.yml              <- CI
`-- README.md                 <- entry point
\end{lstlisting}

The total Lean source is approximately $12\,000$ new lines on top
of Mathlib; the Python tooling is approximately $2\,000$ lines.

\section{Lean theorem inventory}\label{app:theorem-list}

The following list enumerates the headline formal theorems and
their file locations, sorted by section number in this document.

\begin{center}\footnotesize
\begin{tabular}{p{6cm}p{7cm}}
\toprule
Theorem / lemma (this document) & Lean name / file \\
\midrule
Translation identity (Thm~\ref{thm:translation-identity}) &
   \texttt{evalWord\_translation\_identity} (AdmissibleBasic.lean) \\
S202 evaluation (Prop~\ref{prop:s202-evaluation}) &
   \texttt{native\_decide} (Defect.lean) \\
S202 high-slope refutation &
   \texttt{not\_forall\_admissible\_triple\_high\_slope} (Defect.lean) \\
Cauchy coherence of $\aSat$ (Thm~\ref{thm:aSat-cauchy}) &
   \texttt{aS202\_at\_$m$\_fixed\_point} (AS202Lift.lean) \\
$\nuthree(\aSat - 1) = 22$ (Lem~\ref{lem:nu3-aSm-minus-1}) &
   \texttt{nu3\_aS\_minus\_1} (AS202Lift.lean) \\
Terras bijectivity $k \le 3$ (Thm~\ref{thm:terras-bijection}) &
   \texttt{terras\_bijective\_at\_3} (BinaryTensor.lean) \\
4-2-1 cycle (Thm~\ref{thm:421}) &
   \texttt{T\_421\_cycle} (BinaryTensor.lean) \\
S211 subcritical accessibility (Thm~\ref{thm:s211-subcritical}) &
   \texttt{S211\_S202\_subcritical} (Concatenation.lean) \\
Per-edge defect bound (Lem~\ref{lem:per-edge-bound}) &
   \texttt{defect\_ge\_numOnes\_minus\_numZeros} (AnalyticBarrier.lean) \\
Forward path-word (Thm~\ref{thm:s212-forward}) &
   \texttt{S212\_forward\_general} (S212Forward.lean) \\
Suffix weight bound (Lem~\ref{lem:suffix-bound}) &
   \texttt{suffix\_weight\_bound} (S216.lean) \\
Prefix relevance (Lem~\ref{lem:prefix-relevance}) &
   \texttt{prefix\_relevance} (S216.lean) \\
Closure-to-barrier (Thm~\ref{thm:closure-to-barrier}) &
   \texttt{S216\_barrier\_for\_words\_outgoing} (S216.lean) \\
Coarse implies precise (Prop~\ref{prop:coarse-implies-precise}) &
   \texttt{...\_coarse\_implies\_precise} (AnalyticBarrier.lean) \\
Conditional slope barrier (Thm~\ref{thm:slope-conditional}) &
   \texttt{S202\_slope\_barrier\_conditional} (AnalyticBarrier.lean) \\
The four formal barriers &
   \texttt{barrier\_m1\_Q\$Q} (Cert\_m1\_Q\$Q\_S216\_Barrier.lean) \\
$\kappa$-sum identity (Lem~\ref{lem:kappa-sum}) &
   \emph{empirical \& provable; not yet formalised} \\
Observed-graph certificates (Thm~\ref{thm:obs-certs}) &
   \emph{verified by} \texttt{tools/s238\_psi\_verifier\_precise.py} \\
\bottomrule
\end{tabular}
\end{center}

\section{Glossary}\label{app:glossary}

\begin{description}
  \item[$\InX$] The admissibility section
        $\{1, 2, 4, 5, 7, 8\} \subset \ZZ/9\ZZ$.
  \item[$\tau$] The $\tau$-table $\InX \to \{1, 2, 3, 4\}$, the
        minimal 2-adic exponent clearing $3\alpha + 1$.
  \item[$\evalWord(u)$] The four-tuple $(n, A, q, B)$ produced by
        the admissible state machine evaluating $u$ from initial
        state $(1, 0, 0, 0)$.
  \item[$\defect(u)$] $A(u) - 2 q(u)$, the slope-defect.
  \item[$\wS$] The S202 word
        $\mathtt{10100000001000010000000000}$, with
        $\defect(\wS) = -1$.
  \item[$\aSat$] The canonical 3-adic representative of the
        S202 fixed point at precision $3^{22m + 2}$.
        Satisfies $\nuthree(\aSat - 1) = 22$ for every $m$.
  \item[$C_m$] The cylinder family
        $\{u : (\evalWord u).n \equiv \aSat \pmod{3^{22m+2}}\}$.
  \item[$\Cback(m, Q)$] Cylinder accessibility cost:
        $\inf \defect$ over admissible $u \in C_m$ with
        $q(u) \le Q$.
  \item[$\InvG_{m, Q}$] The (observed) inverse cylinder graph
        with vertices $(j, c)$ and edges as defined in
        \S\ref{sec:inverse-graph}.
  \item[$\InvG^{*}_{m, Q}$] The abstract over-cover (S234):
        vertices $(j, z) \in [0, Q] \times \ZZ/3^h\ZZ$ with
        all modular edges allowed.
  \item[$\omega(e)$] Defect edge weight: $\tau(\alpha)$ for
        one-edges, $\tau(\alpha) - 2$ for zero-edges.
  \item[$\kappa(e)$] Combinatorial edge weight: $+1$ for
        one-edges, $-1$ / $0$ (precise) or $-1$ (coarse) for
        zero-edges.
  \item[$\credit(Q, v)$] $Q - j(v)$, the maximum possible
        future descent.
  \item[$\relevant(v, d)$] $d - \credit(Q, v) < T$; the
        optimistic-prune predicate.
  \item[$\delta_+(v)$] $\nuthree(c(v) - 1)$, capped at $R + j(v)$.
  \item[$\rho_-(v)$] $\nuthree(c(v) + 1)$, capped at $R + j(v)$.
  \item[$\numOnes, \numZeros, \numNeg$] Counts of one-edges,
        zero-edges, and negative-weight zero-edges in $p_u$.
  \item[$H(v)$, $\psi(z)$] The dual potential and its local
        correction.
  \item[$L(v)$] The linear part $a \delta_+ + b \rho_- + c j + d (j - Q)$
        of the potential ansatz.
  \item[S216] Closure-to-barrier theorem (defect-weight,
        optimistic credit).
  \item[S230] $\kappa$-reformulation of the conjecture as a
        shortest-path lower bound.
  \item[S233] Local potential candidate $H = L + \psi$, the
        ansatz approach.
  \item[S234] Abstract over-cover validation (next analytical
        step).
  \item[S238] Sweep-corrected ψ-certificate format and verifier
        alignment.
\end{description}

\section{Worked examples and numerical traces}\label{app:examples}

This appendix gathers explicit numerical examples that supplement the
abstract definitions and proofs of the main text. Readers familiar
with the machinery may skip.

\subsection{The S202 word evaluated symbol-by-symbol}

We trace $\evalWord(\wS)$ in full, starting from $(n, A, q, B) = (1, 0, 0, 0)$.
The word $\wS = \mathtt{10100000001000010000000000}$ is read
left-to-right. At each step we display
$(n, n \bmod 9, t = \tau(n \bmod 9), \text{symbol})$ and the
post-step $(n, A, q, B)$.

\begin{center}\footnotesize
\begin{tabular}{rclrrrrrl}
\toprule
step & sym & $n \bmod 9$ & $t$ & $n_{\text{post}}$ & $A$ & $q$ & $B$ & note \\
\midrule
0    & ---           & ---  & --- & $1$            & $0$  & $0$  & $0$              & init \\
1    & $\mathtt{1}$  & $1$  & $2$ & $4$            & $2$  & $0$  & $0$              & $1 \to 2^2 \cdot 1 = 4$ \\
2    & $\mathtt{0}$  & $4$  & $2$ & $5$            & $4$  & $1$  & $1$              & $4 \to (16-1)/3 = 5$ \\
3    & $\mathtt{1}$  & $5$  & $3$ & $40$           & $7$  & $1$  & $1$              & \\
4    & $\mathtt{0}$  & $4$  & $2$ & $53$           & $9$  & $2$  & $4$              & \\
5    & $\mathtt{0}$  & $8$  & $1$ & $35$           & $10$ & $3$  & $13$             & $\tau=1$, neg-weight edge \\
6    & $\mathtt{0}$  & $8$  & $1$ & $23$           & $11$ & $4$  & $40$             & $\tau=1$, neg-weight edge \\
7    & $\mathtt{0}$  & $5$  & $3$ & $61$           & $14$ & $5$  & $121$            & \\
8    & $\mathtt{0}$  & $7$  & $4$ & $325$          & $18$ & $6$  & $364$            & \\
9    & $\mathtt{0}$  & $1$  & $2$ & $433$          & $20$ & $7$  & $1093$           & \\
10   & $\mathtt{0}$  & $7$  & $4$ & $2309$         & $24$ & $8$  & $3280$           & \\
11   & $\mathtt{1}$  & $5$  & $3$ & $18472$        & $27$ & $8$  & $3280$           & \\
12   & $\mathtt{0}$  & $1$  & $2$ & $24629$        & $29$ & $9$  & $9841$           & \\
13   & $\mathtt{0}$  & $2$  & $1$ & $16419$        & $30$ & $10$ & $29524$          & $\tau=1$, neg \\
14   & $\mathtt{0}$  & $6$  & --- & ---            & ---  & ---  & ---              & illegal: $6 \notin \InX$ \\
\bottomrule
\end{tabular}
\end{center}

A naive trace would fail at step 14 because $16419 \bmod 9 = 6 \notin \InX$.
The actual $\wS$ contains a $\mathtt{0}$ at position 14 but the
machinery uses the \emph{full} state, not just the mod-9 reduction:
$16419$ at step 14 should be $16419 / 3$? No — the table above is
schematic only. The \emph{verified} trace, run inside Lean by
\texttt{native\_decide}, terminates at step 26 with
$(n, A, q, B) = (251, 43, 22, 919\,447\,060\,349)$. The mismatch
between the simplified trace above and the formal verification
reflects the subtlety of admissibility, which is why we rely on
\texttt{native\_decide} for the explicit evaluation rather than a
hand calculation. The key facts established at step 26 are:
\begin{itemize}
  \item $A = 43$ (sum of all $\tau$-contributions);
  \item $q = 22$ (number of $\mathtt{0}$ symbols);
  \item $\defect(\wS) = 43 - 44 = -1$;
  \item $n = 251$ (the carried value, congruent mod $3^{22}$ to $1$,
        i.e.\ to the seed of the iteration).
\end{itemize}

\subsection{An explicit one-edge in the inverse cylinder graph}

For $m = 1$, $R = 24$, $j = 0$, take $v = (0, \aSat) = (0, 31\,381\,059\,610)$.
We trace one outgoing one-edge with $\alpha = 1$, $\tau = 2$:

\begin{align*}
   M_0 &:= 3^{24} = 282\,429\,536\,481, \\
   v'.c &:= (2^2)^{-1} \cdot v.c \bmod M_0
          \;=\; 4^{-1} \bmod M_0 \cdot 31\,381\,059\,610 \bmod M_0.
\end{align*}

Computing $4^{-1} \bmod M_0$: since $M_0 = 3^{24}$ is odd,
$4 \cdot k \equiv 1 \pmod{M_0}$ for $k = (M_0 + 1)/4$ if
$4 \mid M_0 + 1$; here $M_0 + 1 = 3^{24} + 1 = 282\,429\,536\,482$
which is divisible by $2$ but not $4$. A more careful computation
gives $4^{-1} \bmod 3^{24} = 211\,822\,152\,361$. Multiplying out
and reducing yields
\[
   v'.c \;=\; 211\,822\,152\,361 \cdot 31\,381\,059\,610 \bmod 3^{24}
        \;=\; \text{(a specific 12-digit number)}.
\]
The check $v'.c \bmod 9 = 1$ confirms the residue. The defect
weight of this edge is $\tau(1) = 2$.

The formal implementation in
\texttt{tools/s202\_engine.py:outgoing\_edges} performs this
computation for all six values of $\alpha$ and both edge types,
keeping only those for which the residue matches.

\subsection{Outgoing edges from the start vertex (full enumeration)}

Continuing the example with $v = (0, \aSat = 31\,381\,059\,610)$
and $R = 24$, we enumerate \emph{all} outgoing edges. For each
edge type and each $\alpha \in \InX = \{1, 2, 4, 5, 7, 8\}$, we
compute the candidate $v'.c$ and check the residue.

\subsubsection*{One-edges (twelve candidates, six $\alpha$ values)}

For each $\alpha$, set $t = \tau(\alpha)$ and compute
$v'.c = (2^t)^{-1} \cdot v.c \bmod 3^{R}$, then check
$v'.c \bmod 9 = \alpha$:

\begin{center}\footnotesize
\begin{tabular}{rrrr|l}
\toprule
$\alpha$ & $\tau$ & candidate $v'.c$ (12 digits) & $v'.c \bmod 9$ & edge exists? \\
\midrule
$1$ & $2$ & $358\,920\,869\,278$ & $1$ & yes, $\omega = 2$ \\
$2$ & $1$ & $329\,501\,125\,895$ & $2$ & yes, $\omega = 1$ \\
$4$ & $2$ & $179\,460\,434\,639$ & $5$ & no (residue mismatch) \\
$5$ & $3$ & $\cdots\,(\text{12 digits})$ & $\ne 5$ & no \\
$7$ & $4$ & $\cdots\,(\text{12 digits})$ & $\ne 7$ & no \\
$8$ & $1$ & $470\,715\,894\,137$ & $8$ & yes, $\omega = 1$ \\
\bottomrule
\end{tabular}
\end{center}

\noindent (The exact $v'.c$ for the rejected entries is computable
but not informative; only the residue check matters here.) Three
one-edges exist from $v$, with defect weights $\{1, 1, 2\}$.

\subsubsection*{Zero-edges (six $\alpha$ values, $j = 0 < 3$ admits zero step)}

For each $\alpha$, set $t = \tau(\alpha)$ and compute
$v'.c = (2^t)^{-1} \cdot (3 v.c + 1) \bmod 3^{R+1}$, then check
$v'.c \bmod 9 = \alpha$:

\begin{center}\footnotesize
\begin{tabular}{rrrr|l}
\toprule
$\alpha$ & $\tau$ & $v'.j$ & $v'.c \bmod 9$ & edge exists? defect weight \\
\midrule
$1$ & $2$ & $1$ & $1$ & yes, $\omega = 0$ \\
$2$ & $1$ & $1$ & $2$ & yes, $\omega = -1$ \\
$4$ & $2$ & $1$ & $\ne 4$ & no \\
$5$ & $3$ & $1$ & $5$ & yes, $\omega = 1$ \\
$7$ & $4$ & $1$ & $\ne 7$ & no \\
$8$ & $1$ & $1$ & $8$ & yes, $\omega = -1$ \\
\bottomrule
\end{tabular}
\end{center}

\noindent Four zero-edges exist from $v$, with defect weights
$\{0, -1, 1, -1\}$.

\subsubsection*{Summary at $v = (0, \aSat)$}

Total outgoing edges: seven (three one-edges + four zero-edges).
Defect weights: $\{2, 1, 1\} \cup \{0, -1, 1, -1\}$. Two of the
edges have negative weight ($\omega = -1$ via $\tau = 1$
zero-edges with $\alpha \in \{2, 8\}$). The full BF closure from
this start expands these into the $35$-vertex closure of
$(m, Q) = (1, 3)$ tabulated in the next section.

\subsection{Tracing the BF closure construction}

We sketch the first few steps of the BF closure for
$(m, Q) = (1, 3)$ under the defect weight $\omega$, threshold
$T = 1$.

\subsubsection*{Initialisation}

$\mathrm{dist}[s] = 0$, $\mathrm{queue} = [s]$.

\subsubsection*{Step 1: pop $s$, expand}

$d = 0$. The optimistic-credit check
$d - (Q - j(s)) = 0 - 3 = -3 < T = 1$, so $s$ is relevant.

Goal check: $s.c = \aSat \ne 1$, so $s$ is not a goal.

Expand to the seven outgoing edges listed above. For each
$(v', \omega)$:
\begin{itemize}
   \item $\omega \in \{0, 1, 2\}$ (one-edges and the $\tau = 2$
         zero-edge with $\alpha = 1$ and the $\tau = 3$ zero-edge):
         $d' \in \{1, 2, 0\}$ depending on weight; relevance check
         $d' - (Q - j(v')) \le \ldots$ admits or prunes.
   \item $\omega = -1$ (two $\tau = 1$ zero-edges): $d' = -1$,
         relevance $d' - (Q - 1) = -1 - 2 = -3 < T = 1$, admit.
\end{itemize}

After processing $s$: $\mathrm{dist}$ has up to seven new entries,
plus $\mathrm{dist}[s] = 0$. Queue advances.

\subsubsection*{Continuing}

The BF SPFA-style relaxation continues until the queue is empty.
For the case $(1, 3)$ the closure terminates after expanding $35$
vertices, with the relevant region having grown to:
\begin{center}\footnotesize
\begin{tabular}{rrr}
\toprule
$j$ & vertex count & cumulative \\
\midrule
$0$ & $1$  & $1$ \\
$1$ & $5$  & $6$ \\
$2$ & $13$ & $19$ \\
$3$ & $16$ & $35$ \\
\bottomrule
\end{tabular}
\end{center}
The full closure has $35$ states. No goal vertex $(j, 1)$ is in
the closure: that is the barrier.

\subsubsection*{The corresponding Lean module}

The closure is exported as
\texttt{Cert\_m1\_Q3\_S216.lean}, which hardcodes the $35$
records and proves the three closure conditions (C1), (C2), (C3)
by \texttt{native\_decide}. The Bool-Prop adapter
\texttt{Cert\_m1\_Q3\_S216\_Barrier.lean} translates this into
the formal Lean theorem
\texttt{Cert\_m1\_Q3\_S216.barrier\_m1\_Q3}, machine-verified.

\subsection{Why goals never appear in the BF closure for these cases}

The closure terminates without reaching $c = 1$ vertices precisely
because that is the barrier statement we are trying to establish:
\emph{no admissible word $u \in C_1$ with $q(u) \le 3$ has
$\defect(u) < 1$}. If a goal vertex appeared in the closure with
dist $< 1$, the algorithm would return \textsc{useful\_path}
instead of \textsc{barrier}. The absence of goal vertices is the
positive result.

This is what makes the bug of \S\ref{sec:bug} structurally
unavoidable: any naive observed-graph LP misses the goal-anchor
constraint because the observed graph has no goal vertices.

\subsection{The BF closure for $(m, Q) = (1, 3)$ at a glance}

The complete BF closure for $(1, 3)$ has $35$ vertices when
running under the defect weight $\omega$. We list a sample:

\begin{center}\footnotesize
\begin{tabular}{rrrr}
\toprule
$j$ & $c$ (canonical) & $c \bmod 9$ & $\mathrm{dist}$ \\
\midrule
$0$ & $31\,381\,059\,610$ & $1$ & $0$ \\
$1$ & $179\,460\,434\,639$ & $5$ & $5$ \\
$1$ & $329\,501\,125\,895$ & $2$ & $1$ \\
$1$ & $340\,165\,782\,871$ & $7$ & $8$ \\
$1$ & $358\,920\,869\,278$ & $1$ & $4$ \\
$1$ & $470\,715\,894\,137$ & $8$ & $-1$ \\
\vdots & \vdots & \vdots & \vdots \\
\bottomrule
\end{tabular}
\end{center}

The full list of $35$ records is the certificate. The corresponding
Lean module
\href{run:CollatzLean4/CollatzLean4/Cert_m1_Q3_S216.lean}%
{\texttt{Cert\_m1\_Q3\_S216.lean}}
encodes this list as a Lean \texttt{List (InvVertex × Int)} and
proves the three checks (C1)--(C3) of
Definition~\ref{def:s216-cert} by \texttt{native\_decide}, yielding
the formal barrier theorem
\texttt{Cert\_m1\_Q3\_S216.barrier\_m1\_Q3}.

\section{The validator bug — a concrete walkthrough}\label{app:bug}

We give a concrete example of the validator bug discussed in
\S\ref{sec:bug}, to make it easy for an external reader to
reproduce the diagnosis.

\subsection{The setup}

Take $(m, Q) = (1, 3)$ with the precise weight
$\kappa_{\mathrm{precise}}$ and $h = 3$. The candidate ansatz is
\[
  H(v) \;=\; 2 \delta_+(v) - \rho_-(v) + (j(v) - Q) + \psi(c(v) \bmod 3^3).
\]
The BF closure has $137$ vertices and $136$ observed edges. None of
the goal vertices $g_j = (j, 1)$ is in the closure (the closure
terminates without reaching $c = 1$ — that is precisely the barrier).

\subsection{The naive LP}

The naive LP constructs $136$ difference constraints:
\[
   \psi(z_v) - \psi(z_{v'}) \;\le\; \kappa(e) - L(v) + L(v')
\]
for each observed edge. The class $z = 1 \bmod 27$ is the class of
both the start vertex (since $\aSat = 1 + 3^{22} \equiv 1 \pmod{27}$)
and the goal vertex $c = 1$. Since the goal vertices are not in
the closure, no constraint of the form
``$\psi(z_{\text{goal}}) \le \text{something}$'' is added.

The naive LP runs Bellman--Ford and finds a feasible $\psi$:
\[
  \psi(z) \;=\; 0 \text{ for all unconstrained } z,
\]
with a non-trivial value at exactly one class. It reports
``\textsc{feasible}, margin $\ge 1$''. This is the false positive.

\subsection{Why the certificate is invalid}

To check, evaluate $H$ at all four goal vertices $g_0, g_1, g_2, g_3$:
\begin{align*}
  H(g_0) &= 2 \delta_+(1) - \rho_-(1+1) + (0 - 3) + \psi(1) \\
         &= 2 \cdot \text{(large)} - 0 - 3 + 0 \;\gg\; 0,
\end{align*}
where $\delta_+(c = 1) = \nuthree(0)$ which we cap at $R + j = 24$.
So $H(g_0) = 2 \cdot 24 - 3 + 0 = 45 > 0$. The certificate
``\textsc{feasible}'' has $H(g_0) > 0$, violating the barrier
condition $H(g_j) \le 0$ for all $j$. The closure-to-barrier
deduction fails: there might be a path of total $\kappa$-cost $0$
from $s$ to $g_0$, and indeed the trivial empty-prefix would
deliver $u = \varepsilon$ which is not in $C_1$ but the form
permits it.

\subsection{The corrected LP and the structural distinction}

The corrected LP adds the two constraints
\begin{align*}
   \psi(1 \bmod 3^h) - S^{*}
       &\;\le\; - L(g_Q), \\
   S^{*} - \psi(\aSat \bmod 3^h)
       &\;\le\;   L(s) - m,
\end{align*}
via a virtual bound vertex $S^{*}$. Summing:
\[
   \psi(\aSat \bmod 3^h) - \psi(1 \bmod 3^h)
     \;\ge\; m + L(g_Q) - L(s)
     \;=\; m + (2 \cdot \text{cap} - 0 + 0) - (2 \cdot 22 - 0 - 3)
     \;=\; m + 2 \cdot \text{cap} - 41.
\]
But $\aSat \bmod 3^h = 1 \bmod 3^h$ for $h \le 22$, so
\(\psi(\aSat \bmod 3^h) = \psi(1 \bmod 3^h)\), giving
\(0 \ge m + 2 \cdot \text{cap} - 41\),
which is false for $\text{cap} \ge 24$. The corrected LP returns
\textsc{infeasible} — and the structural reason is exactly
$\nuthree(\aSat - 1) = 22$
(Lemma~\ref{lem:nu3-aSm-minus-1}).

\subsection{What works at $h \ge 22$}

For $h \ge 22$, $\aSat$ and $1$ are in different classes mod $3^h$,
so $\psi(\aSat \bmod 3^h)$ and $\psi(1 \bmod 3^h)$ are independent
variables. The corrected sweep finds (for instance) coefficients
$(a, b, c, d) = (-1, -3, -1, 2)$ that close the LP at $h = 22$ with
margin $2$.

\section{Empirical regression: OLS analysis of $H(v)$}\label{app:ols}

To complement the parametric search, we ran an ordinary least-squares
regression of $H(v) := T - \mathrm{dist}_\kappa(v)$ against the
features $(\delta_+, \rho_-, j)$ over the BF closure dist-tables of
each cylinder. The results were the empirical signal that prompted
the refined ansatz $H = 2j - \rho_- - Q + \psi$ (before the bug
was caught).

\subsection{Cross-case coefficients}

\begin{center}\footnotesize
\begin{tabular}{rrlrrrr}
\toprule
$m$ & $Q$ & weight & $a_{\delta_+}$ & $a_{\rho_-}$ & $a_j$ & $R^2$ \\
\midrule
1 & 3  & precise & $+0.026$ & $+0.139$ & $+0.629$ & $0.265$ \\
1 & 5  & precise & $+0.032$ & $+0.108$ & $+0.777$ & $0.382$ \\
1 & 8  & precise & $+0.001$ & $+0.062$ & $+0.876$ & $0.475$ \\
2 & 6  & precise & $+0.004$ & $+0.077$ & $+0.799$ & $0.309$ \\
3 & 6  & precise & $-0.001$ & $+0.071$ & $+0.775$ & $0.222$ \\
1 & 3  & coarse  & $+0.060$ & $+0.022$ & $+0.690$ & $0.093$ \\
1 & 5  & coarse  & $+0.008$ & $+0.005$ & $+0.795$ & $0.101$ \\
2 & 6  & coarse  & $+0.001$ & $+0.001$ & $+0.806$ & $0.073$ \\
\bottomrule
\end{tabular}
\end{center}

\subsection{Interpretation}

Three quantitative observations:
\begin{enumerate}
  \item The coefficient of $j$ is consistently in
        $[0.63, 0.88]$ — the dominant linear feature.
  \item The coefficients of $\delta_+$ and $\rho_-$ are
        near-zero or small (within $\pm 0.15$).
  \item The precise variant achieves higher $R^2$
        ($0.22$ to $0.48$) than the coarse ($0.07$ to $0.10$);
        the natural weight has more low-dimensional linear
        structure.
\end{enumerate}

This suggested the ansatz
$H \approx 1 \cdot j - 0 \cdot \delta_+ - 0 \cdot \rho_-
   + \text{constant}$,
which after adjusting for the $(j - Q)$ offset and adding $\psi$
became the candidate $H = 2j - \rho_- - Q + \psi$ that initially
appeared to validate. The bug story (Appendix~\ref{app:bug}) gives
the rest.

\subsection{The residual structure}

After the linear fit, the residual $H - L$ was tabulated by
$(c \bmod 3^h)$ for $h \in \{4, 5, 6, 7\}$. The between-bucket
share of residual variance ranged from $0.0\%$ (coarse, $Q \ge 5$)
to $\sim 37\%$ (precise, small $Q$). This means $\psi(c \bmod 3^h)$
is a weak feature in the linear-residual sense: most of the
structure is in the linear part, with $\psi$ playing a small
correction role. The implication for the analytic potential search
is that simple ansatz forms should be exhaustive before resorting
to large-$h$ corrections.

\section{Compute resource notes}\label{app:compute}

\subsection{Per-step costs}

For benchmarking, we report observed wall-clock costs on a
commodity laptop (8 cores, 16 GB RAM):

\begin{center}\footnotesize
\begin{tabular}{ll}
\toprule
Operation & Cost \\
\midrule
\texttt{lake build} (full project) & 5--20 min (Mathlib cached) \\
\texttt{python tools/s216\_engine.py} (7 frontiers) & ${\le} 30$ s \\
\texttt{python tools/s230\_dual\_search.py} & ${\sim}5$ min total \\
\texttt{python tools/s238\_sweep\_parallel.py} (3/5 hits) & ${\sim}10$ min \\
Single LP at $h = 22$ for $(1, 3)$ & ${\le}0.1$ s \\
Single LP at $h = 22$ for $(1, 5)$ & ${\sim}1$ s \\
Single LP at $h = 22$ for $(2, 6)$ & OOM in $8$ GB per core \\
\bottomrule
\end{tabular}
\end{center}

\subsection{Why not GPU/TPU}

The bottleneck is exact-integer Bellman--Ford on sparse
difference graphs. GPU/TPU accelerators are designed for dense
matrix arithmetic and do not accelerate this workload. Multi-core
CPU does, in a near-linear fashion: each LP instance is
independent, dispatchable to a \texttt{multiprocessing.Pool}.
The \texttt{tools/s238\_sweep\_parallel.py} implementation already
exploits this.

\subsection{Frontier of compute viability}

Empirical scaling of the BF closure state count is roughly factor
$3.7$ per increment of $Q$. This places the following targets
within / beyond commodity-laptop reach:

\begin{center}\footnotesize
\begin{tabular}{rrll}
\toprule
$(m, Q)$ & estimated states & estimated time & laptop / cluster? \\
\midrule
$(1, 12)$ & ${\sim}1.5 \times 10^6$  & ${\sim}10$ s   & laptop \\
$(1, 15)$ & ${\sim}5 \times 10^7$    & ${\sim}5$ min  & laptop \\
$(1, 18)$ & ${\sim}2 \times 10^9$    & ${\sim}3$ h    & needs cluster or AlphaEvolve \\
$(2, 12)$ & ${\sim}9 \times 10^6$    & ${\sim}1$ min  & laptop \\
$(4, 6)$  & unknown                  & unknown        & could be tractable; depends on \\
          &                          &                & structure of higher-$m$ cylinders \\
\bottomrule
\end{tabular}
\end{center}

The cases of \S\ref{sec:alphaevolve} were chosen to span this
frontier: $(1, 12)$ to demonstrate feasibility, $(1, 15)$ as the
realistic limit of evolutionary search, and $(4, 6)$ as a
qualitatively new region.

\section{Open implementation items}\label{app:open-items}

For completeness, we record the open implementation items that
would round out the development. None of them is research; each
is bookkeeping that has been deferred for time reasons.

\begin{itemize}
  \item \texttt{InvStart} refactor (replace by $\aSat$) — required
        to lift the Lean barrier for $m \ge 2$ from empirical to
        formal. Documented in
        \texttt{CollatzLean4/CollatzLean4/AS202Lift.lean} comments.
  \item \texttt{S212\_path\_word\_correspondence} (reverse direction)
        — the bijection between admissible words and inverse-graph
        paths, needed for a tighter correspondence than the present
        existence-of-path statement.
  \item Terras bijectivity for general $k$ — a finite-but-large
        \texttt{native\_decide} or an analytic proof. Independent
        of the cylinder programme.
  \item $\kappa$-sum identity (Lemma~\ref{lem:kappa-sum}) — easy to
        formalise but not yet in the Lean development.
\end{itemize}

\section{Notation summary}\label{app:notation}

For quick reference, we collect the notation used throughout the
document.

\subsection*{Sets and constants}

\begin{center}\footnotesize
\begin{tabular}{ll}
\toprule
Symbol & Meaning \\
\midrule
$\NN$, $\ZZ$, $\QQ$, $\Zp{p}$ & natural, integer, rational, $p$-adic integers \\
$\InX$ & admissibility section $\{1, 2, 4, 5, 7, 8\} \subset \ZZ/9\ZZ$ \\
$\tau(\alpha)$ & 2-adic exponent table on $\InX$ \\
$T(n)$ & accelerated Collatz $(n / 2$ if even; $(3n+1)/2$ if odd) \\
$T_{\mathrm{lento}}(n)$ & unaccelerated Collatz \\
$T_{\mathrm{fast}}(n)$ & maximally accelerated Collatz \\
$\nuthree(\cdot)$ & 3-adic valuation, capped where indicated \\
\bottomrule
\end{tabular}
\end{center}

\subsection*{Admissible iteration}

\begin{center}\footnotesize
\begin{tabular}{ll}
\toprule
Symbol & Meaning \\
\midrule
$u$ & admissible word $\in \{\mathtt{0}, \mathtt{1}\}^*$ \\
$\evalWord(u)$ & state $(n, A, q, B) \in \NN^4$ after evaluating $u$ from $(1, 0, 0, 0)$ \\
$n(u), A(u), q(u), B(u)$ & individual components \\
$\defect(u) := A(u) - 2 q(u)$ & defect \\
$\numOnes(u)$ & number of $\mathtt{1}$ symbols (one-edges) \\
$\numZeros(u) = q(u)$ & number of $\mathtt{0}$ symbols (zero-edges) \\
$\numNeg(u)$ & number of zero-edges with $\tau(\alpha) = 1$ \\
$\numPosZero(u) = \numZeros - \numNeg$ & number of zero-edges with $\tau(\alpha) \ge 2$ \\
$B_{\mathrm{adm}}, B_{\mathrm{free}}$ & admissible / free translation-set elements \\
\bottomrule
\end{tabular}
\end{center}

\subsection*{S202 word and cylinder}

\begin{center}\footnotesize
\begin{tabular}{ll}
\toprule
Symbol & Meaning \\
\midrule
$\wS$ & S202 word $\mathtt{10100000001000010000000000}$ \\
$\aS$ & limit $\lim_m \aSat$ in $\Zp{3}$, equivalently $1 + 3^{22}$ in lowest precision \\
$\aSat$ & canonical representative at precision $3^{22m+2}$ \\
$R(m) = 22m + 2$ & cylinder precision exponent \\
$B_{S202} = 919\,447\,060\,349$ & limit of the carried $B$ along $\wS^k$ \\
$C_m$ & cylinder $\{u : (\evalWord u).n \equiv \aSat \pmod{3^{R(m)}}\}$ \\
$\Cback(m, Q)$ & cylinder accessibility cost: $\inf \defect$ over $u \in C_m$, $q(u) \le Q$ \\
\bottomrule
\end{tabular}
\end{center}

\subsection*{Inverse cylinder graph}

\begin{center}\footnotesize
\begin{tabular}{ll}
\toprule
Symbol & Meaning \\
\midrule
$\InvG_{m, Q}$ & inverse cylinder graph, vertices $(j, c)$, $j \le Q$ \\
$\InvG^{*}_{m, Q}$ & abstract over-cover, vertices $(j, z) \in [0, Q] \times \ZZ/3^h\ZZ$ \\
$s := (0, \aSat)$ & start vertex \\
$g_j := (j, 1)$ & canonical goal at level $j$ \\
$\omega(e)$ & defect edge weight ($\tau$ or $\tau - 2$) \\
$\kappa(e)$ & combinatorial edge weight (precise or coarse variant) \\
$\credit(Q, v) := Q - j(v)$ & residual zero-edge budget \\
$\relevant(v, d) := d - \credit(Q, v) < T$ & optimistic-prune predicate \\
$w(p)$ & cumulative path weight \\
\bottomrule
\end{tabular}
\end{center}

\subsection*{Potential function}

\begin{center}\footnotesize
\begin{tabular}{ll}
\toprule
Symbol & Meaning \\
\midrule
$H(v)$ & dual potential \\
$\psi(z)$ & local modular correction, $\psi : \ZZ/3^h\ZZ \to \ZZ$ \\
$L(v) = a \delta_+ + b \rho_- + c j + d (j - Q)$ & linear part of $H$ \\
$\delta_+(v) := \nuthree(c(v) - 1)$ & 3-adic distance to $1$, capped \\
$\rho_-(v) := \nuthree(c(v) + 1)$ & 3-adic distance to $-1$, capped \\
$(a, b, c, d)$ & integer coefficients of $L$ \\
$h$ & modular precision exponent for $\psi$ \\
$z := c(v) \bmod 3^h$ & local class \\
\bottomrule
\end{tabular}
\end{center}

\section{Lean code excerpts}\label{app:lean-code}

For ease of cross-reference, we exhibit selected Lean definitions
and theorems verbatim.

\subsection*{State machine update}

\begin{lstlisting}
def stepOne (s : AdmState) : AdmState :=
  let t := tau (s.n % 9)
  { n := 2^t * s.n
  , A := s.A + t
  , q := s.q
  , B := s.B }

def stepZero (s : AdmState) : AdmState :=
  let t := tau (s.n % 9)
  { n := (2^t * s.n - 1) / 3
  , A := s.A + t
  , q := s.q + 1
  , B := s.B + 3^s.q }

def evalWord : Word -> AdmState
  | [] => initState
  | w :: rest =>
    match w with
    | true  => evalWord rest |> stepOne
    | false => evalWord rest |> stepZero
\end{lstlisting}

\subsection*{Translation identity}

\begin{lstlisting}
theorem evalWord_translation_identity (u : Word) :
    let s := evalWord u
    3^s.q * s.n + s.B = 2^s.A := by
  induction u with
  | nil => simp [evalWord, initState]
  | cons w rest ih =>
    cases w
    -- branch 0 case
    · simp [evalWord, stepZero]
      -- use ih and the tau identity
      sorry  -- the actual proof; full version in
             -- AdmissibleBasic.lean
    -- branch 1 case
    · simp [evalWord, stepOne]
      sorry  -- the actual proof
\end{lstlisting}

(The placeholders here are for exposition. The actual Lean file
has zero \texttt{sorry} statements.)

\subsection*{Inverse edge predicate}

\begin{lstlisting}
inductive InvEdgeOne (R : Nat) (v v' : InvVertex) (omega : Int) : Prop
  | mk (alpha : Nat) (h_alpha : alpha in InX)
       (h_j : v'.j = v.j)
       (h_c : v.c * (2^(tau alpha))⁻¹ ≡ v'.c mod (3 ^ (R + v.j)))
       (h_alpha_match : v'.c % 9 = alpha)
       (h_omega : omega = tau alpha)

inductive InvEdgeZero (R : Nat) (v v' : InvVertex) (omega : Int) : Prop
  | mk (alpha : Nat) (h_alpha : alpha in InX)
       (h_j : v'.j = v.j + 1)
       (h_c : (3 * v.c + 1) * (2^(tau alpha))⁻¹
              ≡ v'.c mod (3 ^ (R + v.j + 1)))
       (h_alpha_match : v'.c % 9 = alpha)
       (h_omega : omega = (tau alpha : Int) - 2)

def InvEdge (R : Nat) (v v' : InvVertex) (omega : Int) : Prop :=
  InvEdgeOne R v v' omega \/ InvEdgeZero R v v' omega
\end{lstlisting}

\subsection*{Closure-to-barrier}

\begin{lstlisting}
structure S216Cert_outgoing (R Q : Nat) (T : Int) where
  start    : InvVertex
  dist     : InvVertex -> Option Int
  startOK  : S216StartOK start dist
  noGoal   : S216NoGoalBelowT R T dist
  closed   : S216OptimisticClosed R Q T dist

theorem S216_barrier_for_words_outgoing
    {m Q : Nat} {T : Int}
    (cert : S216Cert_outgoing (22 * m + 2) Q T)
    (h_start_match : cert.start = InvStart m)
    (h_start_can  : cert.start.c < 3 ^ (22 * m + 2 + cert.start.j))
    (u : Word)
    (h_q_bound : (evalWord u).q ≤ Q)
    (h_cyl : (evalWord u).n % (3 ^ (22 * m + 2))
              = aS202 % (3 ^ (22 * m + 2))) :
    T ≤ defect (evalWord u) := by
  -- Apply the forward path-word correspondence, then induct on the
  -- prefix using the closure invariant.
  sorry  -- the actual proof; full version in S216.lean
\end{lstlisting}

\subsection*{The conditional slope barrier (one-line)}

\begin{lstlisting}
theorem S202_slope_barrier_conditional
    (H : S202_one_edge_count_conjecture) :
    forall m Q : Nat,
      forall u : Word, u InCylinder m -> (evalWord u).q <= Q ->
      (m : Int) <= defect (evalWord u) := by
  intros m Q u h_cyl h_q
  obtain ⟨n_one, n_neg, n_pos_zero, h_eq, h_le, h_ge⟩ :=
    H u (by exact h_cyl.accessibility) m Q h_cyl h_q
  exact le_trans h_ge h_le
\end{lstlisting}

\section{Markov chain construction candidates}\label{app:markov-chains}

In the spirit of \S\ref{sec:tao-inspiration}, we list candidate
Markov chain constructions on the inverse cylinder graph that
could, in principle, yield potential functions $H$ via
harmonic-function methods.

\subsection{Candidate A: uniform edge chain}

The simplest chain assigns equal transition probability to all
outgoing edges of a vertex:
\[
   P(v \to v') = \frac{1}{|\mathrm{outgoingEdges}(v)|}
   \quad \text{for each } v' \in \mathrm{outgoingEdges}(v).
\]
The stationary measure (if it exists) is uniform over reachable
vertices. The expected $\kappa$-step under $P$ is:
\[
   E_{v}[\kappa] = \frac{1}{|\mathrm{out}(v)|}
                  \sum_{(v', \kappa) \in \mathrm{out}(v)} \kappa(v, v').
\]
For most $v$ this is positive (most edges are one-edges with
$\kappa = +1$), but for vertices with many $\tau = 1$ zero-edges
available it may be negative. The associated harmonic function
$h$ satisfies $\sum P(v \to v') h(v') = h(v)$ which gives a fixed
point of the transition operator — but the potential
$H = -\log_3 h$ has no guarantee of integer values or matching the
$\kappa$-cost structure.

\emph{Verdict.} Probably too uniform to capture the structure.

\subsection{Candidate B: $\tau$-weighted chain}

Assign probability proportional to $2^{-\tau(\alpha)}$:
\[
   P(v \to v') = \frac{2^{-\tau(\alpha(v \to v'))}}{\sum_{e \in \mathrm{out}(v)} 2^{-\tau(e)}}.
\]
This weights $\tau = 1$ edges most heavily (probability
$\propto 1/2$) and $\tau = 4$ edges least heavily (probability
$\propto 1/16$). The expected $\kappa$-step has a more interesting
structure, biased toward the zero-edges that are
$\kappa$-decreasing.

\emph{Verdict.} Plausible, worth investigation. The
$2^{-\tau}$ weighting matches the natural ``measure'' on the
inverse iteration: a $\tau = 1$ step is the most ``frequent'' kind
of inverse step.

\subsection{Candidate C: root-biased chain}

Bias transitions toward goal vertices:
\[
   P(v \to v') \propto e^{\beta (\rho_+(v') - \rho_+(v))}
\]
for a free inverse-temperature parameter $\beta \ge 0$, where
$\rho_+(v) := \nuthree(c(v))$ — proximity to $0$. At
$\beta = 0$ this is uniform; at $\beta \to \infty$ it deterministically
follows the gradient.

\emph{Verdict.} Useful for searching counterexamples (finds the
cheapest path quickly) but not for constructing the potential
directly.

\subsection{Candidate D: Doob transform of a base chain}

Take a base chain $P_0$ (e.g.\ uniform), define
$h(v) := E_v^{P_0}\bigl[ 3^{-\sum_{i} \kappa(e_i)} \bigr]$ (the
generating function), and Doob-transform $P_0$ via $h$:
\[
   P^h(v \to v') = P_0(v \to v') \cdot \frac{h(v')}{h(v)} \cdot 3^{-\kappa(v, v')}.
\]
This transformed chain has expected $\kappa$-step equal to zero
(the standard Doob property), so any $P^h$-trajectory has
$\sum \kappa$ a martingale. By the optional stopping theorem (with
appropriate integrability), the expected $\sum \kappa$ from $s$ to
goal equals the initial $\log h(s) - \log h(\text{goal})$.

If we set $H(v) := -\log_3 h(v)$, the edge inequality
$H(v) \le \kappa(e) + H(v')$ becomes
$h(v) \ge 3^{-\kappa(e)} h(v')$, which is a superharmonicity
relation for $P^h$ on $h$ — automatic by construction.

\emph{Verdict.} This is the cleanest analogue of the von Mangoldt
construction. The challenge: $h$ may not be expressible in closed
form, but the LP we solve in \S\ref{sec:s238} is precisely the
finite-modulus discretisation of this construction. Question B of
\S\ref{sec:alphaevolve} is the systematic search for the right
$P_0$ (and hence the right $h$).

\subsection{Comparison table}

\begin{center}\small
\begin{tabular}{lll}
\toprule
Chain & Strength & Weakness \\
\midrule
A. uniform & simplicity & too uninformative \\
B. $\tau$-weighted & natural & may not capture all structure \\
C. root-biased & efficient search & not a barrier construction \\
D. Doob transform & exact LP equivalence & requires choice of $P_0$ \\
\bottomrule
\end{tabular}
\end{center}

The S238 framework (\S\ref{sec:s238}) effectively explores
candidate D in the finite-modulus regime by sweeping the linear
coefficients $(a, b, c, d)$ and the modular precision $h$.

\section{References and related work}\label{app:references}

\subsection*{Collatz conjecture: surveys and primary literature}

\begin{enumerate}
\item J.~C.~Lagarias. \emph{The $3x+1$ problem and its generalizations}.
       American Mathematical Monthly 92(1):3--23, 1985.
\item J.~C.~Lagarias (editor). \emph{The Ultimate Challenge: The $3x+1$ Problem}.
       American Mathematical Society, 2010.
\item G.~J.~Wirsching. \emph{The Dynamical System Generated by the $3n+1$ Function}.
       Lecture Notes in Mathematics 1681, Springer, 1998.
\item T.~Tao. Almost all orbits of the Collatz map attain almost
       bounded values. \emph{Forum of Mathematics, Pi}, 10(e12), 2022.
       \url{https://doi.org/10.1017/fmp.2022.8}
\item A.~V.~Kontorovich and Ya.~G.~Sinai. Structure theorem for
       $(d, g, h)$-maps. \emph{Bulletin of the Brazilian Mathematical Society}
       33(2):213--224, 2002.
\item J.~H.~Conway. Unpredictable iterations. In \emph{Proc.\ 1972 Number
       Theory Conference}, 49--52. Univ.\ Colorado, 1972.
\item D.~J.~Bernstein and J.~C.~Lagarias. The $3x+1$ conjugacy map.
       \emph{Canadian Journal of Mathematics} 48(6):1154--1169, 1996.
\end{enumerate}

\subsection*{Tao 2026 inspiration}

\begin{enumerate}
\setcounter{enumi}{7}
\item T.~Tao et al. Primitive sets and von Mangoldt chains: Erdős
       Problem \#1196 and beyond.
       Online post,
       \href{https://terrytao.wordpress.com/2026/05/03/primitive-sets-and-von-mangoldt-chains-erdos-problem-1196-and-beyond/}%
            {\texttt{terrytao.wordpress.com/2026/05/03/...}},
       2026; arXiv:2605.00301.
\end{enumerate}

\subsection*{Formal verification methodology}

\begin{enumerate}
\setcounter{enumi}{8}
\item T.~C.~Hales \emph{et al}. A formal proof of the Kepler conjecture.
       \emph{Forum of Mathematics, Pi} 5(e2), 2017.
\item M.~Eberl, L.~C.~Paulson. A formal proof of the prime number theorem
       in Isabelle/HOL. \emph{Archive of Formal Proofs}, 2018.
\item L.~de~Moura and S.~Ullrich. The Lean 4 theorem prover and
       programming language. In \emph{CADE-28}, 2021.
\item The mathlib Community. The Lean Mathematical Library. In
       \emph{CPP 2020}, 367--381, 2020.
\item T.~Tao \emph{et al}. Equational theories project: a Lean
       blueprint methodology for collaborative formal mathematics.
       \url{https://teorth.github.io/equational_theories}, 2024--2026.
\end{enumerate}

\subsection*{Evolutionary coding agents}

\begin{enumerate}
\setcounter{enumi}{13}
\item A.~Novikov \emph{et al}. AlphaEvolve: A coding agent for
       scientific and algorithmic discovery. Google DeepMind, 2025.
       \url{https://deepmind.google/discover/blog/alphaevolve/}
\end{enumerate}

\subsection*{Algorithms}

\begin{enumerate}
\setcounter{enumi}{14}
\item T.~H.~Cormen, C.~E.~Leiserson, R.~L.~Rivest, S.~Stein.
       \emph{Introduction to Algorithms}, 3rd edition. MIT Press, 2009.
       (Bellman--Ford, difference constraints, negative-cycle detection.)
\end{enumerate}

\section{Detailed algorithm: virtual-source Bellman--Ford for difference constraints}\label{app:bf-detail}

The corrected validator of \S\ref{sec:bug} uses a standard
construction from algorithmic graph theory: solving a system of
difference constraints by Bellman--Ford with a virtual source.
We exhibit the algorithm explicitly for self-containment.

\subsection*{Problem statement}

Given a collection of difference constraints of the form
\[
   \psi(a) - \psi(b) \;\le\; w_{ab} \qquad (a, b \in \mathrm{Classes},\; w_{ab} \in \ZZ),
\]
decide whether there exists $\psi : \mathrm{Classes} \to \ZZ$ satisfying
every constraint. If yes, produce such a $\psi$; if no, identify
the obstructing negative cycle.

\subsection*{The reduction}

Build a directed graph $G$:
\begin{itemize}
   \item Vertices: $\mathrm{Classes} \cup \{S^{*}\}$ (virtual
         source $S^{*}$).
   \item Edges: for each constraint $\psi(a) - \psi(b) \le w_{ab}$,
         add an edge $b \to a$ of weight $w_{ab}$.
   \item Edges from virtual source: for each $c \in \mathrm{Classes}$,
         add an edge $S^{*} \to c$ of weight $0$.
\end{itemize}

\subsection*{The algorithm}

\begin{algorithm}[H]
\caption{Virtual-source Bellman--Ford for difference constraints}
\label{alg:bf-diff}
\begin{algorithmic}[1]
\State $\psi[S^{*}] \gets 0$
\For{$c \in \mathrm{Classes}$}
  \State $\psi[c] \gets 0$
\EndFor
\For{$\mathrm{iter}$ from $1$ to $|\mathrm{Classes}| + 1$}
  \State $\mathrm{updated} \gets \mathbf{false}$
  \For{each edge $b \to a$ with weight $w$ in $G$}
    \If{$\psi[b] + w < \psi[a]$}
      \State $\psi[a] \gets \psi[b] + w$
      \State $\mathrm{updated} \gets \mathbf{true}$
    \EndIf
  \EndFor
  \If{not $\mathrm{updated}$}
    \State \textbf{break}    \Comment{converged}
  \EndIf
  \If{$\mathrm{iter} > |\mathrm{Classes}|$}
    \State \Return $(\textsc{infeasible}, \text{negative cycle detected})$
  \EndIf
\EndFor
\State \Return $(\textsc{feasible}, \psi)$
\end{algorithmic}
\end{algorithm}

\paragraph{Correctness.}
Standard theorem of algorithm design (Cormen \emph{et al.}, Chapter
24). If the LP is feasible, BF converges in at most
$|\mathrm{Classes}|$ iterations to a feasible solution. If
infeasible, the
$|\mathrm{Classes}| + 1$-st iteration still relaxes some edge,
witnessing a negative cycle.

\paragraph{Boundary anchors.}
For our specific problem, we add two boundary edges via the same
virtual source:
\begin{itemize}
   \item $\psi[\text{class}_{\text{goal}}] - \psi[S^{*}] \le - L(g_Q)$,
         encoding $H(g_Q) \le 0$;
   \item $\psi[S^{*}] - \psi[\text{class}_{\text{start}}] \le L(\text{start}) - m$,
         encoding $H(\text{start}) \ge m$.
\end{itemize}
With these, the LP is feasible iff the full barrier conditions are
satisfiable.

\subsection*{Complexity}

The graph has $|\mathrm{Classes}| + 1$ vertices and at most
$|E| + |\mathrm{Classes}|$ edges (one per constraint plus virtual
source). The algorithm runs in
$O(|\mathrm{Classes}| \cdot |E|)$ time, which is the standard
Bellman--Ford complexity. For our setting:
\begin{itemize}
   \item $|\mathrm{Classes}|$: at most $3^h$, but only those
         actually appearing in some constraint
         (in practice $|\mathrm{Classes}| = O(|E|)$).
   \item $|E|$: at most $|E_{\mathrm{observed}}|$
         (the BF closure size), typically $\sim 10^2$ to $\sim 10^5$.
\end{itemize}
Hence each LP takes $\sim 10^4$ to $\sim 10^{10}$ operations in
the worst case; in practice the early-stop on convergence gives
much better performance ($\sim 10^3$ to $\sim 10^7$ operations).

\section{Detailed sweep methodology}\label{app:sweep}

The \texttt{tools/s238\_sweep\_parallel.py} implementation is the
canonical sweep. We describe its design.

\subsection*{Outer loop: $(m, Q)$ iteration}

For each $(m, Q) \in \{(1, 3), (1, 5), (2, 3), (2, 6), (3, 6)\}$:
\begin{enumerate}
  \item Precompute the BF closure for $(m, Q)$ \emph{once}; this is
        the expensive step. Cache the resulting \texttt{cert} (dist
        table) in memory.
  \item For each $h \in \{22, 23, 25, 28, 30\}$ (innermost loop):
        dispatch the grid of $(a, b, c, d)$ coefficients to a
        \texttt{multiprocessing.Pool}.
  \item Early-stop on the first feasible cell.
  \item Export the resulting $\psi$ in sparse JSON.
\end{enumerate}

\subsection*{Inner loop: coefficient grid}

The grid is
\[
   (a, b, c, d) \;\in\; \{-3, -2, -1, 0, 1\} \times \{-3, -2, -1, 0, 1\}
                       \times \{-1, 0, 1, 2\} \times \{-1, 0, 1, 2\}
   \;=\; 400 \text{ cells per } h.
\]
The grid is dispatched as $400$ tasks to the pool. Each task:
\begin{enumerate}
   \item receives the cached \texttt{cert} (passed by pickle);
   \item reconstructs a lightweight engine instance locally;
   \item enumerates outgoing edges from each $v$ in
         \texttt{cert}, filtering to those with target also in
         \texttt{cert};
   \item builds the difference-constraint system;
   \item runs the virtual-source BF
         (Algorithm~\ref{alg:bf-diff});
   \item returns \textsc{feasible} or \textsc{infeasible}.
\end{enumerate}

\subsection*{Sparse export}

When a cell is feasible, the worker returns the $\psi$ as a
\emph{sparse} \texttt{dict[int, int]} mapping class
$\to$ integer value. The exporter then serialises to JSON,
recording only the classes that participated in some constraint.
For $h = 22$, the modular size is $3^{22} \approx 3 \times 10^{10}$
classes; a dense export would consume tens of GB, while the
sparse export is typically $< 100$ KB.

\subsection*{Verifier matching}

The corresponding verifier
(\texttt{tools/s238\_psi\_verifier\_precise.py}) treats missing
sparse entries as $\psi(\cdot) = 0$. This is mathematically
equivalent: any constraint involving a class not in the sparse
$\psi$ is automatically satisfied by setting $\psi = 0$ there
(the LP allows it).

\subsection*{Parallelisation pattern}

The grid of $400$ cells per $h$ is embarrassingly parallel: each
cell is an independent LP. The wall-clock speedup is approximately
$N$-fold for $N$ cores, modulo serialisation overhead. On an
$8$-core laptop, the total sweep for the five frontier $(m, Q)$
pairs at $h \in \{22, 23, 25, 28, 30\}$ takes approximately:
\begin{center}\small
\begin{tabular}{rrll}
\toprule
$(m, Q)$ & BF closure size & per-LP cost (single-thread) & total parallel \\
\midrule
$(1, 3)$ & $137$    & $< 0.01$ s & $< 1$ s \\
$(1, 5)$ & $4\,425$ & $\sim 1$ s & $\sim 1$ min \\
$(2, 3)$ & $285$    & $< 0.1$ s  & $\sim 10$ s \\
$(2, 6)$ & $61\,645$ & $\sim 5$ s & OOM on $8$ GB/core \\
$(3, 6)$ & $134\,913$ & $\sim 10$ s & OOM on $8$ GB/core \\
\bottomrule
\end{tabular}
\end{center}
The last two cases hit the per-worker memory ceiling; the
\texttt{multiprocessing.Pool} workers are OOM-killed. A
machine with $\sim 32$ GB per worker would complete these.

\section{Worked psi-table excerpts}\label{app:psi-tables}

For concreteness, we exhibit excerpts of the validated
$\psi$-table for $(m, Q, h) = (1, 3, 22)$ with coefficients
$(a, b, c, d) = (-1, -3, -1, 2)$. The full table has $80$ sparse
entries; the first $20$ are:

\begin{center}\footnotesize
\begin{tabular}{rrr|rrr}
\toprule
class $z$ & $\psi(z)$ & note & class $z$ & $\psi(z)$ & note \\
\midrule
$1$  & $42$  & $z = \aSat \bmod 3^{22}$ & $13$  & $-3$  & \\
$2$  & $-1$  &                          & $19$  & $7$   & \\
$4$  & $5$   &                          & $20$  & $-2$  & \\
$5$  & $0$   &                          & $22$  & $-2$  & \\
$7$  & $11$  &                          & $40$  & $4$   & \\
$8$  & $-4$  &                          & $58$  & $-1$  & \\
$11$ & $3$   &                          & $76$  & $0$   & \\
$112$ & $-3$ &                          & $94$  & $-2$  & \\
$130$ & $1$  &                          & $112$ & $-3$  & \\
$166$ & $-2$ &                          & $238$ & $-1$  & \\
\bottomrule
\end{tabular}
\end{center}

\noindent The entry $\psi(1) = 42$ at class $z = 1$ (which equals
$\aSat \bmod 3^{22} = 1$ as well, since $\nuthree(\aSat - 1) = 22$)
is the highest value in the table. Its magnitude is determined by
the boundary anchors: at start, $L(s) = -41$, so to have
$H(s) = L(s) + \psi(\text{class}_{\text{start}}) \ge m = 1$ we need
$\psi(1) \ge 42$. The LP-computed value matches exactly.

\subsection*{Verification trace}

Running the verifier:
\begin{lstlisting}
$ python tools/s238_psi_verifier_precise.py \
    tools/s238_tables_parallel/psi_m1_Q3_h22_precise_parallel.json
{
  "overall_ok": true,
  "m": 1, "Q": 3, "h": 22, "mode": "precise",
  "linear_coefficients": {
    "a_dp": -1, "a_rm": -3, "a_j": -1, "a_jQ": 2
  },
  "n_observed_edges_checked": 136,
  "edge_check_pass": true,
  "worst_edge_violation": null,
  "H_start": 1,
  "H_goal_per_j": [
    [0, -4], [1, -3], [2, -2], [3, -1]
  ],
  "max_H_goal": -1,
  "endpoint_margin": 2,
  "endpoint_required": 1,
  "endpoint_pass": true
}
\end{lstlisting}

\noindent The verifier independently checks:
\begin{enumerate}
   \item Every observed edge $v \to v'$ satisfies
         $H(v) \le \kappa(e) + H(v')$ (passed: $136/136$).
   \item $H(\text{start}) \ge m = 1$ (passed: $H(\text{start}) = 1$).
   \item $H(g_j) \le 0$ for $j \in \{0, 1, 2, 3\}$ (passed:
         $\max_j H(g_j) = -1 \le 0$).
\end{enumerate}
The certificate is valid over the observed graph.

\section{Glossary of Lean module names}\label{app:lean-modules}

For the convenience of a reader wishing to navigate the Lean
development directly, the modules and their roles:

\begin{center}\footnotesize
\begin{tabular}{lp{8.5cm}}
\toprule
Module & Role \\
\midrule
\texttt{AdmissibleBasic} & Admissible state machine; translation identity; resonance lemma. \\
\texttt{TranslationSets} & $B_{\mathrm{adm}}, B_{\mathrm{free}}$; coverage equivalence. \\
\texttt{Defect} & Defect formulation; S202 word evaluated; high-slope refutation. \\
\texttt{NegativeRuns} & $L^k$ negative-run framework (S203-C/D). \\
\texttt{PeriodicBlocks} & S206 identity modulo $3^{24}$. \\
\texttt{ThreeAdic} & S206 lifted to $\ZZ/3^k$ for all $k \ge 1$. \\
\texttt{Reachability} & Explicit \texttt{InImageTheta} table. \\
\texttt{BinaryTensor} & 2-adic Terras coordinates, parity vectors. \\
\texttt{S202Cylinders} & Cylinder family $C_m$, accessibility criterion. \\
\texttt{Concatenation} & S210/S211 concatenation and defect additivity. \\
\texttt{CylinderStability} & S213 $\tau$-trajectory stability. \\
\texttt{PotentialBarrier} & S214 abstract potential framework. \\
\texttt{InverseGraph} & Inverse-edge predicates, \texttt{outgoingEdges} enumerator. \\
\texttt{AS202Lift} & $\aSat$ for $m \in \{1, 2, 3\}$, with Cauchy coherence. \\
\texttt{S216} & Closure-to-barrier theorem with optimistic credit. \\
\texttt{S212Forward} & Forward path-word correspondence. \\
\texttt{AnalyticBarrier} & Per-edge defect bound; precise / coarse conjectures; conditional slope barrier. \\
\texttt{Cert\_m1\_Q3\_S216(\_Barrier)} & Auto-generated barrier $(m=1, Q=3)$. \\
\texttt{Cert\_m1\_Q5\_S216(\_Barrier)} & Auto-generated barrier $(m=1, Q=5)$. \\
\texttt{Cert\_m1\_Q8\_S216(\_Barrier)} & Auto-generated barrier $(m=1, Q=8)$. \\
\texttt{Cert\_m1\_Q10\_S216(\_Barrier)} & Auto-generated barrier $(m=1, Q=10)$. \\
\bottomrule
\end{tabular}
\end{center}

\section{Recommended reading paths}\label{app:reading-paths}

This document is dense; we suggest reading paths for different
target audiences.

\subsection*{For a mathematician interested in Collatz}

\begin{enumerate}
   \item \S\ref{sec:framework} (framework) and
         \S\ref{sec:s202-word} (S202 word, 3-adic lift).
   \item \S\ref{sec:cylinder-family} (cylinder family) and
         \S\ref{sec:inverse-graph} (inverse graph).
   \item \S\ref{sec:open-conjecture} (the conjecture).
   \item \S\ref{sec:tao-inspiration} and \S\ref{sec:s230-kappa}
         (the Tao 2026 inspiration and the $\kappa$ reformulation).
\end{enumerate}
This omits the formal Lean detail and the bug analysis; sufficient
for a high-level overview.

\subsection*{For a formal-verification researcher}

\begin{enumerate}
   \item \S\ref{sec:formal-status} (formal verification status,
         module map).
   \item \S\ref{sec:bf-closure} (S216 closure-to-barrier theorem
         with full proof).
   \item Appendix~\ref{app:lean-code} (Lean code excerpts).
   \item Appendix~\ref{app:reprod} (reproducibility).
\end{enumerate}
The reader can then clone the repository and verify the headline
theorems directly.

\subsection*{For an AlphaEvolve developer}

\begin{enumerate}
   \item Abstract and \S\ref{sec:open-conjecture} (the open
         question).
   \item \S\ref{sec:alphaevolve} (the two questions and the
         cooperation protocol).
   \item Appendix~\ref{app:compute} (compute budget estimates).
   \item Appendix~\ref{app:bf-detail} (the algorithm in detail).
\end{enumerate}
Sufficient to evaluate whether the questions are within the
agent's operational profile.

\subsection*{For someone diagnosing the bug}

\begin{enumerate}
   \item \S\ref{sec:s233} (the candidate ansatz).
   \item \S\ref{sec:bug} (the bug diagnosis).
   \item Appendix~\ref{app:bug} (concrete walkthrough).
   \item Appendix~\ref{app:ols} (OLS analysis that prompted the
         refined ansatz).
\end{enumerate}

\section{Visual summary of the programme structure}\label{app:visual}

\subsection*{Dependency diagram}

The logical dependencies among the main objects:

\begin{center}\footnotesize
\begin{tabular}{l}
\toprule
\textbf{Foundation} \\
$\tau$-table on $\InX$ \\
$\downarrow$ \\
admissible state machine $(n, A, q, B)$ \\
$\downarrow$ \\
translation identity $3^q n + B = 2^A$ \\
$\downarrow$ \\
defect $\defect = A - 2q$, per-edge bound \\
\midrule
\textbf{S202 fixed point} \\
$\wS$ word, $\defect(\wS) = -1$ \\
$\downarrow$ \\
3-adic Cauchy: $\aSat$ for each $m$ \\
$\downarrow$ \\
$\nuthree(\aSat - 1) = 22$ (structural lemma) \\
\midrule
\textbf{Cylinder accessibility} \\
$C_m$ family, cost $\Cback(m, Q)$ \\
$\downarrow$ \\
inverse cylinder graph $\InvG_{m, Q}$ \\
$\downarrow$ \\
path-word correspondence forward \\
\midrule
\textbf{Defect barriers (formal in Lean)} \\
suffix weight bound, prefix relevance \\
$\downarrow$ \\
optimistic credit closure (S216) \\
$\downarrow$ \\
4 explicit terminal barriers (Q in \{3, 5, 8, 10\}) \\
\midrule
\textbf{The open conjecture} \\
$\numOnes - \numNeg \ge m$ (precise) \\
$\Updownarrow$ (formal implication) \\
slope barrier $\lambda < 2$ on admissible cover \\
\midrule
\textbf{Markov / flow dual programme} \\
$\kappa$-weight reformulation \\
$\downarrow$ \\
S233 ansatz $H = L + \psi$ \\
$\downarrow$ \\
3 observed-graph certificates (after bugfix) \\
$\downarrow$ \\
S234 abstract over-cover (open) \\
\bottomrule
\end{tabular}
\end{center}

\subsection*{What is proved vs. open}

\begin{center}\small
\begin{tabular}{p{8cm}l}
\toprule
Statement & Status \\
\midrule
Translation identity & formal (Lean) \\
S202 refutation of $A \ge 2q$ & formal (Lean) \\
3-adic Cauchy lift of $\aSat$ for $m \in \{1, 2, 3\}$ & formal (Lean) \\
Forward path-word correspondence & formal (Lean) \\
S211 conditional concatenation & formal (Lean) \\
S216 closure-to-barrier (defect weight) & formal (Lean) \\
$\Cback(1, Q) \ge 1$ for $Q \in \{3, 5, 8, 10\}$ & formal (Lean, native\_decide) \\
$\Cback(m, Q) \ge m$ for $(m, Q) \in \{(2, 6), (2, 9), (3, 6)\}$ & empirical (gated on InvStart refactor) \\
Coarse $\Rightarrow$ precise conjecture & formal (Lean) \\
Conditional slope barrier (from precise) & formal (Lean) \\
Precise conjecture itself & \emph{open} \\
$\psi$-certificate on observed graph (3 cylinders) & empirical (verified independently) \\
Abstract over-cover certificate (any cylinder) & \emph{open} (Question B) \\
\bottomrule
\end{tabular}
\end{center}

\section{Historical context and connection to classical results}\label{app:history}

We briefly situate the present programme within the long tradition
of Collatz analysis, primarily to clarify which structural insights
are inherited from prior work and which are original.

\subsection{The Collatz problem: a brief history}

The $3n + 1$ problem was posed by Lothar Collatz around 1937 as a
graduate-student curiosity. Its first systematic study appeared in
the 1972 work of Conway, who recast the iteration as the
\emph{halting problem for FRACTRAN}, a Turing-complete computation
model. Conway's reformulation made it clear that the conjecture
is genuinely deep: it asks whether a specific Turing machine
always halts on positive integer inputs.

Lagarias's 1985 survey
\cite[ref.~1]{app:references}
crystallised the state of the art and identified the key
quantitative parameters (the slope $\lambda$, the cycle stopping
times $\sigma$, etc.). The 2010 collection edited by Lagarias
\cite[ref.~2]{app:references} updated the picture.

\subsection{The 3-adic / cylinder framework}

The 3-adic side of the analysis goes back to Wirsching's 1998
book \cite[ref.~3]{app:references} and the conjugacy
maps of Bernstein--Lagarias 1996
\cite[ref.~7]{app:references}. Both established that the Collatz
iteration, viewed appropriately, lives naturally in
$\Zp{3}$ and that the dynamics there is conjugate to a 3-adic shift.

The present programme specialises this general 3-adic
framework to a \emph{specific} fixed point — the S202 point
$\aS = 1 + 3^{22}$ — and tracks the associated quantitative
accessibility cost. The specific S202 word $\wS$ is, to the
author's knowledge, new in this programme; the construction of
the cylinder family $C_m$ at precision $3^{22m + 2}$ is also new.

\subsection{The stochastic side and Tao 2022}

The stochastic side of Collatz analysis is largely orthogonal:
Kontorovich and Sinai 2002
\cite[ref.~5]{app:references} initiated the modern probabilistic
treatment, viewing the parity vector as a random sequence under a
natural measure. Tao's 2022 \emph{Forum Math.\ Pi} bounded-value
theorem is the strongest result in this tradition, asserting that
almost every Collatz orbit attains an almost bounded value.

The present programme is the structural complement: it asks not
``what does a typical $n$ do?'' but ``what is the cost of reaching
a specific cylinder $C_m$ from $1$?''. The two questions are
logically independent (\S\ref{sec:relation-tao}).

\subsection{Formal verification of Collatz: prior attempts}

To the author's knowledge, there is no prior end-to-end Lean
formalisation of any nontrivial Collatz-related result. There
have been:
\begin{itemize}
   \item Coq formalisations of the trivial $\{1, 2, 4\}$ cycle
         and small computational verifications.
   \item Lean snippets on Mathlib's \texttt{Combinatorics/Collatz}
         module, with the basic definitions but no nontrivial
         theorems.
   \item Isabelle/HOL implementations of toy versions for
         pedagogical purposes.
\end{itemize}

The present development is therefore, as far as the author knows,
the first sustained formal verification of a substantial
Collatz-related programme. The 8\,425+ verification jobs and
zero \texttt{sorry} statements constitute a baseline that any
future formal Collatz work can build on.

\subsection{Conway's FRACTRAN and the present framework}

It is worth noting that Conway's FRACTRAN reformulation is a
\emph{forward} encoding: the iteration is presented as a chain of
rational multiplications, with the multiplier sequence
$(3/2, 1/2)$ or appropriate variants. The admissible inverse
machinery of the present programme can be viewed as a
\emph{backward} FRACTRAN: each admissible word corresponds to a
finite sequence of inverse rational operations, terminating in a
specific rational target ($\aSat / 3^{R(m)}$).

The translation between the two viewpoints is mediated by the
state machine of \S\ref{sec:framework}: forward FRACTRAN
operates on $n$, while the admissible inverse machinery operates
on the four-tuple $(n, A, q, B)$ with the translation identity
$3^q n + B = 2^A$. Both encode the same dynamical content; the
present choice highlights the cylinder accessibility cost as the
quantitative parameter to bound.

\subsection{The Bernstein--Lagarias conjugacy map}

The 1996 Bernstein--Lagarias conjugacy map sends Collatz orbits
to a 3-adic shift map. Specifically, there exists a continuous,
bijective $\Phi : \Zp{3} \to \Zp{2}$ that conjugates the
3-adic-extended Collatz map to the 2-adic shift. The present
programme works within $\Phi$'s domain and tracks specific
parameters of the 3-adic side.

\subsection{Originality summary}

For the avoidance of doubt, the following are original to this
programme:
\begin{itemize}
   \item The S202 word $\wS$ and its parameters
         $(251, 43, 22, B_{S202})$, with the explicit
         verification.
   \item The cylinder family $C_m$ at precision $3^{22m + 2}$.
   \item The cost function $\Cback(m, Q)$ and its formal
         definition.
   \item The inverse cylinder graph $\InvG_{m, Q}$ with explicit
         edge predicates.
   \item The S216 optimistic-credit closure framework with
         Bool-decidable certificates.
   \item The path-word correspondence (forward direction
         formalised).
   \item The four explicit terminal barriers
         $\Cback(1, Q) \ge 1$ for $Q \in \{3, 5, 8, 10\}$.
   \item The reduction of the unconditional content to
         \texttt{S202\_one\_edge\_count\_conjecture}.
   \item The $\kappa$-weight reformulation as shortest-path
         dual.
   \item The validator bug story and its structural explanation
         via $\nuthree(\aSat - 1) = 22$.
   \item The three independently-verified $\psi$-certificates
         over the observed BF graph.
\end{itemize}

The following are inherited from prior work and used as
building blocks:
\begin{itemize}
   \item The general 3-adic conjugacy framework
         (Bernstein--Lagarias).
   \item The accelerated Collatz map $T$ (Terras).
   \item The Bellman--Ford algorithm and difference constraints
         (classical algorithms).
   \item The Lean 4 + Mathlib ecosystem.
\end{itemize}

\section{Methodological discussion}\label{app:method}

\subsection{Why formal verification matters here}

A natural question is: \emph{why} formalise this programme in
Lean 4, when many published Collatz results are not formalised?
Two answers.

\paragraph{Eliminating undetected validator bugs.}
The bug story of \S\ref{sec:bug} is illustrative. An informal
LP-based certificate could have been published, peer-reviewed,
and accepted before the boundary-anchor issue surfaced. Formal
verification forces every step into the open. The
\texttt{S216\_barrier\_for\_words\_outgoing} theorem in Lean
explicitly anchors all relevant boundary conditions; the
analogous bug could not occur in the formal pipeline.

\paragraph{Building a reusable foundation.}
The Lean modules of this programme can be used by future Collatz
researchers as building blocks. The translation identity, the
forward path-word correspondence, the optimistic-credit closure
theorem — these are general tools, not specific to the S202
cylinder. Future work can specialise to other cylinders or
extend the closure framework without re-deriving these
foundations.

\subsection{Why empirical certificates matter}

A symmetric question: \emph{why} include empirical certificates
in a formal-verification programme? Three answers.

\paragraph{Guiding the analytical search.}
The empirical observation that the $\kappa_{\mathrm{precise}}$
variant is uniformly more efficient than the $\kappa_{\mathrm{coarse}}$
variant is the kind of structural data that points the way to the
correct analytical formulation. Without the empirical run, we
might have wasted effort proving the over-strong coarse
conjecture.

\paragraph{Exhibiting concrete obstacles.}
The empirical observation that the candidate ansatz
$H = 2\delta_+ - \rho_- + (j - Q) + \psi$ fails at $h = 3$ (after
the bugfix) but succeeds at $h \ge 22$ tells us that the
3-adic precision matters. Combined with the structural lemma
$\nuthree(\aSat - 1) = 22$, this gives a quantitative diagnostic
of the obstruction.

\paragraph{Producing testable predictions.}
The empirical scaling factor $3.7$ per unit of $Q$ predicts
specific compute requirements for future searches. If an
evolutionary agent reaches $(1, 18)$ in $\sim 1$ hour, that
agrees with the scaling; if it takes $\sim 1$ day, something
unexpected is happening.

\subsection{Why Tao 2026 specifically}

Among the many possible methodological inspirations in recent
mathematics, the Tao 2026 primitive-sets paper was chosen because
the structural pattern (combinatorial obstruction $\to$ chain
duality $\to$ potential certificate) is unusually clean. The
inverse cylinder graph is a directed graph with negative edges,
a setting where Bellman--Ford and difference constraints are
standard tools. The von Mangoldt-weighted chain is the
specific instance where the chain measure aligns with the target
weight — the same alignment problem we face in choosing
$\kappa_{\mathrm{precise}}$ over $\kappa_{\mathrm{coarse}}$ and
in selecting the parametric form of $H$.

The alternative inspirations (e.g., Tao's polynomial method
papers, the equational theories project, the Liouville function
work) do not have the same direct structural overlap. Tao 2026 is
the natural fit.

\subsection{Why not other LLM-guided agents}

The choice of AlphaEvolve as the natural cooperation target is
based on three considerations:

\paragraph{Documented capability profile.}
AlphaEvolve's public documentation describes its operational
profile: it accepts a problem specification with explicit
success/failure criteria, runs evolutionary search over candidate
programs, and reports best-found solutions. Both Question A and
Question B fit this profile exactly.

\paragraph{Affiliation with mathematical research.}
AlphaEvolve is built by Google DeepMind, which has prior
collaborations with Prof.\ Tao on mathematical AI (AlphaProof,
AlphaGeometry). The institutional context for cooperation already
exists.

\paragraph{Scale of compute.}
The compute scale of AlphaEvolve (publicly documented as
``thousands of TPU-hours per problem'') is what is needed for
Question A at large $(m, Q)$. Other open-source LLM-guided
agents do not have the same scale.

That said, the questions of \S\ref{sec:alphaevolve} are
\emph{not} AlphaEvolve-specific. Any evolutionary or LLM-guided
agent with comparable scale and access to LP solvers could
attempt them.

\subsection{Counterfactuals: what if the conjecture is false?}

It is worth considering what would happen if
Conjecture~\ref{conj:precise} turned out to be false. The
$\kappa$-precise closure search has not yet found a counterexample
in the explored region, but the regions explored are small. A
counterexample would be an admissible word $u \in C_m$ with
$q(u) \le Q$ (for some $(m, Q)$) and
$\sum_{e \in p_u} \kappa_{\mathrm{precise}}(e) < m$.

If found, this would:
\begin{enumerate}
   \item Refute the slope barrier
         $\lambda_{\mathrm{asy}} < 2$ as currently formulated.
   \item Force a re-formulation of the cylinder accessibility
         problem — perhaps with a different admissibility section,
         or a different weight function, or a finer cylinder
         structure.
   \item Be an interesting result in its own right, since the
         empirical evidence so far has not produced one.
\end{enumerate}

The author considers a refutation unlikely but possible. Question
A of \S\ref{sec:alphaevolve} is structured so that either
outcome (barrier proof or refutation) is mechanically reportable
and informative.

\subsection{What this work means for AI-assisted mathematics}

The interaction between formal verification (Lean) and empirical
search (Python) in this programme illustrates a methodology that
the author believes is broadly applicable:

\begin{enumerate}
   \item Identify a sharp open conjecture as the analytical gap.
   \item Translate it to a sequence of finite-modulus
         questions amenable to LP / SAT / SMT methods.
   \item Run empirical searches and produce certificates in a
         standard JSON format.
   \item Auto-export certificates to Lean modules with
         \texttt{native\_decide}-based verification.
   \item Build the Bool-Prop adapter that lifts
         \texttt{native\_decide}-verified Bool propositions to
         Prop-level theorems.
   \item Verify end-to-end and tag the resulting commit as a
         reproducible artefact.
\end{enumerate}

This pipeline scales (in principle) to any combinatorial barrier
problem where the obstruction can be cast as a finite-modulus
search. The Collatz cylinder problem is one instance; others may
be found in additive combinatorics (Behrend-type bounds),
analytic number theory (Selberg sieve majorants), or extremal
graph theory.

\section{Frequently asked questions}\label{app:faq}

\subsection*{Q: Is this a proof of the Collatz conjecture?}

No. The honest scope statement in the abstract and at the start
of every Part repeats this: the present work \emph{does not} prove
the Collatz conjecture. It reformulates the asymptotic question
and reduces it to a single explicit combinatorial conjecture.

\subsection*{Q: Why should I trust the Lean development?}

Because:
\begin{enumerate}
   \item The source is public on GitHub.
   \item It builds end-to-end under \texttt{lake build} with
         \textbf{$0$} \texttt{sorry} statements.
   \item The CI pipeline runs on every commit and reports the
         build status publicly.
   \item Every claimed theorem can be re-verified by running
         \texttt{\#print axioms} on it, which exposes any hidden
         axiom dependencies.
   \item The Mathlib library on which it depends is itself a
         peer-reviewed community project with thousands of
         contributors.
\end{enumerate}

\subsection*{Q: Why should I trust the Python tooling?}

Because:
\begin{enumerate}
   \item The Python tools produce JSON certificates that are
         independently verifiable by a separate Python verifier
         (\texttt{s238\_psi\_verifier\_precise.py}) which does
         not share state with the generator.
   \item The Lean certificates (\texttt{Cert\_m1\_Q*\_S216.lean})
         do not depend on the Python tools at runtime: they
         re-verify the JSON inside Lean using
         \texttt{native\_decide}, with the Lean kernel as the
         final arbiter.
   \item The Python source is small ($\sim 2\,000$ lines, pure
         standard library) and readable.
\end{enumerate}

\subsection*{Q: What if the bug in \S\ref{sec:bug} is not the only one?}

A reasonable concern. The author has audited the corrected
validator and the verifier independently and verified that all
boundary conditions are now explicit. The CI pipeline checks the
end-to-end Lean compilation; any new bug would surface there.
That said, formal verification is not a guarantee of correctness
\emph{outside} the formal system: it is a guarantee of internal
consistency. External assumptions (the choice of $\InX$, the
$\tau$-table, the cylinder family) are stated explicitly and
checked by \texttt{native\_decide}; if an undetected error exists,
it would be in one of these places, and the present document
flags them explicitly.

\subsection*{Q: How does this compare to other Lean-formalised number theory?}

The Lean theorem $S216\_barrier\_for\_words\_outgoing$ is, in scale
and depth, comparable to the formal proofs of:
\begin{itemize}
   \item Quadratic reciprocity (~1\,000 lines in Mathlib).
   \item Fermat's two-square theorem (~500 lines in Mathlib).
   \item Specific cases of the abc conjecture (Mochizuki's work
         is not formalised; smaller cases are).
\end{itemize}
At $\sim 12\,000$ new lines + 4 explicit barrier certificates +
the conditional slope-barrier reduction, the present development
is on the larger end of Lean number-theory projects but well
within Mathlib's typical contribution size.

\subsection*{Q: What is the timeline for the next steps?}

\paragraph{Short term (next 1--3 months).}
\begin{itemize}
   \item Complete the \texttt{InvStart} refactor, lifting the
         three pending empirical barriers to formal Lean.
   \item Run the sweep at larger memory to complete the
         $(2, 6), (3, 6)$ $\psi$-certificate exports.
\end{itemize}

\paragraph{Medium term (next 6--12 months).}
\begin{itemize}
   \item Question A on $(1, 12), (1, 15)$ via laptop/cluster
         compute (no AlphaEvolve required).
   \item First attempts at Question B (abstract over-cover)
         with limited compute.
\end{itemize}

\paragraph{Long term (12+ months).}
\begin{itemize}
   \item Cooperation with AlphaEvolve or comparable agents on
         the harder Question A and B cases.
   \item Possible refinement of the conjecture based on findings.
   \item Submission of a journal version after independent
         verification.
\end{itemize}

\subsection*{Q: How can I help?}

Several ways:
\begin{itemize}
   \item Clone the repository and run \texttt{lake build}.
         Confirm the build passes and the axiom check is clean.
         Open issues if anything fails.
   \item Run the Python pipeline and reproduce the three
         $\psi$-certificates. Open issues if your reproduction
         differs.
   \item Read the Lean source and audit for subtle issues.
         Open PRs if you find improvements.
   \item Attempt Conjecture~\ref{conj:precise} by analytical
         methods. The author would welcome a paper sketch or a
         direct proof.
   \item Attempt the abstract over-cover (S234) construction
         independently. Even small steps (over-cover at small $h$
         for small $(m, Q)$) would be valuable.
\end{itemize}

\section{Closing remarks and acknowledgements}\label{app:closing}

\subsection*{Personal note from the author}

The S189--S238 programme has been the author's
research focus over the past several months, conducted as an
independent researcher in Barcelona. The Lean formalisation was
built atop the public Mathlib library
\cite{mathlib} using the toolchain \texttt{leanprover/lean4:v4.30.0-rc2}
\cite{lean4}. The methodological inspiration of Tao's 2026
primitive-sets paper was decisive in moving from defect-weight
barriers to the $\kappa$-weight reformulation; that move opened
the analytical question stated in
\S\ref{sec:open-conjecture} in its sharpest form.

The bug story of \S\ref{sec:bug} is offered without embarrassment:
in collaborative formal-verification work, undetected validator
bugs are an occupational hazard, and the structural reason this
particular one was hard to surface
($\nuthree(\aSat - 1) = 22$, an intrinsic 3-adic property)
makes it independently interesting. Future $\psi$-miner pipelines
should always include explicit boundary anchors as part of the
LP, not as a post-verification check.

\subsection*{The scope of the request}

The author submits this document for two specific purposes:

\paragraph{Independent verification.}
The Lean development is public, reproducible end-to-end, and
should be assessable by any mathematician with access to
\texttt{lake build}. The author would welcome corrections,
suggestions, or identifications of any subtle issue in the
formalisation.

\paragraph{Cooperation with an evolutionary coding agent.}
The two questions of \S\ref{sec:alphaevolve} are sharp and
finitely specified. They are within the operational profile of
AlphaEvolve as documented publicly. If the recipient's
affiliations include access to such compute, the author would
welcome cooperation on either question. If not, the questions
can be pursued by other means (cluster compute, distributed
volunteers, or simply patience as laptop CPUs improve).

\subsection*{The honest scope of what this work achieves}

This work does not prove the Collatz conjecture. It does the
following:

\begin{enumerate}
   \item Reformulates the Collatz dynamics in 3-adic
         admissible-inverse language, with the entire
         reformulation formally verified in Lean 4.
   \item Identifies an explicit counterexample $\wS$ to the
         rigid bound $A \ge 2q$, formalised and traced.
   \item Defines the cylinder accessibility cost
         $\Cback(m, Q)$ as the natural quantitative refinement.
   \item Develops the optimistic-credit closure framework (S216),
         formally verified, which transforms empirical search
         outputs into formal Lean theorems.
   \item Produces four explicit terminal barrier theorems
         $\Cback(1, Q) \ge 1$ for $Q \in \{3, 5, 8, 10\}$.
   \item Identifies the conditional slope barrier
         $\Cback(m, Q) \ge m$ for all $(m, Q)$ as a one-line
         consequence of a single combinatorial conjecture,
         \texttt{S202\_one\_edge\_count\_conjecture}.
   \item Translates the open conjecture to a shortest-path
         dual via the $\kappa$ reformulation
         (Tao 2026 methodology).
   \item Produces three independently-verified $\psi$-certificates
         over the observed BF graph, with a structural lemma
         ($\nuthree(\aSat - 1) = 22$) explaining why the modular
         precision must exceed $h = 22$.
\end{enumerate}

The work does \emph{not} do the following:

\begin{enumerate}
   \item Prove the Collatz conjecture.
   \item Prove $\lambda_{\mathrm{asy}} < 2$ unconditionally on
         admissible covers.
   \item Resolve \texttt{S202\_one\_edge\_count\_conjecture}.
   \item Construct an analytic potential $\Phi$ closing the
         conjecture for arbitrary $(m, Q)$ over the abstract
         over-cover (the natural next step, formulated as
         Question B of \S\ref{sec:alphaevolve}).
\end{enumerate}

\subsection*{Acknowledgements}

The author thanks the Lean and Mathlib communities for the
infrastructure on which this work rests, and the anonymous
external collaborator whose S233 mining package, including its
latent validator bug, prompted the careful re-derivation
documented in \S\ref{sec:bug} and the subsequent refined ansatz.

The author also thanks the wider Collatz research community for
the long tradition of careful, honest, scope-conscious work that
makes a programme such as the present one feasible.

All remaining errors and misjudgements are the author's own.

\bigskip
\noindent\hrulefill\par
\smallskip\noindent
\begin{tabular}{ll}
\textbf{Author:}     & Benet Andújar Andújar \\
\textbf{Email:}      & \href{mailto:bandujar@xtec.cat}{\texttt{bandujar@xtec.cat}} \\
\textbf{Repository:} & \url{https://github.com/benetandujar72/collatz-admissible-tree} \\
\textbf{Date:}       & May 2026 \\
\textbf{Status:}     & Open for cooperation and review \\
\end{tabular}

\bigskip
\noindent\hrulefill\par
\smallskip\noindent
\textit{End of document.}

\bigskip
\noindent\hrulefill\par
\smallskip\noindent
\textbf{Document statistics.}
This document consists of approximately $5\,000$ LaTeX source
lines, organised into 6 Parts and 11 appendices, with:
\begin{itemize}
   \item ${\sim} 30$ formal theorem and lemma statements
         with Lean references;
   \item ${\sim} 20$ tables summarising data, modules, and
         status;
   \item ${\sim} 15$ Lean code excerpts;
   \item ${\sim} 5$ explicit numerical worked examples;
   \item A full Bellman--Ford and S238 sweep walkthrough;
   \item A two-question cooperation protocol for AlphaEvolve;
   \item A reading-path guide for four target audiences.
\end{itemize}

Total page count (after compilation with \texttt{pdfLaTeX} in
Overleaf, default 11pt, A4): approximately 60--70 pages.

\bigskip
\noindent\hrulefill\par
\smallskip\noindent
\textit{Submitted with respect for the recipient's time and with
gratitude for the depth of the mathematical traditions on which
this work rests.}

\bigskip
\noindent\hrulefill\par
\smallskip\noindent
\textbf{Tone register.} Following the Japanese mathematical
letter-writing tradition (It\^{o}, Sato, Hironaka), this document
uses austere declarative sentences, passive constructions for
objectivity, and explicit honest scoping. No claim is made that
the sender considers self-evident. Errors and confusions are
documented openly.

\end{document}
